%
\documentclass[11pt]{article}
\usepackage{amsmath,amsfonts,amsthm,amssymb,verbatim}
\usepackage{epsfig}
%\usepackage{showkeys}

\newcommand\version{11-19, 2026}
\numberwithin{equation}{section}

\def\lanbox{\hbox{$\, \vrule height 0.25cm width 0.25cm depth 0.01cm \,$}}
\newcommand{\e}{\varepsilon}
\newcommand{\rd}{{\rm d}}
\newcommand{\x}{{\rm x}}
\newcommand{\wh}{\widehat}
\newcommand{\ann}{a^{\hphantom{+}}}
\newcommand{\cre}{a^{\dagger}}
\renewcommand{\qedsymbol}{\lanbox}
\newcommand{\beq}{\begin{equation}}
\newcommand{\eeq}{\end{equation}}


\newcommand{\beqa}{\begin{eqnarray}}
\newcommand{\eeqa}{\end{eqnarray}}
\newcommand\Z{{\mathbb Z}}
\newcommand\N{{\mathbb N}}
\newcommand\C{{\mathbb C}}
\newcommand\SsFF{{\widetilde S}}
\newcommand\R{{\mathbb R}}
\newcommand\eps{\varepsilon}
\newcommand\half{\mbox{$\frac 12$}}
\newcommand{\mfr}[2]{{\textstyle\frac{#1}{#2}}}
\newcommand\rhov\varrho
\newcommand\Ci{ C_0^\infty }
\newcommand\rmor{ \,\,\,{\rm or}\,\,\, }
\newcommand\rmand{ \,\,\,{\rm and}\,\,\, }
\newcommand\kf{k_{\rm F}}
\newcommand\rhop{\rho_{\,\eps}}
\newcommand\rhom{\rho_{\,-\eps}}
\newcommand\const{{\rm const.\, }}
\newcommand\Tr{{\rm Tr}}
\newcommand\tr{{\rm tr}}
\newcommand\al{{\alpha}}
\newcommand\umt{\left\{u_1, m_1\right\},\cdots,\left\{u_t,m_t\right\}}
\newcommand\Ff{{\cal F}}
\newcommand\Ecal{{\cal E}}
\newcommand\be{\beta}
\newcommand\Hh{{\cal H}}
\newcommand\hN{{\hat N}}
\newcommand\Pp{{\cal P}}
\newcommand\Da{\downarrow}
\newcommand\mcut{m_c}
\newcommand\gamo{\gamma}
\newcommand\upa{\uparrow}
\newcommand\doa{\downarrow}
\newcommand\per{{\rm per}}
\newcommand\vu{\nu}
\newcommand\rrho{\rho}
\renewcommand\kappa\varkappa
\renewcommand\rho\varrho
\newcommand\AaFF{{\mathbb A}}
\newcommand{\id}{{\mathord{\;\rm d}}}


\newtheorem{thm}{THEOREM}[section]
\newtheorem{prop}{PROPOSITION}[section]
\newtheorem{lem}{Lemma}[section]
\newtheorem{cor}{COROLLARY}[section]
\newtheorem{mydef}{DEFINITION}[section]
\theoremstyle{definition}
\newtheorem{rem}{Remark}[section]


\begin{document}

%\markboth{\scriptsize{\version}}{\scriptsize{\version}}


\title{The Second Order Upper Bound for the Ground Energy of a Bose Gas \thanks{Partially supported
by NSF grants DMS-0757425, 0804279}
}

\author{Horng-Tzer  Yau  and  Jun  Yin \\ \\
Department of Mathematics, Harvard University\\
Cambridge MA 02138, USA \\ \\
\\}


%\title{ }

%\author{   \vspace{5pt}  Horng-Tzer  Yau  and  Jin  Yin 
%\vspace{-2pt}\small

%\\ \vspace{-2pt}\small
%\\ {\small Email: \texttt  {jyin@math.harvard.edu}} 

%\date{\small \version}

\maketitle


\begin{abstract} 
Consider $N$  bosons in a finite box $\Lambda= [0,L]^3\subset \mathbf R^3$
interacting via a two-body  smooth repulsive  short range potential. 
We construct a variational state which gives the following 
upper bound on the ground state energy
per particle 
\[
\overline\lim_{\rho\to0} \overline \lim_{L \to \infty, \, N/L^3 \to \rho} \left(
\frac{e_0(\rho)- 4 \pi a \rho }{(4 \pi a)^{5/2}( \rho )^{3/2} }\right)\leq \frac{16}{15\pi^2},
\]
where $a$ is the scattering length of the potential.  Previously, an upper bound 
of the form $C 16/15\pi^2$ for some constant $C > 1$ was obtained in \cite{ESY}.
Our result  proves the upper bound of the  prediction by Lee-Yang \cite{LYang} and  Lee-Huang-Yang \cite{LHY}. 


\end{abstract} 

\renewcommand{\thefootnote}{${\,}$}
%\footnotetext{ Work Partially supported by NSF grants DMS-0757425, 0804279}
\renewcommand{\thefootnote}{${\, }$}
\footnotetext{\copyright\,2009 by the author.
This paper may be reproduced, in its entirety, for non-commercial
purposes.}
\footnotetext{This version (November 2026) corrects a number of details in the
arguments of arXiv:0903.5347 and writes out the steps that were only sketched there
(Section~\ref{sec:lemma22} and Appendix~\ref{app:props}); the main result,
Theorem~\ref{mainthm}, is unchanged. A complete formalization of the proof in the
Lean~4 proof assistant accompanies this version.}
\renewcommand{\thefootnote}{\arabic{footnote}}





\bigskip

\noindent {\bf AMS 2000 Subject Classification:} 82B10

\bigskip

\noindent {\it Keywords:} Bose gas, Bogoliubov transformation, variational principle.



%\maketitle


%
%\newpage 
% \tableofcontents

\newpage

\section{Introduction}

The ground state energy  is a fundamental property of a quantum system and 
it has been intensively studied since the invention of the quantum mechanics. 
The recent progresses  in experiments for the 
 Bose-Einstein condensation have inspired re-examination of the theoretic foundation  
 concerning the Bose system and, in particular, its ground state energy.    
 In the low density limit, the leading term of the ground state energy per particle 
 was identified  rigorously by   Dyson (upper
bound) \cite{D} and Lieb-Yngvason (lower bound) \cite{LY} to be $4\pi a \varrho$, 
where $a$ is the scattering length of the two-body
potential and $\varrho$ is the density. The famous second order
correction to this leading term was first computed by Lee-Yang \cite{LYang} (see also Lee-Huang-Yang \cite{LHY} and  the recent paper 
by Yang \cite{Y} for results in other dimensions). To describe this prediction, 
we now fix our notations:  Consider $N$ interacting bosons in a finite box $\Lambda= [0,L]^3\subset \mathbf R^3$
with periodic boundary conditions.  
The two-body interaction  is given by a smooth nonnegative potential $V$ 
of fast decay. 
The Lee-Yang's prediction of the energy per particle 
up to the second order  is given by
\beq\label{nextorder}
e_0(\rho)  = 4\pi\varrho a \Big [1+
\frac{128}{15\sqrt{\pi}}( \varrho a^3)^{1/2}
 + \cdots  \Big ] \; .
\eeq
The approach by Lee-Yang \cite{LYang} is based on the pseudo-potential approximation \cite{HY, LHY}
and the ``binary collision expansion method'' \cite{LHY}. One can also obtain \eqref{nextorder} by performing the 
Bogoliubov \cite{B} approximation and then replacing the integral of the potential 
by its scattering length \cite{Landau}.  Another derivation of \eqref{nextorder} 
was later given  by Lieb \cite{L} using 
a self-consistent closure assumption  for the hierarchy of correlation functions.
%Bogoliubov \cite{B} is based on an approximation of
%the interacting Hamiltonian by a quadratic  one  consisting only of
%interactions between the condensate and the non-condenstate particle
%pairs with momenta $k, -k$.  


In the recent paper \cite{ESY},  the potential $V$ was replaced by $\lambda V_0$ for some fixed function $V_0$ and 
$\lambda$ is small.  A variational state was constructed to yield  the rigorous
upper bound 
\beq
e_0(\rho) \le 4\pi\varrho a \Big [1+
\frac{128}{15\sqrt{\pi}}( \varrho a^3)^{1/2}S_\lambda \Big ]
 + O(\varrho^2|\log \varrho| )  \;
\label{E-1}
\eeq
with $S_\lambda \leq 1+C\lambda$.
In the limit $\lambda \to 0$, one recovers the prediction of Lee-Yang \cite{LYang} and  Lee-Huang-Yang \cite{LHY}. 
The trial state in \cite{ESY}  does not have a
fixed number of particles, and is a state in the Fock space
with expected number of particles  $N$ (Presumably a  trial state with a  fixed number of particles
can be constructed with a similar idea).
The trial state in \cite{ESY} is similar to the trial state used by Girardeau and Arnowitt \cite{GA} 
and  recently by Solovej \cite{Sol}; it is of the form 
\beq\label{trialstate0}
 \exp \left [ |\Lambda|^{-1}  \sum_k  {c_k}  \cre_k\cre_{-k}\ann_0\ann_0  +  \sqrt{N_0}\cre_0  \right  ] |0\rangle
\eeq 
where  $c_k$ and $N_0$  have to be chosen carefully to give the correct asymptotic in energy.  This state captures  the idea that 
particle pairs of opposite momenta  are created from the sea of condensate consisting of  zero momentum particles.  It is believed that this type of trial state gives the ground state energy consistent with the 
Bogoliubov approximation. In the case of Bose gas,  the Bogoliubov approximation yields  
the correct energy up to  the order $\rho^{3/2}$, but the constant is correct only in the semiclassical limit---
consistent with the calculation using the trial state \eqref{trialstate0}.   
It should be noted that the Bogoliubov approximation gives the  correct ``correlation  energy"  in several 
setting including the one and  two component charged Boson gases  \cite{LS1, LS2, Sol} and 
the Bose gas  in  large density-weak potential limit \cite{GS}.



For the Bose gas in low density, the result of \cite{ESY} suggests to correct the error by 
renormalizing the propagator.  Unfortunately,  it is difficult to implement this idea. 
Our main observation is to relax the concept of condensates 
by allowing particle pairs to have nonzero total momenta. More precisely, we consider a trial state of the form 
\beqa\label{trialstate00}
&& \exp \Big  [ |\Lambda|^{-1}  \sum_k \sum_{v \sim \sqrt \rho }  2 \sqrt { \lambda_{k+v/2}  \lambda_{-k+v/2}} \cre_{k+v/2}\cre_{-k+v/2}\ann_v\ann_0   \nonumber \\ 
&& +  |\Lambda|^{-1}  \sum_k  {c_k}  \cre_k\cre_{-k}\ann_0\ann_0  +  \sqrt{N_0}\cre_0  \Big  ] |0\rangle
\eeqa
for suitably chosen $c$ and $\lambda$. 
Notice that the total momentum of the pair, $v$, is required to be of  order $\rho^{1/2}$ and the constant $2$ comes from the ordering of 
$\ann_v \ann_0$. We shall make further simplification that $\lambda_k = c_k$. Even with this simplification, 
however,  
this state  is still too complicated. We will extract some properties from this representation and define an 
$N$ particle trial state whose energy is given by the Lee-Yang's prediction up to the second order term.  Details will be given in Section \ref{trialstate}.  Our result shows that, in order to obtain the second order energy,  the typical ansatz for the Bogoliubov approximation should be extended to allow pair particles with nonzero  momenta. This also suggests that the Bogoliubov approximation has to be  modified in order to yield the correct energy of the low density  Bose gas to the second order. 

 
\section{Notations  and Main Results}

Let  $\Lambda= [0,L]^3\subset \R^3$ be a cube 
with periodic boundary conditions with the dual space  
$\Lambda^*:= (\frac{2\pi}{L} \Z)^3$. 
The Fourier transform is defined as
$$
      W_p:= \hat W(p)  = \int_{x\in \R^3} e^{-ipx} W(x) \rd x, \qquad
     W(x) = \frac{1}{|\Lambda|} \sum_{p\in \Lambda^*}e^{ipx}  W_p \; .
$$ 
Here we have used the convention to denote the Fourier transform of a function $W$ at the momentum $p$ by $W_p$ 
instead of $\wh W (p)$  to avoid too heavy notations. 
%The inversion formula is given by 
%\[
%     \frac{1}{|\Lambda|} \sum_{p\in \Lambda^*} e^{ipx} = \delta_{\R^3}(x), \qquad \int_\Lambda e^{ipx}\rd x = \delta_{\Lambda^*}(p) 
%\]
%where $\delta_{\R^3}(x)$ is the usual continuum delta function and the function 
%$\delta_{\Lambda^*}(p)=
%|\Lambda|= L^d$ if $p=0$ (otherwise it is zero) is the lattice delta 
%|function.
%
%{\it Remark:} For technical
%reason, one may introduce a discrete state space, then
%$\Lambda := [0,L]^3\cap (\e \Z)^3 = (\e \Z)^d/(L\Z)^3$
%and $\Lambda^*:= (\frac{2\pi}{L}\Z)^3/(\frac{2\pi}{\e}\Z)^3$
%and then all integrals become finite sums and one can later remove $\e\to0$. 
%This can easily be made rigorous and we neglect it in
%this discussion
%We will neglect the subscript, the argument indicates whether it is the momentum or position space delta function. In general 
%we will also neglect the hat in the Fourier
%transform. To avoid confusion, we follow the convention that the variables $x,y,z$ etc denote position space, the variables $p,q,r$ etc. momentum space. 
Since the summation of $p$ is always restricted to $\Lambda^*$, we will  not explicitly specify it. 
%$$
%\sum_p:=   \sum_{p\in \Lambda^*}
%$$
%i.e. momentum summation is always on the whole $\Lambda^*$.
%
%Notice that the choice of the $\delta_{\Lambda^*}$ ensures
%that in the $L\to\infty$ limit, this delta function converges to the usual continuum delta function $\delta(p)$ in momentum space with respect to the measure $\rd^3 p/(2\pi)^3$: $$
%       \lim_{L\to\infty}
%       \frac{1}{|\Lambda|}\sum_{p\in\Lambda^*} F(p) \delta_{\Lambda^*}(p-q)
%       = \int_{\R^3} \frac{\rd^3p}{(2\pi)^3} F(p)\delta(p-q)
%     = F(q)
%$$
%(where $\delta(p)$ is defined by the last formula).

%

%
%We consider the canonical
%  bosonic operators, $\ann_k \cre_k$, for any $k\in\Lambda^*$, with the usual CCR: $$
%      [\ann_p, \cre_q] = \ann_p \cre_q -\cre_q\ann_p = \delta(p-q)
%     = \left\{
%\begin{array}{ll}
%  L^d & \mbox{ if } p=q \\ 
%0 & \mbox{ otherwise.}
%\end{array} 
%\right.
%$$
%Then the Hamiltonian is given by
%$$
%      H = \frac{1}{|\Lambda|}
%\sum_p p^2 \cre_p\ann_p + \frac{1}{|\Lambda|^3} \sum_{p,q,u}
%     V_u \cre_p\cre_q\ann_{p-u}\ann_{q+u}
%$$
%
%However, it will be more convenient to redefine the bosonic operators as $$
%    \ann_k\to \frac{1}{\sqrt{|\Lambda|}} \;\ann_k, 
%\qquad \cre_k\to \frac{1}{\sqrt{|\Lambda|}}\; \cre_k,
%$$
%(without changing the notation)
%i.e. from now on we assume that

We will use the bosonic operators with the commutator relations
$$
      [\ann_p, \cre_q] = \ann_p \cre_q -\cre_q\ann_p
     = \left\{
\begin{array}{ll}
  1 & \mbox{ if } p=q \\ 
0 & \mbox{ otherwise.}
\end{array} 
\right.
$$
The two body interaction is given by  a smooth, symmetric (i.e., $V(-x)=V(x)$)  non-negative function $V(x)$  of fast decay.
Clearly, in the Fourier space, we have 
$    V_u=V_{-u}= \bar V_u$. 
Furthermore,  we assume that the potential $V$ is small so that the Born series converges.
The  Hamiltonian of the many-body systems with the potential $V$ and the periodic boundary condition is thus  given by
\beq
      H = 
\sum_p p^2 \cre_p\ann_p + \frac{1}{|\Lambda|} \sum_{p,q,u}
     V_u \cre_p\cre_q\ann_{p-u}\ann_{q+u}
\label{ham}
\eeq
In position space, $H$ acts on symmetric functions of $x_1, \ldots, x_N\in \Lambda$ as
\[
H=\sum_{i=1}^N-\Delta_i+\sum_{i\neq j}V^P(x_i-x_j), \qquad V^P(x)=\sum_{n\in\Z^3}V(x+nL),
\]
where $\Delta_i$ has periodic boundary conditions; the periodization $V^P$ of $V$ has the Fourier coefficients
$\int_\Lambda e^{-iux}V^P(x)\rd x=V_u$, $u\in\Lambda^*$.
Thus there is no factor $1/2$ in front of the interaction in \eqref{ham}: each pair of particles is counted twice,
in accordance with the scattering equation below.

Let $1-w$ be the zero energy scattering solution 
$$
    -\Delta(1-w) +  V(1-w)=0
$$
with $0\leq w<1$ and $w(x)\to 0$ as $|x|\to \infty$.
Then the scattering length is given by the formula 
$$
   a: =\frac{1}{4\pi}\int_{\R^3} V(x)(1-w(x))\rd x
$$
Introduce $g_0$, whose meaning will be explained later on,  to denote the  quantity
\[
g_0 = 4 \pi a.
\]
%We will give a meaning of $g_0$ later on.  
%and $w(x)\sim 1-a/|x|$ when $|x|\geq R_0$.

Let $\mathcal H_N$ be the Hilbert space of  $N$ bosons. Denote by $\rho_N=N/|\Lambda|$ the density of the system.
The ground state energy of the Hamiltonian \eqref{ham} in  $\mathcal H_N$  
is given by 
\[
E^{P}_0(\rho,\Lambda) = \inf spec H_{\mathcal H_N}
\]
and the ground state energy per particle is $e_0^P(\rho,\Lambda) =  E_0^P(\rho,\Lambda)/N$. We can also consider other boundary 
conditions, e.g., $e_0^D(\rho,\Lambda)$ is the Dirichlet boundary condition ground state energy per particle,
i.e., the ground state energy per particle of $\sum_{i=1}^N-\Delta_i+\sum_{i\neq j}V(x_i-x_j)$ on symmetric functions
with Dirichlet boundary conditions on $\Lambda$.


In this paper, we will always take the limit $L \to \infty$ so that the density $\rho_N \to \rho$ 
for some fixed density $\rho$. {F}rom now on,  we will use $\lim_{L \to\infty}$  for the more complicated notation 
$\lim_{L \to \infty, \, N/L^3 \to \rho}$. We now state the main result of this paper. 


\begin{thm}\label{mainthm} 
Suppose the potential $V$ is  smooth, symmetric, nonnegative with fast decay
and sufficiently small so that the Born series converges.\footnote{We also assume that $V$ is not identically zero. Then
$g_0=4\pi a=\int V(1-w)>0$, and the quotient below is defined.}
Then the ground state energy per particle  satisfies the upper bound
\beq
\overline\lim_{\rho\to0} \overline \lim_{L \to\infty}\left(
\frac{e_0^P(\rho,\Lambda)- g_0 \rho }{g_0^{5/2} \rho ^{3/2} }\right)\leq \frac{16}{15\pi^2}
\eeq
\end{thm}
Although we state the theorem in the form of limit $\rho \to 0$, an error bound  is available from the proof. 
We avoid stating such an estimate to simplify the notations and proofs. 
Our result  holds also for Dirichlet boundary condition.

\subsection{Reduction to Small Torus with Periodic Boundary Conditions}

To prove Theorem \ref{mainthm}, since $e_0^P(\rho,\Lambda)\le \langle H_N\rangle_\Psi N^{-1}$ by the variational principle,
we only need to construct a trial state $\Psi(\rho, \Lambda)$ satisfying  the boundary condition and
\beq\label{resultpro1}
\overline\lim_{\rho\to0}\overline\lim_{L\to\infty}\left(
\frac{\langle H_N\rangle_\Psi N^{-1}-g_0\rho }{ g_0^{5/2} \rho ^{3/2}}\right)\leq \frac{16}{15\pi^2}
\eeq
%
The first step is to construct a trial state with a Dirichlet boundary condition in a cube  of order slightly bigger than $\rho^{-1}$. 
%It is easy to see the trial state in the thermodynamic limit can be construct with the trial state in the \textit{small} box, i.e., the following propostion does imply our main result theorem one.
\begin{lem}\label{thm1}
For density $ \rho$ small enough, there exist $ L \sim  \rho^{\,-25/24}$ and 
 a trial state $\Psi$ of $ N$ $( N= \rho  L^3)$ particles on $ \Lambda = [0,   L]^3$ satisfying  the Dirichlet boundary condition and 
 \beq
 \overline\lim_{ \rho\to0}\left(
\frac{\langle H_{ N}\rangle_\Psi  { N}^{-1}-g_0 \rho }{ g_0^{5/2}  \rho ^{\,3/2}}\right)\leq \frac{16}{15\pi^2}.
 \eeq 
% such that 
%the energy per particle  w.r.t. this trial state can be bounded above by 
%\beq
%\overline\lim_{\rho\to0}\left(
%\frac{\langle H_N\rangle_\Psi N^{-1}-g_0\rho }{ (g_0\rho)^{3/2}}\right)\leq \frac{16}{15\pi^2}
%\eeq
\end{lem}
Once we have a trial state with the Dirichlet boundary condition, we can duplicate it so that a trial state can be constructed for cubes
with linear dimension $  \gg  \rho^{-25/24}$. This proves Theorem \ref{mainthm}. The duplication is carried out in
Section \ref{subsec:dup}, where we also specify the rate of decay of $V$ that it uses.

The next lemma shows that a Dirichlet boundary condition trial state with correct energy can be obtained from a   
periodic one.

\begin{lem}\label{lemma1}
Recall the ground state energies  per particle  $e^{D}_0(\rho,\Lambda)$ and $e^{P}_0(\rho,\Lambda)$ for 
the Dirichlet and periodic boundary condition. 
Let   $\Lambda = [0, L]^3$ and  $L = \rho^{-25/24}$.\footnote{Since the number of particles $N=\rho L^3=\rho^{-17/8}$
is an integer, here and in Theorems \ref{lemma2} and \ref{lemma2.1} the density runs through the values
$\rho=N^{-8/17}$, $N\in\N$, and $\rho\to0$ means $N\to\infty$.} Suppose  the energy for the periodic boundary condition
satisfies that 
\beq
\overline\lim_{\rho\to0}\left(
\frac{e^P_0(\rho,\Lambda)-g_0\rho }{ g_0^{5/2} \rho ^{3/2}}\right)\leq \frac{16}{15\pi^2}
\eeq
Then for $\widetilde\Lambda=[0, \widetilde L ]^3$, $\widetilde L=L(1+2\rho^{25/48})$ and $\widetilde \rho=\rho L^3/\widetilde L^3$, the following estimate for the energy of the Dirichlet boundary condition  holds: 
\beq\label{resultlemma1}
\overline\lim_{\rho\to0}\left(
\frac{e^D_0(\widetilde \rho,\widetilde\Lambda)-g_0\widetilde \rho }{ g_0^{5/2} \widetilde \rho ^{\,3/2}}\right)\leq \frac{16}{15\pi^2}
\eeq
%Note: $\widetilde L\sim\widetilde \rho^{\,-25/24}$
\end{lem}


The construction of  a periodic trial state yielding the correct energy upper bound is the core of this paper. 
We state it as the following theorem. 
 
\begin{thm}\label{lemma2}
There exists a periodic trial state $\Psi$ of $N$ particles on $\Lambda = [ 0, L]^3$,  $L = \rho^{-25/24}$ such that ($N=|\Lambda|\rho$) 
%\eqref{resultpro1} holds. 
%the energy per particle for  periodic boundary condition can be bounded above by  
\beq
\overline\lim_{\rho\to0}\left(
\frac{\langle H_N\rangle_\Psi N^{-1}-g_0\rho }{ g_0^{5/2}\rho^{3/2}}\right)\leq \frac{16}{15\pi^2}
\eeq
\end{thm}


This paper is organized as follows: In Section 3, we define rigorously the trial state and state the estimate
on its energy, Theorem \ref{lemma2.1}.
In Section 4, we outline the Lemmas needed to prove Theorem \ref{lemma2.1}, and deduce Theorem \ref{lemma2} from it.
In Section 5, we estimate the number of particles in the condensate and various momentum regimes.
These estimates are the building blocks for all other estimates later on. In Section 6, we estimate the kinetic energy. The potential energy is estimated in Sections 7--11. Finally in Section 12, we prove the reduction
to the periodic boundary condition, i.e., Lemma \ref{lemma1}, derive Lemma \ref{thm1} from Theorem \ref{lemma2} and
Lemma \ref{lemma1} (by the variational principle, Theorem \ref{lemma2} implies the assumption of Lemma \ref{lemma1}),
and prove Theorem \ref{mainthm} by duplicating the trial state of Lemma \ref{thm1} (Section \ref{subsec:dup}).
The proofs of Propositions \ref{prop1} and \ref{prop2} are given in Appendix \ref{app:props}.




\section{Definition of the Trial State}\label{trialstate}

We now give a formal definition of the trial state. This somehow abstract definition will be explained later on.  
We first identify  four regions in the momentum space  $\Lambda^*$  which are  relevant to the construction of the trial state: $P_0$ for the condensate,  $P_L$ for the low momenta, which are of the order  $\rho^{1/2}$;  $P_H$ for momenta of order one, and $P_I$ the region between 
$P_L$ and $P_H$. 

\begin{mydef} \label{def1} Define four subsets of momentum space: $P_0$, $P_L$, $P_I$ and $P_H$.
\beqa\nonumber
P_0&&\equiv\left\{p=0\right\}\\\nonumber
P_L&&\equiv\left\{p\in \Lambda^*| \eps_L\rho^{1/2}\leq|p|\leq \eta^{-1}_L\rho^{1/2}\right\}\\\nonumber
P_I&&\equiv\left\{p\in \Lambda^*|\eta^{-1}_L\rho^{1/2} <   |p|\leq \eps_H\right\}\\
P_H&&\equiv\left\{p\in \Lambda^*|\eps_H <   |p|\right\}\, , 
\eeqa
where the parameters are chosen so that  
\beq\label{defepseta}
\eps_L, \eta_L, \eps_H\equiv\rho^{\eta}\rmand \eta\equiv1/200
\eeq
Denote by $P = P_0\cup P_L \cup P_I \cup P_H$.
\end{mydef}
We remark that the momenta between $P_0$ and $P_L$ are irrelevant to our construction. 
Next, we need a notation for the collection of states with $N$ particles. 
\begin{mydef}
Let  $\widetilde M$ be  the set of all functions $\al: P \rightarrow \N\cup 0$ such that
\beq
\sum_{k\in P}\al(k)=N
\eeq
For any $\al \in \widetilde M$,  denote by $|\al\rangle\in \mathcal H_N$  the unique state (in this case, an  $N$-particle  wave function) 
defined by the map $\alpha$ 
\[
|\al\rangle =  C \prod_{k\in P} (a^\dagger_k)^{\alpha(k)} | 0 \rangle\, , 
\]
where the positive constant $C$ is chosen so that $|\al\rangle $ is $L_2$ normalized.  
Define $\al_{free}$ as $\al_{free}(k)=N\delta_{0,k}$.
\end{mydef}

Clearly, we have 
\beq\label{number}
a^\dagger_k a_k|\al\rangle=\al(k)|\al\rangle,\,\,\, \forall k\in P
\eeq



%We need notations connecting two wave functions. 

% related by two operations: 1. Creating 
%a pair of particles with momenta $k$ and $-k$. 2. Creating two high  momentum particles whose 
%total momentum is small, but not zero. The precise definition is given in the following definition. 

%
%\begin{mydef}{$\mathcal A^k$ and $\mathcal A^{u,\,k'}$}
%\par
%\begin{enumerate}
%	\item We say $\be=\mathcal A^k\al$, if $\al, \be\in \mathcal B_N$  {\bf or we use $M$?} and $$|\be\rangle =\const \cre_k\cre_{-k}\ann_0\ann_0|\al\rangle$$.
%\item For $(u,k)$ satisfying  $u\in P_L$ and $\pm k+u/2\in P_H$, we say $\beta=\mathcal A^{u,\,k}\al$, if $\al, \be\in \mathcal B_N$ and $$|\be\rangle =\const \cre_{k+v/2}\cre_{-k+v/2}\ann_v\ann_0|\al\rangle$$.
%\end{enumerate}
%\end{mydef}


\begin{mydef}
We define two  relations between functions in $\widetilde M$: 
\begin{enumerate}
	\item Strict pair creation of momentum $k$: Denote by  $
\be: =  \mathcal A^k \al $
 if $\beta$ is generated by creating 
a pair of particles with momenta $k$ and $-k$, i.e., 
\beq	
\be(p)= \left\{ 
\begin{array}{ll}
\al(p)-2,& p=0 \\ 
\al(p)+1, & p=\pm k\\
\al(p), & others
\end{array}\right.
\eeq
In terms of states, we have 
\[
|\be\rangle =  C a^+_k a^+_{-k} a_0^2 | \alpha\rangle 
\]
where $C$ is a positive constant so that the state $|\be\rangle$ is normalized. 


\item Soft pair creation with total momentum $u$ and difference $ 2 k$:   
Denote by  $\beta=\mathcal A^{u,\,k}\al$ if $\beta$ is generated by creating two  particles with high momenta $\pm k+u/2\in P_H$ 
so that the total momentum $u$ is in $P_L$, i.e., 
\beq	
\be(p)= \left\{ 
\begin{array}{ll}
\al(p)-1,& p=0\,\,\, {\rm or}\,\,\, u \\ 
\al(p)+1, & p=\pm k+u/2\\
\al(p), & others
\end{array}\right.
\eeq
Notice that $\mathcal A^{u,\,k}\al$ is defined only if $\pm k+u/2\in P_H$. 
In terms of states, we have 
\[
|\be\rangle = C a^+_{k+ u/2} a^+_{-k + u/2} a_0 a_u | \alpha\rangle 
\]
where $C$ is the normalization constant. 
Since $\beta(p)$ 
has to be  nonnegative,  the state  $\mathcal A^{k}\al$ or  $\mathcal A^{u,\,k}\al$ is not defined for all $\alpha$ or $k, u$. 
\end{enumerate}
Define $D_\alpha$ to be  the set all possible derivations of $\alpha$ from the previous two operations: 
\beq
D_\alpha=\left\{\mathcal A^{u,\,k}\al \in \widetilde M \right\}\cup\left\{\mathcal A^{k}\al \in \widetilde M\right\}
\eeq
\end{mydef}


Our trial state will be of the form  $\sum_{\al\in \widetilde  M}f(\al)|\al\rangle$ where $f$ is  supported in a 
subset of $\widetilde M$ which we now define. 


\begin{mydef}\label{M}
Fix a large real number $k_c$. 
We define $M$ as the smallest subset of $\widetilde M$ such that
\begin{enumerate}
	\item $\al_{free} \in M$.
	\item  $M$ is closed under strict pair creation provided the momentum  $u\in P_I\cup P_H$, i.e.,  if  $\al \in M$ and $\mathcal A^{u}\al\in \widetilde M$ then $\mathcal A^{u}\al\in M$.
	\item $M$ is closed under strict pair creation provided the momentum  $u\in P_L$  and  $\max\{\al(u), \al(-u)\}<\mcut$, i.e., 
	if  $\al \in M$ and $\mathcal A^{u}\al\in \widetilde M$, then $\mathcal A^{u}\al\in M$. Here we choose $m_c$ as 
	\beq\label{defmc}
	m_c\equiv\rho^{-\eta}=\rho^{-1/200}
	\eeq
	\item   M is closed under soft pair creation from states with perfect pairing of momenta $u$ and $-u$. More precisely, 
	for $u \in P_L$ with $\al(u)=\al(-u)$, if $\al \in M$, $\mathcal A^{u,\,k}\al\in \widetilde M$ and 
	$$ \e_H < |\pm k+u/2|\leq k_c,%\,\,\,\,\,\,|w_{\pm k+u/2}|\geq \rho^{1/3},
	$$ 
	then $\mathcal A^{u,\,k}\al\in M$. 
\end{enumerate}
\end{mydef}
The set $M$ is unique since the intersection of two such sets $M_1$ and $M_2$ satisfies all four conditions. 
%Clearly,  for any $\al \in M$,  there exists $\al_1$, $\al_2$, $\cdots$ $\al_m$ such that $\al_{i+1}\in D_{\al_{i}}$ (for $1\leq i<m$) and $\al_1=\al_{free}$, $\al_m=\al$.

For any $u\in P_L$, we define the set of states with symmetric (asymmetric resp.)  pair particles of momenta $u, -u$ 
by  $M_u^s$ ($M_u^a$ resp.): 
\beqa\label{defMsau}
&&M_u^s\equiv\{\al\in M|\al(u)=\al(-u)\}\\\nonumber
&&M_u^a\equiv\{\al\in M|\al(u)\neq\al(-u)\}.
\eeqa
Denote by  $\al^*(u)$  the maximum of $\al(u)$ and $ \al(-u)$:
\beq\label{def*}
\al^*(u)=\max\{\al(u), \al(-u)\}
\eeq 
Since soft pair creation was allowed only from momenta in $P_L$ and the final momenta are in $P_H$, we have 
\[
\al^*(u)- \al(u)\in \{0, 1 \} \text{ for all }  u \in P_L, \quad  \al(-u)= \al(u) \text{ for all }  u \in P_I
\]


Before defining the weight $f(\alpha)$, we introduce several quantities related 
to the scattering equation. In the momentum space, the  scattering equation is 
given by ($p\in \R^3$)
%
%Then we have 
%\beq
%     4\pi a = V_0 - \int V_p w_p  \frac{\rd^3p}{(2\pi)^3}
%     = V_0 - \lim_{L\to\infty}\frac{1}{|\Lambda|}\sum_p V_pw_p  
%     \eeq
%      and 
\beq\label{av}
   -p^2 w_p+ V_p -\int_{\R^3} V_{p-r} w_r\,\frac{\rd^3 r}{(2\pi)^3}=0, \qquad \forall p \neq 0
\eeq 
 Let $g$ be the function 
\beq
     g(x): = V(x) (1- w(x))
\label{gdef}
\eeq
%in Fourier space:
%$$
 % f_p = (V\ast w)_p = \frac{1}{|\Lambda|}
 % \sum_r V_{p-r} w_r
%$$
Then the scattering equation in momentum space takes the form
\beq\label{gp}
     g_p= p^2w_p  \qquad \forall p \neq 0
\eeq
(Since $w(x)$ decays only like $a|x|^{-1}$, $w$ is not integrable; $w_p$ denotes its Fourier transform
in the sense of distributions, which is the locally integrable function $w_p=g_p|p|^{-2}$, $p\neq 0$.)
%In particular, from \eqref{av}
One can check $ 4\pi a =g_0$  %\label{ag} and 
this explains the notation $g_{0}$ used in Theorem \ref{mainthm} and Theorem \ref{lemma2}. 
%If $V(x)$ is a nice function, then so is $w(x)$, so for simplicity we assume in the sequel that 
%\beq
%    V_p, g_p  \quad \mbox{are bounded, continuous, sufficiently 
%decaying}
%\label{cond}
%\eeq




\begin{mydef}
Define  for all $ \eps\neq 0$
\beq\label{defrho0}
\rhop\equiv\rho_0+\eps\rho^{3/2}\, , \, \rho_0:= \rho-\frac{1}{3\pi^2}(g_0)^{3/2}\rho^{3/2}, 
\eeq
where  $\rho_0$ will be the approximate density of the condensate.
Define the ``chemical potential"   $\lambda$ by 
\beq\label{lambda}
\lambda_k= \left\{ 
\begin{array}{ll}
\frac{1-\sqrt{1+4\rho g_0|k|^{-2}}}{1+\sqrt{1+4\rho g_0|k|^{-2}}}\,\rho^{-1},& k\in P_L \\ 
-w_k, & k\in P_I\cup P_H
\end{array}\right.
\eeq
\end{mydef}
One can check that, to the leading order, $\lambda$ is given by 
\beq\label{lambda0}
\rho\lambda_k\equiv\frac{1-\sqrt{1+4\rho g_k|k|^{-2}}}{1+\sqrt{1+4\rho g_k|k|^{-2}}}
\eeq
Notice that $\lambda_k$ is real number and can be negative.

\begin{mydef}{\bf The Trial State}
\par Let  $\Psi$ be defined by 
\beq
\Psi\equiv\sum_{\al\in M}f(\al)|\al\rangle
\eeq
where the coefficient $f$ is given by 
 \beq\label{deffal}
%f(\al)=f(\al_{free})\sqrt{\frac{N!/\al(0)!}{|\Lambda|^{N-\al(0)}}}\prod_{k\in P_I\cup P_H}{(\lambda_{k})^{\frac{\al(k)}{2}}}\prod_{u\in P_L}{(\lambda_{u})^{\frac{\al^*(u)}{2}}}\prod_{u:u\in P_L}^{\al\in M^u_a }\left(\frac{4\al^*(u)}{|\Lambda|}\right)^{1/4}
f(\al)=C_N %\sqrt{\frac{N! }{|\Lambda|^{N}} }
\sqrt { \frac  {|\Lambda|^{\alpha(0)}} {\al(0)!} }\prod_{k\neq 0}(\sqrt{\lambda_{k}})^{\al(k)}\prod_{u\in P_L, \alpha^\ast (u) - \alpha (u) = 1}\sqrt{\frac{4\al^*(u)\lambda_u}{|\Lambda|}}
\eeq
Here we follow the convention $\sqrt{x}=\sqrt{|x|}i$ for $x<0$. For convenience, we define $f(\al)=0$ for $\al\notin M$.
The constant  $C_N$ is  chosen so that $\Psi$ is $L_2$ normalized,\footnote{Such a constant exists: with $C_N=1$,
the sum $\sum_{\al\in M}|f(\al)|^2$ contains $|f(\al_{free})|^2>0$, and it is finite. Indeed,
$N!/\al(0)!\le N^{N-\al(0)}$ and $\al^*(u)\le N$ give
$|f(\al)|^2\le |f(\al_{free})|^2\prod_{k\neq 0}(\rho|\lambda_k|)^{\al(k)}\prod_u 4\rho|\lambda_u|$, the last product
over $u\in P_L$ with $\al^*(u)-\al(u)=1$; soft pair creations produce momenta of size at most $k_c$, so
$\al(k)=\al(-k)\le N$ for $|k|>k_c$; and $\sum_{|k|>k_c}\rho^2|\lambda_k|^2<\infty$, since there
$|\lambda_k|=|w_k|\le g_0|k|^{-2}$ by \eqref{gp} and $|g_k|\le\int g=g_0$ ($g\ge0$). Since $\al$ takes
values in $\{0,\ldots,N\}$ on the finite set $\{0<|k|\le k_c\}$, the sum is at most a finite constant times
$\prod_{\{k,-k\},\,|k|>k_c}\sum_{n=0}^{N}(\rho|\lambda_k|)^{2n}<\infty$.} i.e.,
\[\left\langle \Psi|\Psi\right\rangle=1.\]  
\end{mydef}



\begin{thm}\label{lemma2.1} Suppose  $\Lambda = [ 0, L]^3$ and $L = \rho^{-25/24}$. 
Then the  trial state $\Psi$ in \eqref{deffal} satisfies the  estimate 
\beq
\overline\lim_{k_c \to \infty} \overline\lim_{\rho\to0}\left(
\frac{\langle H_N\rangle_\Psi N^{-1}-g_0\rho }{ g_0^{5/2}\rho^{3/2}}\right)\leq \frac{16}{15\pi^2},
\eeq
where $k_c$  is given in Definition \ref{M}. We recall that $m_c^{-1}, \e_L, \eta_L, \e_H$ are chosen as a small power of $\rho$ in \eqref{defepseta} and \eqref{defmc}.
\end{thm}
Theorem \ref{lemma2} follows from Theorem \ref{lemma2.1} by choosing the cutoff $k_c$ as a function of $\rho$;
see the end of Section 4.
%The last factor in \eqref{deffal} can be understood as follows: We have  $\al^*(u)$  of pairs of particles. Now we choose one of them can make a soft pair creation with probability $\lambda_u$. The entropy factor to to choose from these $\al^*(u)$ pair gives the factor 
%$\al^*(u)$. Since we can still choose one from each pair, we have 
%the factor $2$  


\subsection{Heuristic Derivation of the Trial State}

We now give a heuristic idea for the construction of  the trial state. Fix an ordering of momenta in $\Lambda^*$ so that the first one 
is the zero momentum.  We will  use the occupation number representation so that 
\beq
|n_1,\,\,n_2,\,\,\cdots\rangle
\eeq
represents the normalized state with  $n_i$ particles of momentum $k_i$. For example, 
\[
|N,\; 0, 0 ,\,\cdots\rangle = \frac 1 { \sqrt {N!}}( \cre_0)^N |0\rangle 
\]
%Let  $\mathcal B$ denote the set of all these states   and $\mathcal B_m$  consisting of those with total number of particles $m$, i.e., 
%\beq\label{defB}
%\mathcal B\equiv \big\{|n_1,\,\,n_2,\,\,\cdots\rangle:\,\,\,n_i\in \Z,\,\,\,n_i\geq 0\big\},\,\,\,\mathcal B_m\equiv\mathcal H_m\cap\mathcal B
%\eeq
%For any state $|\al\rangle\in\mathcal B$ and for any $k\in\Lambda^*$, let  
%$\alpha(k)$  denote the expected number of particle with momentum $k$: 
%\beq\label{defal}
%\al(k)=  \langle \al |  \cre_k\ann_k|\al\rangle =  \langle   \cre_k\ann_k \rangle_{\alpha}
%\eeq
Recall that we would like to generate a state of the form in \eqref{trialstate00}. A slightly modified one is 
\beqa\label{trialstate001}
&& \exp \Big  [ |\Lambda|^{-1}  \sum_k \sum_{v \sim \sqrt \rho }  2 \sqrt {\lambda_{k+v/2}  \lambda_{-k+v/2}} \cre_{k+v/2}\cre_{-k+v/2}\ann_v\ann_0   \nonumber \\ 
&& + |\Lambda|^{-1} \sum_k  {\lambda_k}  \cre_k\cre_{-k}\ann_0\ann_0   \Big  ] |N,\; 0, 0 ,\,\cdots\rangle
\eeqa

We now expand the exponential and require that $\cre_{k+v/2}\cre_{-k+v/2}\ann_v\ann_0$ to  appear at most once. 
The rationale of this assumption is that the soft pair creation is a rare event and thus we can neglect higher order terms.  
Our trial state is thus a sum  of the following  state parametrized by  $k_1,\cdots,k_s$, $n_1,\cdots,n_s$, $k'_1,\cdots,k'_t$
and $v_1,\cdots,v_t$: 
\beq\label{defmn}
 \const \prod_{j=1}^t    \sqrt { 4 \lambda_{k'_j+v_j/2}  \lambda_{-k'_j+v_j/2}}  \prod_{i=1}^s\left(\lambda_{k_i} \right)^{n_i} |\al\rangle
\eeq
where 
\beqa
|\al\rangle = && \const |\Lambda|^{-t-\sum_{i=1}^s n_i}  \prod_{j=1}^t    \cre_{\frac{v_j}{2}+k'_j}\cre_{\frac{v_j}{2}-k'_j}\ann_{v_j}  \ann_0  \nonumber   \\
&& \times \prod_{i=1}^s   \frac 1 {n_i !} \left ( \cre_{k_i}\cre_{-k_i}\ann_{0}\ann_0\right)^{n_i}|N,0,\cdots\rangle
\eeqa
Here we have chosen the constant so that the norm of $|\al\rangle$ is one. 
We  also require that  $v_i+v_j\neq0$ for $1\leq i,j\leq t$ since $v_i+v_j= 0$  is a higher order event. 

%Recall the definition of $P$ for momenta in $\Lambda^*$. 
We  further make  the simplifying  assumption that $v_i \in P_L$.  
Observe now that the state $|\al\rangle$ can be obtained from strict and soft pair creations. This explains the core  idea behind the definition of  $M$ in Definition \ref{M}. Other restrictions in the definition were mostly due to various cutoffs needed in the estimates. 
Finally, up to factors depending only on $\Lambda$ and $N$, the coefficient  in \eqref{defmn} gives $f(\alpha)$ in \eqref{deffal}. Notice all factors depending on $s, t, n_i$ were already included in $|\alpha\rangle $.




The choice of $\lambda$ is much more complicated. To the first approximation, $\lambda$ can be obtain
from the work of \cite{ESY}. We thus use this choice to identify the error terms.  Once this is done, we
optimize the main terms and this leads to the current definition of $\lambda$. Notice that, since our trial state
is different, there are more main terms than in \cite{ESY}.

%We know that $\rho \lambda_k$ is near the $c_k$ 
%in the paper by \cite{ESY}. Here the factor $\rho$ appears since our normalization is slightly different from \cite{ESY}. 
%The basic idea to get the correct $ \lambda_k$ is first to use the choice of \cite{ESY} to identify the main terms
%of the energy. Since our trial state is different, there will be more main terms than in \cite{ESY}. Collecting these main terms 
%and optimizes the choice of $\lambda_k$ gives the choice  \eqref{lambda}. More details will be given later on. 




\section{Proof of Theorem \ref{lemma2}}
\begin{proof}
We first prove
\beq\label{Maindesired}
\overline
\lim_{k_c\to\infty}
  \left(\overline\lim_{\rho\to0}
     \left(
         \frac{|\Lambda|^{-1}\langle H\rangle_\Psi-g_0\rho^2}
         {\rho^{5/2}}
      \right)
  \right)
\leq \frac{16}{15\pi^2}g_0^{5/2}
\eeq
Here $g_0=4\pi a$, $\langle H\rangle_\Psi=\langle\Psi|H|\Psi\rangle$. We decompose the Hamiltonian as follows:
\beq
H=\sum_{i=1}^N-\Delta_i+H_{S1}+H_{S2}+H_{S3}+H_{A1}+H_{A2},
\eeq
where 
\begin{enumerate}
	\item $H_{S1}$ is the part of interaction that annihilates two particles and creates the same two particles, i.e.,
	\beq
	H_{S1}= |\Lambda|^{-1} \sum_{u}V_{0}a^\dagger_u a^\dagger_u a_u a_u+ |\Lambda|^{-1} \sum_{u\neq v}(V_{u-v}+V_{0})a^\dagger_u a^\dagger_v a_u a_v
	\eeq
	\item $H_{S2}$ is the interaction between the condensate and strict pairs, i.e., 
	\beq
 H_{S2}=	|\Lambda|^{-1}\sum_{u\neq 0}V_{u}a^\dagger_u a^\dagger_{-u} a_0 a_0+C.C.
	\eeq
	\item $H_{S3}$ is the part of interaction that strict pairs are involved, i.e., 
	\beq
	H_{S3}=|\Lambda|^{-1} \sum_{u, v\neq 0, u\neq \pm v} V_{u-v}a^\dagger_u a^\dagger_{-u} a_v a_{-v}
	\eeq
	\item $H_{A1}$ is the part of the  interaction that one and only one condensate particle is involved i.e., 
	\beq
	H_{A1}=|\Lambda|^{-1} \sum_{v_1,v_2,v_3 \neq 0,\; v_1=v_2+v_3}2V_{v_2}a^\dagger_0 a^\dagger_{v_1} a_{v_2} a_{v_3}+C.C.
	\eeq
	\item $H_{A2}$ is the part of the  interaction which is not counted in $H_{S1}$ and there is no condensate nor strict pair involved i.e., 
	\beq
	H_{A2}= |\Lambda|^{-1} \sum_{\substack{v_i \neq 0,\; v_1+v_2=v_3+v_4\neq 0\\ \{v_1,v_2\}\neq \{v_3,v_4\}}}V_{v_1-v_3}a^\dagger_{v_1} a^\dagger_{v_2} a_{v_3} a_{v_4}
	\eeq
\end{enumerate} 
The estimates for the energies of these components are stated as the following lemmas, which will be proved 
in later sections. 

\begin{lem}\label{lemkinetic}The total kinetic energy is bounded  above by
\beq\label{proofthem1-1}
\begin{split}
&\overline\lim_{k_c\to\infty}\overline\lim_{\rho\to0}
\left(\frac{1}{|\Lambda|}\left\langle\sum_{i=1}^N-\Delta_i\right\rangle_{\Psi}-\rho_0^2 \|\nabla w\|_2^2\right)\rho^{-\frac{5}{2}}\\
&\qquad\leq \frac{4\|\nabla w\|_2^2g_0^{3/2}}{3\pi^2}  -\frac{8g_0^{5/2}}{5\pi^2}
\end{split}
\eeq
\end{lem}
\begin{lem}\label{lemHS1}
The expectation value of $H_{S1}$ is bounded  above by,
\beq\label{proofthem1-2}
\overline\lim_{k_c\to\infty}\overline\lim_{\rho\to0}
\left(\frac{1}{|\Lambda|}\left\langle H_{S1}\right\rangle_{\Psi}-\rho_0^2 V_0\right)\rho^{-5/2}
\leq \frac{4V_0g_0^{3/2}}{3\pi^2}
\eeq
\end{lem}
\begin{lem}\label{lemHS2}
The expectation value of $H_{S2}$ is bounded above by,
\beq\label{proofthem1-3}
\overline\lim_{k_c\to\infty}\overline\lim_{\rho\to0}
\left(\frac{1}{|\Lambda|}\left\langle H_{S2}\right\rangle_{\Psi}+2\rho_0^2 \|Vw\|_1\right)\rho^{-5/2}
\leq \frac{2V_0g_0^{3/2}}{\pi^2}
\eeq
\end{lem}
\begin{lem}\label{lemHS3}
The expectation value of $H_{S3}$ is bounded above by,
\beq\label{proofthem1-4}
\overline\lim_{k_c\to\infty}\overline\lim_{\rho\to0}
\left(\frac{1}{|\Lambda|}\left\langle H_{S3}\right\rangle_{\Psi}-\rho_0^2 \|Vw^2\|_1\right)\rho^{-5/2}
\leq \frac{-2\|Vw\|_1g_0^{3/2}}{\pi^2}
\eeq
\end{lem}
\begin{lem}\label{lemHAS1}
The expectation value of $H_{A1}$ is bounded above by,
\beq\label{proofthem1-5}
\overline\lim_{k_c\to\infty}\overline\lim_{\rho\to0}
\left(\frac{1}{|\Lambda|}\left\langle H_{A1}\right\rangle_{\Psi}\right)\rho^{-5/2}
\leq \frac{-8\|Vw\|_1g_0^{3/2}}{3\pi^2}
\eeq
\end{lem}
\begin{lem}\label{lemHAS2}
The expectation value of $H_{A2}$ is bounded above by,
\beq\label{proofthem1-6}
\overline\lim_{k_c\to\infty}\overline\lim_{\rho\to0}
\left(\frac{1}{|\Lambda|}\left\langle H_{A2}\right\rangle_{\Psi}\right)\rho^{-5/2}
\leq \frac{4\|Vw^2\|_1g_0^{3/2}}{3\pi^2}
\eeq
\end{lem}
%Combine the results in \eqref{proofthem1-1}-\eqref{proofthem1-6}.  

By definitions of $g_0$ and $w$ \eqref{av}, \eqref{gdef}, we have
\beq
 \|\nabla w\|_2^2-\|Vw\|_1+\|Vw^2\|_1=0, \;  
V_0-\|Vw\|_1=g_0
\eeq
Summing  \eqref{proofthem1-1}-\eqref{proofthem1-6}, we have 
\beq
\overline\lim_{k_c\to\infty}\overline\lim_{\rho\to0}
\left(\frac{1}{|\Lambda|}\left\langle H_{N}\right\rangle_{\Psi}-\rho^2_0g_0\right)\rho^{-5/2}
\leq \frac{26g_0^{5/2}}{15\pi^2}
\eeq
By definition of $\rho_0$ \eqref{defrho0}, we have proved (\ref{Maindesired}).
Since $N=\rho|\Lambda|$ and $g_0>0$, \eqref{Maindesired} is the estimate of Theorem \ref{lemma2.1}.

We now deduce Theorem \ref{lemma2} by choosing $k_c$ as a function of $\rho$.
Denote by $\Psi_{k_c,\rho}$ the trial state \eqref{deffal} with cutoff $k_c$ at density $\rho$ (and $L=\rho^{-25/24}$),
by $F(k_c,\rho)$ the quotient in Theorem \ref{lemma2.1} for $\Psi=\Psi_{k_c,\rho}$, and let $c_*=16/(15\pi^2)$.
By Theorem \ref{lemma2.1}, for every integer $n\geq 1$ there is $K_n\geq n$ with
$\overline\lim_{\rho\to0}F(K_n,\rho)<c_*+1/n$, hence there is $\rho_n>0$ such that $F(K_n,\rho)<c_*+1/n$
for all $\rho<\rho_n$; we may choose $\rho_1>\rho_2>\cdots$ with $\rho_n\to0$.
Let $k_c(\rho)=K_n$ for $\rho_{n+1}\leq\rho<\rho_n$, and $k_c(\rho)=K_1$ for $\rho\geq\rho_1$.
If $\rho<\rho_n$, then $\rho_{m+1}\leq\rho<\rho_m$ for some $m\geq n$, and
\[
F(k_c(\rho),\rho)=F(K_m,\rho)<c_*+\frac 1m\leq c_*+\frac 1n .
\]
Hence $\overline\lim_{\rho\to0}F(k_c(\rho),\rho)\leq c_*$, i.e., the periodic trial state
$\Psi=\Psi_{k_c(\rho),\rho}$ of $N=\rho|\Lambda|$ particles on $\Lambda=[0,L]^3$ satisfies the estimate of
Theorem \ref{lemma2}.
\end{proof}
%\par \textbf{Note:} The contribution of $\left\langle H_{A1}\right\rangle_{\Psi}$ comes from the interaction that annihilates one condensate particle and a particle with momentum in $P_L$ and creates a soft pair of particles.
%And the contribution of $\left\langle H_{A2}\right\rangle_{\Psi}$ comes from the interaction that annihilates a soft pair of particles and creates a soft pair of particles. 

\section{Estimates on the Numbers of Particles}

The first step to prove the Lemma \ref{lemkinetic} to Lemma \ref{lemHAS2} is to estimate the number of 
particles in the condensate, $P_{L}, P_{I}$, and $P_{H}$. This is the main task of this section and 
we start with  the following  notations.

\begin{mydef}\label{1}
Suppose $u_i, k_j \in P$ for $i = 1, \ldots, t$, $j = 1, \ldots, s$.
\begin{enumerate}
	\item  The expectation of the  product of  particle numbers with momenta  $k_1$, $\cdots$, $k_s$:
	\beqa\nonumber
&&Q_\Psi  \left(k_1,k_2,\cdots,k_s\right)
=\left\langle\prod_{j=1}^sa^\dagger_{k_j} a_{k_j}\right\rangle_{\Psi}
=\sum_{\al\in M} \left(\prod_{j=1}^s{\al(k_j)}\right)|f(\al)|^2
\eeqa

\item The probability to have $m_i$ particles with momentum 
	$u_i, i= 1 \ldots, s$: 
%$\Psi\left(\umt\right)$,  here $m_i\in \Z$ and $u_i\in P_L\cup P_I\cup P_H\cup P_0$
\beqa\label{2.1}
&& Q_\Psi\left(\umt\right)
\equiv\sum_{\al\in A} |f(\al)|^2\\\nonumber
 &&{\rm Here}\,\,\,A=\left\{\al\in M| \al(u_1)=m_1,\cdots,\al(u_t)=m_t\right\}
\eeqa
\item The expectation of the  product of  particle numbers with momenta  $k_1$, $\ldots$, $k_s$,
conditioned that there are  $m_i$ particles with momentum $u_i$:
\beqa\nonumber
&&Q_\Psi\left(k_1,\cdots,k_s \, | \, \umt\right)\\\nonumber
\equiv&&\left(\sum_{\al\in A} \left(\prod_{j=1}^s{\al(k_j)}\right)|f(\al)|^2\right)
\left(\sum_{\al\in A} |f(\al)|^2\right)^{-1}, \\\nonumber
%{\rm Here}\,\,\, A=&&\left\{\al\in M| \al(u_1)=m_1,\cdots,\al(u_t)=m_t\right\}
\eeqa
where $A$ is the same as in item 2; if $Q_\Psi\left(\umt\right)=0$, we set $Q_\Psi\left(k_1,\cdots,k_s \, | \, \umt\right)=0$.
\end{enumerate}
\end{mydef}





The following theorem provides the main estimates on the number of particles. 


\begin{thm}\label{thmA} 
In the limit $\lim_{k_c \to\infty}\lim_{\rho\to0}$, $  Q_\Psi(u)$  can be estimated as follows
\beqa\label{boundspsiuPMPHtotal}
&&\lim_{k_c \to\infty}\lim_{\rho\to0} \left(\rho^{-3/2}|\Lambda|^{-1}\sum_{u\in P_I\cup P_H} Q_\Psi(u)\right)=0\\
\label{boundspsiuPLtotal}
&&\lim_{k_c \to\infty}\lim_{\rho\to0}\left(\rho^{-3/2}|\Lambda|^{-1}\sum_{u\in P_L} Q_\Psi(u)\right)=\frac1{3\pi^2}g^{3/2}_0
\eeqa
\end{thm}




We first collect a few obvious identities of $f$ into the following lemma. 

\begin{lem} \label{f}
\begin{enumerate}
	\item  If $k\in P_I\cup P_H$ and $\al, \mathcal A^k \al  \in M$, then 
	\beq \label{properf1}f( \mathcal A^k \al )=\sqrt{\frac{\al(0)}{|\Lambda|}}\sqrt{\frac{\al(0)-1}{|\Lambda|}}\lambda_k f(\al)
	\eeq
	\item If $k\in P_L$, $\al\in M_k^s$ and $\al, \mathcal A^k \al  \in M$, then 
	\beq \label{properf2}f( \mathcal A^k \al )=\sqrt{\frac{\al(0)}{|\Lambda|}}\sqrt{\frac{\al(0)-1}{|\Lambda|}}\lambda_k f(\al)
	\eeq
  \item If $k\in P_L$, $\al\in M^a_k$ and $\al, \mathcal A^k \al  \in M$, then
 	\beq \label{properf3}f( \mathcal A^k\al )=\sqrt{\frac{\al(0)}{|\Lambda|}}\sqrt{\frac{\al(0)-1}{|\Lambda|}}\sqrt{\frac{\al^*(k)+1}{\al^*(k)}}\lambda_k f(\al)
	\eeq
	\item If $\al\in M_u^s$ and $ \mathcal A^{u, k}\al  \in M$, then 
	\beq\label{properf4}
f( \mathcal A^{u, k}\al )=2\sqrt{\frac{\al(0)}{|\Lambda|}}\sqrt{\frac{\al(u)}{|\Lambda|}}\sqrt{\lambda_{k+\frac{u}{2}}}\sqrt{\lambda_{-k+\frac{u}{2}}}f(\al)
	\eeq
	\item If $\al\in M_u^a$ and $ \mathcal A^{u, k}\al  \in M$, then (recall that $u\in P_L$, so that $\lambda_u<0$ by \eqref{lambda} and $g_0>0$)
	\beq\label{properf5}
f( \mathcal A^{u, k}\al )=\frac{1}{2\lambda_u}\sqrt{\frac{\al(0)}{|\Lambda|}}\sqrt{\frac{|\Lambda|}{\al(u)}}\sqrt{\lambda_{k+\frac{u}{2}}}\sqrt{\lambda_{-k+\frac{u}{2}}}f(\al)
	\eeq
\end{enumerate}
\end{lem}

In defining the space $M$,   the operation $ \mathcal A^{u, k}\al $  is not allowed when  $\al\in M_u^a$.  However, it is possible through rare coincidences  that  $ \mathcal A^{u, k}\al \in M$ even if $\al\in M_u^a$.  Clearly,   $\al\in M_u^a$ and $ \mathcal A^{u, k}\al  \in M $ imply that $\alpha(u)  = \alpha(-u) + 1$. 
%Furthermore,   $\al\in M_u^a$ and $ \mathcal A^{u, k}\al  \in M $  requires that  the particles created by  the soft pair creation $\pm k + u/2$ can also be obtained from various operations in the definition of $M$. 
The following lemma summarizes some properties we need for $\lambda$. Since $V$ is symmetric, i.e., $V(x)=V(-x)$, the functions $w$ and $g$ are even, so that $w_k=w_{-k}$ and $g_k=g_{-k}$ are real; since $g=V(1-w)\geq 0$, we also have $|g_k|\leq g_0$. Note that $g_k$ may be negative for large $|k|$.
\begin{lem}\label{lem-l}
For $\rho$ small enough (depending on $g_0$), the following holds.
\begin{enumerate}
	\item For any $k\in P_L\cup P_I\cup P_H$, $\lambda_{-k}=\lambda_k$ and
	\beq\label{ublambda}
	|\lambda_k|\leq g_0|k|^{-2}, \,\,\, |\rho\lambda_k|\leq 1-\const\eps_L ,
	 \eeq
	 where the constant depends only on $g_0$. Moreover, $|\lambda_k|=|w_k|=|g_k||k|^{-2}$ for $k\in P_I\cup P_H$.
	 \item For any $k\in P_L$, $\lambda_k$ is negative and
	 \beq\label{lublambdaPL0}
	 -\frac{g_0}{2}\eta_L^2\rho^{-1}\geq \lambda_k\geq-\rho^{-1}
	 \eeq
	 \item For any $k\in P_H$, $|\lambda_k|$ is bounded as
	 \beq\label{ublambdaPH0}
	|\lambda_k|\leq g_0\eps_H^{-2}
	 \eeq
%	\item For any $k$, $\lambda_k$ is very close to $w_k$, 
%	\beq\label{temp2.12}
%0\leq \left(\rho\lambda_k+ \rho w_k\right)\leq 2(\rho w_k)^2\leq \const \rho ^2|k|^{-4} 
%\eeq 
%\beq\label{temp2.13}
% 0\geq (\rho \lambda_k)^2- (\rho w_k)^2\geq -4(\rho w_k)^3 \geq \const\rho ^3|k|^{-6}
%\eeq
\end{enumerate}
\end{lem}
Indeed, for $k\in P_L$ we have $\rho\lambda_k=(1-s)/(1+s)$ with $s=(1+4x)^{1/2}$, $x=\rho g_0|k|^{-2}$. Hence $|\rho\lambda_k|=4x(1+s)^{-2}\leq x$ and $1-|\rho\lambda_k|=2/(1+s)\geq \eps_L/(1+\sqrt{g_0})$, since $x\le g_0\eps_L^{-2}$. Moreover, $|\rho\lambda_k|$ is increasing in $x$, $x\geq g_0\eta_L^2$ on $P_L$, and $4x(1+s)^{-2}\geq x/2$ if $x\leq 2-\sqrt2$; this gives \eqref{lublambdaPL0} once $g_0\eta_L^2=g_0\rho^{2\eta}\leq 2-\sqrt 2$. For $k\in P_I\cup P_H$ we have $\lambda_k=-w_k=-g_k|k|^{-2}$ by \eqref{gp}, hence $|\rho\lambda_k|\leq g_0\rho^{2\eta}$. The evenness of $\lambda$ follows from that of $w$ and from \eqref{lambda}.

%Further more, with  the fact that for any real radial symmetric functions $g_1$, $g_2$ and $h$, if $h$ is non-negative, then
%\beq
%\left|\int dx d y  g_1(x)h(x-y)g_1(y)-\int dx d y g_2(x)h(x-y)g_2(y)\right|  \leq  \left|\|g_1\|^2_2-\|g_2\|^2_2\right|\|h\|_1
%\eeq
%which implies
%\beq\label{lambda2w2}
%\left|V^{-2}\sum_{u,\,v\in P_H}V_{u-v}\left(\lambda_u\lambda_v- w_uw_v\right)\right|\leq \const\rho\eps^{-3}_H
%\eeq
%

%

%At first, we are going to prove the number of the zero-mementum-particles is equal to $\rho_0\Lambda$ in proper limit, i.e.,
%\beqa
%\lim_{m_c\to\infty}\lim_{\rho\to0}\left(Q_\Psi(0)/\Lambda-\rho_0\right)\rho_0^{-3/2}=0
%\eeqa 
%Here $\rho_0$ is defined in (\ref{defrho0}). Recall the definition of $\rhop$ in (\ref{defrho0}), we actually can improve this result as follows,

To prove Theorem \ref{thmA}, we start with the following estimate on the condensate.  


\begin{lem}\label{boundspsi0}
For any $\eps> 0$, when $\rho$ is small enough, 
the expected 
number of zero-momentum particles can be estimated by 
\beqa\label{boundspsi01}
&&|\Lambda|\rhom\leq Q_\Psi\left(0 \right)\leq |\Lambda|\rhop 
\eeqa
\end{lem}
%This Lemma holds also  if  we condition on more than three momenta, but we shall not need thoes results. 


\subsection{A Lower Bound on the Number of Condensates}

Since the total number of particles in fixed to be $N$, upper bound on $Q_\Psi\left(u \right)$ for ($u\neq 0$)
yields   a lower bound for $Q_\Psi\left(0 \right)$.  %The results about the lower bound on $Q_\Psi\left(0 \right)$ can help us to estimate the lower bounds on $Q_\Psi\left(u \right)$ for ($u\neq %0$), which implies the upper bounds on $Q_\Psi\left(0 \right)$. 
%
%We first prove upper bounds on the expectation of particle numbers for a fixed momentum $u \not = 0$.
The following lemma provides the upper bounds for expected number of particles in various momentum space 
regions.  


\begin{lem} \label{lem1} For small enough $\rho$,  
the following upper bounds on $ Q_\Psi(u)$ hold:
\begin{enumerate}
	 \item For $u\in P_I$, 
	 \beq\label{ubpsiuroughPM}
 Q_\Psi(u)\leq \frac{\lambda_u^2\rho^2}{1-\lambda_u^2\rho^2}=\sum_{i=1}^\infty(\lambda_u\rho)^{2i}
\eeq
\item For $u\in P_L$, 
\beq\label{ubpsiuroughPL}
 Q_\Psi(u)\leq  \frac{\lambda_u^2\rho^2}{1-\lambda_u^2\rho^2}\left(1+\const\frac{\rho m_c}{\eps_H}\right)
\eeq
\item For $u\in P_H$,  
	 \beq\label{ubpsiuroughPH}
 Q_\Psi(u)\leq\const \rho^2|u|^{-2}|\lambda_u|
\eeq
\end{enumerate}
\end{lem}

\begin{proof}  The basic idea to prove Lemma \ref{lem1} is the following lemma which  compares, in particular,
$Q_\Psi(\{u,m\})$ and $Q_\Psi(\{u,m-1\})$.

\begin{prop}\label{prop5}
When $\rho$ is small enough,  for any $u \in P_I$, we have
\beq\label{7}
 Q_\Psi(\{u,m\})\leq (\lambda_u\rho)^{2i} Q_\Psi(\{u,m-i\}){\rm \,\,\,for\,\,\,} m\geq i \ge 1
\eeq
\end{prop}


\begin{proof}
We start with the following simple observation, whose proof is given in Appendix \ref{app:props}.

\begin{prop} \label{prop1}  For any $u \in P_I$ fixed and  all  
$\al\in M$ with $\al(u)=m\ge 1$,  there exists  a $\be\in M$ such that  $\mathcal A^{u}\beta=\al$ and $\be(u)=m-1$. 
\end{prop}



{F}rom the property of $f$ in (\ref{properf1}) and $\be(0)\leq N$,  we obtain 
\[
|f(\mathcal A^{u}\beta)|= |\lambda_u|\frac{\sqrt{\be(0)(\be(0)-1)}}{|\Lambda|}|f(\be)|\leq |\lambda_u|\rho\,|f(\be)|
\]
Therefore,  we have for $m\ge 1$
\beqa\label{ubpsiuroughPM1}
Q_\Psi\left(\{u,m\}\right) &&\leq \sum_{\be(u)=m-1}|f(\mathcal A^{u}\beta)|^2\leq\lambda_u^2\rho^2\sum_{\be(u)=m-1}|f(\be)|^2
\\\nonumber
&&=\lambda_u^2\rho^2 Q_\Psi\left(\{u,m-1\}\right)
\eeqa
This proves \eqref{7} for $i = 1$. The general cases follow from iterations. 


\end{proof}


Together with  $\sum_{m=0}^N Q_\Psi\left(\{u,m\}\right)=1$, we have 
\beqa\label{4}
 Q_\Psi(u)=&&\sum_{m=1}^N m  Q_\Psi(\{u,m\})=\sum_{i=1}^N\left(\sum_{m=i}^N
Q_\Psi\left(\{u,m\}\right)\right)\\\nonumber
\leq &&\sum_{i=1}^N(\lambda_u\rho)^{2i}\left(\sum_{m=0}^N Q_\Psi\left(\{u,m\}\right)\right)\leq\frac{\lambda_u^2\rho^2}{1-\lambda_u^2\rho^2}
\eeqa
This proves  \eqref{ubpsiuroughPM}.




\par We now prove (\ref{ubpsiuroughPL}).  Recall that $\rho$ is small,  $1\leq m\leq \lceil m_c\rceil$ and $u\in P_L$.
Since $m_c$ is not an integer in general, rule 3 of Definition \ref{M} allows $\al^*(u)$ to reach $\lceil m_c\rceil$: we have $\al(u)\le\al^*(u)\leq \lceil m_c\rceil\le m_c+1\le 2m_c$ for all $\al\in M$ and $u\in P_L$.
{F}rom the definition of $M$ \eqref{defMsau},  all elements in the asymmetric part,  $M^a_u$, are generated from the symmetric part
$M^s_u$ via soft pair creations: if $\al\in M^a_u$ and $\al(-u)=\al(u)+1$, then $\al=\mathcal A^{u, k}\be$ for some $k$ and some $\be\in M^s_u$ with $\be(u)=\al^*(u)$; if $\al(u)=\al(-u)+1$, then $\al=\mathcal A^{-u, k}\be$ for some $k$ and some $\be\in M^s_u$ with $\be(-u)=\be(u)=\al^*(u)$. Thus
\beq\label{ubpsiuroughPL1.5}
\sum_{\al:\, \al\in M_u^a}^{\al^*(u)=m}|f(\al)|^2\leq\sum_{\be:\, \be\in M_u^s}^{\be(u)=m}\;\sum_{\sigma=\pm1}\left(\sum_{k: \pm k+\sigma u/2\in P_H}|f(\mathcal A^{\sigma u, k} \beta)|^2\right)
\eeq
%Here we have used that for $k, u$ fixed, there is at most one $\alpha$ for each $\beta$.
{F}rom \eqref{properf4} (applied to $\sigma u$; note $M^s_{-u}=M^s_u$), we have,  for $\beta(u)=\beta(-u) \le m$,
\beqa\nonumber
|f(\mathcal A^{\sigma u, k} \beta)|^2&=&  {4} \left|\lambda_{k+\sigma u/2}\lambda_{-k+\sigma u/2}\right|\frac{\be(0)}{|\Lambda|}\frac{\be(\sigma u)}{|\Lambda|}|f(\be)|^2\\\label{ubpsiuroughPL2}
&\leq&   {4} \left|\lambda_{k+\sigma u/2}\lambda_{-k+\sigma u/2}\right|\frac{\rho m}{|\Lambda|}|f(\be)|^2
\eeqa
Using the upper bound of $\lambda_k$ in (\ref{ublambda}) and  $|u|\ll |k|$, we have, for $\sigma=\pm1$,
\beq
\sum_{k:\pm k+\sigma u/2\in P_H}\left|\lambda_{k+\sigma u/2}\lambda_{-k+\sigma u/2}\right|\leq \sum_{p\in P_H}\const|p\,|^{-4}\leq \const\eps_H^{-1} |\Lambda|
\eeq
Inserting these results into \eqref{ubpsiuroughPL1.5}, we obtain 
%that the total probability of the asymmetric  configurations 
%are bounded by the total probability of the symmetric  configurations with a small factor proportional to $\rho$: 
\beq\label{ubpsiuroughPL1}
\sum_{\al:\, \al\in M_u^a, \al^*(u)=m}|f(\al)|^2\leq\const\frac{\rho m}{\eps_H} \sum_{\be:\, \be\in M_u^s,  \be(u)=m}|f(\be)|^2
%\leq \frac13\rho^{3/4} m \sum_{\be:\, \be\in M_u^s}^{ \be(u)=m}|f(\be)|^2
\eeq
Summing the last bound over $1 \le m \le \lceil m_c\rceil$, we have,  for each $u$ fixed,
\beq\label{boundMa}
\sum_{\al:\, \al\in M_u^a}|f(\al)|^2\leq\const\frac{\rho m_c}{\eps_H}
\eeq
Using this method, we can also prove, for $u\neq \pm v$, 
\beq\label{boundMaMa}
\sum_{\al:\, \al\in M_u^a, \al\in M_v^a}|f(\al)|^2\leq\const(\frac{\rho m_c}{\eps_H})^2
\eeq
{F}rom \eqref{ubpsiuroughPL1}, we have, for  $\rho$ is small enough
\beq
 Q_\Psi(u)\leq \sum_{m=1}^{\lceil m_c\rceil} \left(m\sum_{\al:\,\al\in M^s_u}^{\al(u)=m}|f(\al)|^2\right)\left(1+\const\frac{\rho m_c}{\eps_H}\right)
\eeq

Following the proof of (\ref{ubpsiuroughPM1}), we have the bound 
\beq\label{pubrho3}
\left[\sum_{\al:\al(u)=m}^{\al\in M^s_u}|f(\al)|^2\right]\leq \lambda_u^2\rho^2\left[\sum_{\be:\be(u)=m-1}^{\be\in M^s_u}|f(\be)|^2\right]
\eeq
Therefore, we can  prove (\ref{ubpsiuroughPL}) using the  argument of \eqref{4}.


\bigskip 


\par  We now prove  (\ref{ubpsiuroughPH}) by starting with the following proposition. Its proof is given in Appendix \ref{app:props}.

\begin{prop}\label{prop2}
For any $u \in P_H$ fixed and  all
$\al\in M$ with $\al(u)=m\ge 1$,  either there exists   $\be\in M$ such that  $\mathcal A^{u}\beta=\al$ and $\be(u)=m-1$
or there exists $v \in P_L$ and $\be\in M_v^s$  such that  $ \alpha = \mathcal A^{v,\,u-v/2}  \beta$ and $\be(u)=m-1$.
\end{prop}
In the second case, $\mathcal A^{v,\,u-v/2}$ creates one particle with momentum $u$ and one with momentum $v-u$, and $v-u\neq u$ since $|v|<2\eps_H<2|u|$; this is why $\be(u)=m-1$ in both cases.

{F}rom this proposition,  we have
\beq\label{4.1}
 Q_\Psi(\{u,m\})\leq \sum_{\be:\,\be(u)=m-1}\left[\left|f(\mathcal A^{u}\beta)\right|^2+\sum_{v \in P_{L}, \mathcal A^{v,\,u-v/2}  \beta\in M}^{\be\in M_v^s}\left|f(\mathcal A^{v,\,u-v/2}  \beta)\right|^2\right].
\eeq
By  the properties of $f$ in (\ref{properf1}, \ref{properf4}), we obtain 
%\beqa\nonumber
\[
|f( \mathcal A^u\be )|^2+\sum_{v\in P_{L} }|f(\mathcal A^{v,\,u-v/2}  \beta)|^2
\leq \left(\rho^2\lambda_u^2+\sum_{v\in P_{L}}4\rho\frac{\be(v)}{|\Lambda|}\left|\lambda_{u}\lambda_{-u+v}\right|\right)|f(\be)|^2 .
\]
Since  $v \in P_{L}$ and $ u \in P_{H}$, from (\ref{ublambda}) we have $\left|\lambda_u\right|$, 
$\left|\lambda_{-u+v}\right| \le \const |u|^{-2} $.  By definition of $M$, $\beta(v) \le \lceil m_{c}\rceil\le 2m_c$. Thus
\[
\sum_{v\in P_{L}}4\rho \frac{\be(v)}{|\Lambda|}
\le \sum_{v\in P_{L}}4\rho\frac{\lceil m_{c}\rceil}{|\Lambda|} \le \const\eta_{L}^{-3}m_c\rho^{5/2}.
\]
Hence  we have 
\[
|f( \mathcal A^u\be )|^2+\sum_{v\in P_{L}}|f(\mathcal A^{v,\,u-v/2}  \beta)|^2 \leq \const |u|^{-2}\left|\lambda_u\right|\rho^2|f(\be)|^2.
\]
Together with the bound in \eqref{4.1}, we obtain 
\beq\label{temp6.29}
 Q_\Psi(\{u,m\})\leq \const |\lambda_u| |u|^{-2}\rho^2 Q_\Psi(\{u,m-1\}){\rm\,\,\,for\,\,\,} m\geq 1 .
\eeq
Summing the last inequality over $m$, we have proved  (\ref{ubpsiuroughPH}). 

 
 
 
 \end{proof}
 
 
%\par Then combining (\ref{ubpsiuroughPM}), (\ref{ubpsiuroughPL}) and (\ref{ubpsiuroughPH}), one can easily prove that  $\forall \eps> 0$, there exists $m_c^*(\eps)$ and $\rho^*(m_c,\eps)$, when $m_c\geq m_c^*$, $\rho\leq\rho^*(m_c,\eps)$ 
%We claim  the following proposition on estimating the number of particles in $P_L, P_I, P_H$. 

The summations in the inequalities in Lemma \ref{lem1} can be performed; 
we summarize the conclusions in the following lemma. 

\begin{prop}\label{particlebound} %Suppose $\e_L, \eta_L$ are small numbers. 
Recall that $\e_L, \eta_L, \e_H$ are chosen in Definition \ref{def1}
as $\rho^\eta$. Then for any $k_c$ and  small enough $\rho$ we have  
\beqa\label{temp4.17}
&&|\Lambda|^{-1}\sum_{u\in P_I} Q_\Psi(u)  \leq \const\rho^{3/2+\eta}
\\
&&|\Lambda|^{-1}\sum_{u\in P_H} Q_\Psi(u)\leq \rho^{7/4} \label{4.18}
\\
&&|\Lambda|^{-1}\sum_{u\in P_L} Q_\Psi(u) \leq \left(\frac{g_0^{3/2}}{3\pi^2}+\const\rho^\eta\right)\rho^{3/2}
\label{4.19}
\eeqa
\end{prop} 



Assuming this proposition, we have, for any $\eps>0$, when $\rho$ is small enough, 
\beq\label{resultlbpsi0}
Q_\Psi(0)=N-\sum_{u\neq 0} Q_\Psi(u)
\geq \rhom |\Lambda|
\eeq
 This  proves the lower bound in Lemma \ref{boundspsi0}.  We now prove Proposition \ref{particlebound}.



\begin{proof} The upper bound \eqref{4.18} follows from \eqref{ubpsiuroughPH}, 
$|\lambda_u|\leq g_0|u|^{-2}$ \eqref{ublambda} and the assumption  $u \ge \e_H$ for $u \in P_H$.  

To prove the other bounds, we  first  sum over $u \in P_L$ in \eqref{ubpsiuroughPL} to have 
\beq
|\Lambda|^{-1}\sum_{u\in P_L} Q_\Psi(u)\leq    |\Lambda|^{-1} 
\sum_{u\in P_L}\frac{(\rho\lambda_u)^2}{1-(\rho\lambda_u)^2}(1+\rho^{3/4}), 
\eeq
where we have bounded the factor $\rho m_c /\e_H$ in the error term by $\rho^{3/4}$. 
%This can be achieved by choosing 
%$m_c\ge m_c^*$ and  $\rho\leq\rho^*(m_c,\eps)$ for some $m_c^*$ large and $\rho^*(m_c,\eps)$ small. 
%Recall $\rho_0=\rho-O(\rho^{3/2})$ and we have 



%Recall the definitions   $\lambda_u$ \eqref{lambda}. 
Let $h(k)=\sqrt{1+4g_0|k|^{-2}}$ and we can rewrite $\lambda$ as 
\beq \label{defhk}
\rho\lambda_{\sqrt {\rho} k} = 
 \frac {1 - h(k)} { 1 + h(k) }   .
 \eeq
We shall replace lattice sums by integrals several times, always in the following form.

\smallskip\noindent{\it Claim (Riemann sums).} Let $0\leq a<3<b$ and let $F$ be a $C^1$ function on $\R^3\setminus\{0\}$ such that
\[
|F(p)|+|p|\,|\nabla F(p)|\leq C|p|^{-a}(1+|p|)^{a-b},\qquad p\neq 0 .
\]
Then there is a constant $K$, depending only on $a$ and $b$, such that for all $0<\delta\leq1$, all $0\leq r<R\leq\infty$ and every set $S$ with $\{r<|p|<R\}\subset S\subset\{r\leq|p|\leq R\}$,
\begin{equation*}
\left|\left(\frac{\delta}{2\pi}\right)^3\sum_{p\in\delta\Z^3\cap S,\, p\neq0}F(p)-\int_{r\leq|p|\leq R}F(p)\frac{\rd^3p}{(2\pi)^3}\right|\leq KC\,e_a(\delta),
\tag{5.34$'$}\label{riemann}
\end{equation*}
where $e_a(\delta)=\delta$ for $a<2$, $e_a(\delta)=\delta(1+|\log\delta|)$ for $a=2$, and $e_a(\delta)=\delta^{3-a}$ for $2<a<3$.

\smallskip\noindent{\it Sketch of proof.} Compare each term $(\delta/2\pi)^3F(p)$ with $(2\pi)^{-3}$ times the integral of $F$ over the cube of side $\delta$ centered at $p$. For the cubes contained in $\{r<|q|<R\}$ and at distance at least $6\delta$ from the origin, the gradient bound gives a total error at most $KC\delta\int_{|q|\geq\delta}|q|^{-a-1}(1+|q|)^{a-b}\rd^3q\leq KC e_a(\delta)$. The cubes within distance $6\delta$ of the origin contribute at most $KC\delta^{3-a}$. The remaining cubes meet one of the spheres $|q|=s$, $s\in\{r,R\}$, $s\geq4\delta$; there are at most $Ks^2\delta^{-2}$ of them and $|F|\leq KCs^{-a}(1+s)^{a-b}$ on them, so they contribute at most $KC\delta s^{2-a}(1+s)^{a-b}\leq KCe_a(\delta)$.

\smallskip
Since $u\in\Lambda^*=(2\pi/L)\Z^3$, the rescaled momenta $k=\rho^{-1/2}u$ run over $\delta\Z^3$ with $\delta=2\pi/(L\rho^{1/2})=2\pi\rho^{13/24}$, and $|\Lambda|^{-1}\rho^{-3/2}=(\delta/2\pi)^3$. The function $F(k)=(h(k)-1)^2/(4h(k))$ behaves like $\sqrt{g_0}/(2|k|)$ near $0$ and like $g_0^2|k|^{-4}$ at infinity, so it satisfies the assumptions of \eqref{riemann} with $a=1$, $b=4$.
Thus we have
\beqa\label{3.1}
&&|\Lambda|^{-1}\rho^{-3/2}\left(\sum_{u\in P_L}\frac{(\rho\lambda_u)^2}{1-(\rho\lambda_u)^2}\right) \nonumber \\
=  &&\frac{1}{(2\pi)^3}
\int_{\eps_L\leq |k|\leq \eta_L^{-1}}
\frac{(h(k)-1)^2}{4h(k)}dk^3 +  O(\rho^{13/24})  .
\eeqa
The last error comes from replacing the summation by integral.




Due to the choices of  $\eps_L, \eta_L$, and since $(2\pi)^{-3}\int_{\R^3}F(k)\rd^3k=g_0^{3/2}/(3\pi^2)$ while the integrals of $F$ over $\{|k|<\eps_L\}$ and $\{|k|>\eta_L^{-1}\}$ are $O(\eps_L^2+\eta_L)$,
we can continue the computation as
\beqa\label{sum}
  &&\frac{1}{(2\pi)^3}
\int_{\eps_L\leq |k|\leq \eta_L^{-1}}
\frac{(h(k)-1)^2}{4h(k)}dk^3  \nonumber \\
=  &&\frac{1}{3\pi^2}g^{3/2}_0  +  O(\rho^{\eta})
\eeqa
Together with the previous bound, this proves \eqref{4.19},
since $\rho^{13/24}\leq\rho^{\eta}$.



Finally, for $u \in P_{I}$ we have $\lambda_u=-w_u$ and, by \eqref{ublambda}, $|\rho\lambda_u|\leq \rho g_0|u|^{-2}\leq g_0\eta_L^2\leq 1/2$ for $\rho$ small. Hence \eqref{ubpsiuroughPM} and $y^2/(1-y^2)\leq 2y^2$ for $|y|\leq 1/2$ give $Q_\Psi(u)\leq 2(\rho\lambda_u)^2\leq 2g_0^2\rho^2|u|^{-4}$. Comparing each term with the integral of $16|q|^{-4}$ over the cube of side $2\pi/L$ centered at $u$, we have, with a universal constant $K$, 
\[
|\Lambda|^{-1}\sum_{u\in\Lambda^*,\,|u|>R}|u|^{-4}\leq KR^{-1},\qquad R\geq 4\pi\sqrt3/L .
\]
Applying this with $R=\eta_L^{-1}\rho^{1/2}$, we obtain
\beq
|\Lambda|^{-1}\sum_{u\in P_I}Q_\Psi(u)\leq 2g_0^2\rho^2\,|\Lambda|^{-1}\sum_{u\in\Lambda^*,\,|u|>\eta_L^{-1}\rho^{1/2}}|u|^{-4}\leq 2Kg_0^2\,\rho^{3/2+\eta} .
\eeq
This proves \eqref{temp4.17} and concludes Proposition \ref{particlebound}.


\end{proof}

As a corollary to the proof,  we have the following estimates.
\begin{cor}\label{cortail}
\beq\label{tail}
\lim_{ n \to \infty}\lim_{\rho\to 0}|\Lambda|^{-1}\rho^{-3/2}\left(\sum_{u\in P_L}\sum_{m=1}^{n}\left(\rho\lambda_u\right)^{2m}\right)
= \frac{g_0^{3/2}}{3\pi^2}
\eeq
\end{cor}
\begin{proof}
Since $\sum_{m=1}^\infty(\rho\lambda_u)^{2m}=(\rho\lambda_u)^2/(1-(\rho\lambda_u)^2)$, by the previous proof, i.e., by \eqref{3.1}--\eqref{sum}, we only need to prove  the tail terms vanish.
Recall $\left|\rho \lambda_u \right|\le 1 - \const \e_L < 1$ in \eqref{ublambda}. Thus we have
% the factor $\left(\rho\lambda_u\right)^{ 2 \sqrt {m_c} } $ can be made arbitrarily small  for $m_c$ large
\beqa\label{temp6.35}
&&\lim_{n\to \infty}\overline\lim_{\rho\to 0}|\Lambda|^{-1}\rho^{-3/2}\left(\sum_{u\in P_L}\sum_{m=n+1}^{\infty}\left(\rho\lambda_u\right)^{2m}\right)
\\\nonumber
= &&\lim_{n\to \infty}\overline\lim_{\rho\to 0}|\Lambda|^{-1}\rho^{-3/2}    \left(\sum_{u\in P_L}  \frac{(\rho\lambda_u)^{2n+2}}{1-(\rho\lambda_u)^2}\right) \\\nonumber
\le &&\lim_{n\to \infty}\frac{1}{(2\pi)^3}
\int_{\R^3}
H(2n)dk^3  ,
\eeqa
where
$$
H(2n)=\frac{(h(k)-1)^2\left(\frac{1-h(k)}{1+h(k)}\right)^{2n}}{4h(k)}.
$$
Here we have used \eqref{riemann} for $H(2n)$ (with $a=1$, $b=4$, $\delta=2\pi\rho^{13/24}$, for each fixed $n$) and $H(2n)\geq0$.
Since $|1-h(k)|<1+h(k)$, $H(2n)$ decreases to zero pointwise, and $H(0)$ is integrable. By Lebesgue monotone convergence theorem, we have that  $\int H(2n)$ converges to zero. This proves the
Corollary.
\end{proof}

\bigskip 


We note that (\ref{temp6.29}) also shows that, for $u\in P_H$, $ Q_\Psi(\{u,m\})$ is exponentially small with $m$, i.e.,
\beq\label{temp5.19}
 Q_\Psi(\{u,m\})\leq (\const \left|\lambda_u\right|\rho^2|u|^{-2})^m.
\eeq
Furthermore, this result can be generalized to: for $u,v\in P_H$ with $u\neq \pm v$
\beq\label{boundpsiuvmn}
 Q_\Psi(\{u,m\},\{v,n\})\leq (\const |\lambda_u|\rho^2\eps_H^{-2})^m(\const |\lambda_v|\rho^2\eps_H^{-2})^n,
\eeq
which implies, for  $u,v\in P_H$ with $u\neq\pm v$, the following inequality:
\beq\label{boundpsifixeduv}
 Q_\Psi(u,v)\leq \const \left|\lambda_u\lambda_v\right|\rho^4\eps_H^{-4}.
\eeq
(For $u=v$, \eqref{boundpsifixeduv} fails: $Q_\Psi(u,u)\geq Q_\Psi(\{u,1\})$, which is of order $(\rho\lambda_u)^2$.)
To prove \eqref{boundpsiuvmn}, set $a_{m,n}=Q_\Psi(\{u,m\},\{v,n\})$ and $b_u=A\rho^2|\lambda_u|\eps_H^{-2}$, $b_v=A\rho^2|\lambda_v|\eps_H^{-2}$, with a constant $A$ to be chosen. Let $\al(u)=m\geq1$ and $\al(v)=n$. By Proposition \ref{prop2}, either $\al=\mathcal A^u\be$, or $\al=\mathcal A^{v',\,u-v'/2}\be$ with $v'\in P_L$, $\be\in M^s_{v'}$; in both cases $\be(u)=m-1$. Since $v\neq\pm u$, we have $\be(v)=n$, except in the second case with $v'-u=v$, i.e., $v'=u+v\in P_L$, where $\be(v)=n-1$. The terms with $\be(v)=n$ are estimated exactly as in the proof of \eqref{temp6.29} and give at most $b_u a_{m-1,n}$ for $A$ large. In the remaining case, by \eqref{properf4}, $\be(0)\leq N$ and $\be(u+v)\leq\lceil m_c\rceil$,
\begin{align*}
|f(\mathcal A^{u+v,\,(u-v)/2}\be)|^2&=4|\lambda_u\lambda_v|\frac{\be(0)\be(u+v)}{|\Lambda|^2}|f(\be)|^2\\
&\leq 4\rho\lceil m_c\rceil|\Lambda|^{-1}|\lambda_u\lambda_v||f(\be)|^2\leq b_ub_v|f(\be)|^2
\end{align*}
for $\rho$ small, since $|\Lambda|^{-1}=\rho^{25/8}$ and $1+25/8-\eta>4-4\eta$. Hence
\[
a_{m,n}\leq b_u\,a_{m-1,n}+b_ub_v\,a_{m-1,n-1},\qquad m\geq1,
\]
with $a_{m-1,-1}=0$. By \eqref{temp5.19}, $a_{0,n}\leq Q_\Psi(\{v,n\})\leq b_v^n$ for $A$ large, and induction on $m$ gives $a_{m,n}\leq(2b_u)^mb_v^n$, which is \eqref{boundpsiuvmn}. Since $2b_u,b_v\leq 1/2$ for $\rho$ small, summing $mn\,a_{m,n}$ over $m,n\geq1$ gives \eqref{boundpsifixeduv}.


\subsection{Proof of Lemma \ref{boundspsi0}: Upper Bound}


Proposition \ref{particlebound} states that  the density  of particles with momenta in $P_{I}$ and $P_{H}$ are much smaller than $\rho^{3/2}$. And it  implies  
an upper bound on the density  of particles with momenta in $P_{L}$.  We now prove a matching lower bound
\beq\label{upt01}
\sum_{u\in P_L} Q_\Psi(u)\geq \left(\frac{1}{3\pi^2}g^{3/2}_0-\eps\right)\rho^{3/2}|\Lambda|
\eeq
for $\rho$  small enough. Since the total number of particles is fixed, this will provide a upper bound on the number of particles in the condensate and hence proves the upper bound part of Lemma \ref{boundspsi0}.

 We start with the following lemma, which  bounds the average number of particles in the condensate under the condition that there are at most $k$ particles with momentum $u$. 
\begin{prop}\label{prop6} For $u \in P_{L}$ and for any $k$ fixed with $0\leq k\leq \lceil m_c\rceil$ ($m_{c}$ defined in \eqref{defmc}),
we have, for $\rho$ small enough,
\beqa\label{temp6.42}
\frac{\sum_{i=0}^kQ_\Psi(0 | \{u,i\}) Q_\Psi(\{u,i\})}{\sum_{i=0}^k Q_\Psi(\{u,i\})}\geq N-\const N\rho^{1/2}m_c .
\eeqa
\end{prop}


\begin{proof}  By \eqref{ubpsiuroughPL1},  the contribution of $\alpha \in M^a_u$ to $Q_\Psi(\{u,m\})$ for  $1\leq m\leq \lceil m_c\rceil$
is of lower order when compared with the contribution of $\alpha \in M^s_u$. The ratio of the contributions from $\alpha \in M^s_u$
between  $Q_\Psi(\{u,m\}) $ and $ Q_\Psi(\{u,m-1\})$ is estimated in \eqref{pubrho3}. Together with  the upper bound on  
$|\lambda_u|$ in (\ref{ublambda}) and the choices of  $\eps_L, \eps_H$, we have for $\rho$ small enough, 
\beq\label{mcmlbum0}
\frac{ Q_\Psi(\{u,m\})}{ Q_\Psi(\{u,m-1\})}\leq (\rho^2\lambda_u^2)(1+\const \frac{m_c\rho}{ \eps_H})\leq (1-\const (\eps_L- \frac{m_c\rho}{ \eps_H}))< 1.
\eeq
Hence $Q_\Psi(\{u,m\}) $ is monotonic decrease in $m$.  We thus  have for $0\leq k\leq \lceil m_c\rceil$,
\beq\label{mcmlbum1}
\sum_{i=0}^{k}Q_\Psi(\{u,i\})
\ge \frac{k+1}{\lceil m_c\rceil+1} \sum_{i=0}^{\lceil m_c\rceil}Q_\Psi(\{u,i\})  =  \frac{k+1}{\lceil m_c\rceil+1},
\eeq
where the last identity is the normalization of the state $\Psi$ (recall $\al(u)\leq\lceil m_c\rceil$ for $\al\in M$).

By definition of  $Q_\Psi(0 | \{u,i\})$ and \eqref{resultlbpsi0}, we have
$$
\sum_{i=0}^{\lceil m_c\rceil}Q_\Psi(0|\{u,i\}) Q_\Psi(\{u,i\})=
Q_\Psi(0)\geq N-\const N\rho^{1/2}
$$
On the other hand, for any $m$, $Q_\Psi(0|\{u,m\})\leq N$. Hence, the numerator 
on the left side of  \eqref{temp6.42} can be bounded by:
\beqa\nonumber
   &&  \sum_{i=0}^kQ_\Psi(0 | \{u,i\}) Q_\Psi(\{u,i\}) \\\nonumber
    && = \sum_{i=0}^{\lceil m_c\rceil}Q_\Psi(0| \{u,i\}) Q_\Psi(\{u,i\})-\sum_{i=k+1}^{\lceil m_c\rceil}Q_\Psi(0 | \{u,i\}) Q_\Psi(\{u,i\})\\\label{temp6.45}
 && \geq N-\const N\rho^{1/2}-N\sum_{i=k+1}^{\lceil m_c\rceil}   Q_\Psi(\{u,i\}) \\\nonumber
&&  =N{\sum_{i=0}^{k} Q_\Psi(\{u,i\})}-\const N\rho^{1/2}, \\\nonumber
\eeqa
where we have used  $\sum_{i=0}^{\lceil m_c\rceil} Q_\Psi(\{u,i\})=1$ in the last identity.
Finally, we  divide \eqref{temp6.45} by $\sum_{i=0}^{k} Q_\Psi(\{u,i\})$ and use  \eqref{mcmlbum1} to conclude 
\eqref{temp6.42}. 
\end{proof}


Return to the proof of \eqref{upt01} for $u\in P_L$. By rule 3 of Definition \ref{M}, if $\be\in M$, $\be(0)\geq2$ and $\be^*(u)<m_c$, then $\mathcal A^{u}\be\in M$; this is the case if $\be(u)\leq\lceil m_c\rceil-2$, since then $\be^*(u)\leq\be(u)+1\leq\lceil m_c\rceil-1<m_c$. Since $\mathcal A^{u}\beta$ is a one to one map 
(not necessarily surjective), we have 
\beq\label{lastlb2.1}
\sum_{i=1}^{\lceil m_c\rceil} Q_\Psi(\{u,i\}) \ge \sum_{\be(u)=0}^{\lceil m_c\rceil-2}
|f(\mathcal A^{u}\beta)|^2 
\eeq
{F}rom \eqref{properf2} and \eqref{properf3}, the right hand side is bounded below by 
\beq\label{lastlb2}
\lambda_u^2|\Lambda|^{-2}\sum_{\be(u)=0}^{\lceil m_c\rceil-2}
\left(\be(0)^2-\be(0)\right)|f(\be)|^2
\eeq
(the terms with $\be(0)\leq 1$ vanish). By Jensen's inequality and $\be(0)\leq N$, it is bounded below by 
\beq\label{lastlb3}
\lambda_u^2|\Lambda|^{-2}\left(\left(\frac{\sum_{\be(u)=0}^{\lceil m_c\rceil-2}
\be(0)|f(\be)|^2}{\sum_{\be(u)=0}^{\lceil m_c\rceil-2}
|f(\be)|^2}\right)^2-N\right)\sum_{\be(u)=0}^{\lceil m_c\rceil-2}
|f(\be)|^2
\eeq
By definition, 
\beq\label{temp6.49}
\frac{\sum_{\be(u)=0}^{\lceil m_c\rceil-2}
\be(0)|f(\be)|^2}{\sum_{\be(u)=0}^{\lceil m_c\rceil-2}
|f(\be)|^2}=\frac{\sum_{i=0}^{\lceil m_c\rceil-2}Q_\Psi(0 | \{u,i\}) Q_\Psi(\{u,i\})}{\sum_{i=0}^{\lceil m_c\rceil-2} Q_\Psi(\{u,i\})}
\eeq
The term on the right hand side can be estimated by  Proposition \ref{prop6} with $k=\lceil m_c\rceil-2$. Combining all estimates 
up to now and  we obtain 
\beq
\sum_{i=1}^{\lceil m_c\rceil} Q_\Psi(\{u,i\})\geq ((\rho-\rho^{5/4})\lambda_u)^{2}\sum_{i=0}^{\lceil m_c\rceil-2} Q_\Psi(\{u,i\})
\eeq
Finally, using \eqref{mcmlbum1}, we have
\beq
\sum_{i=1}^{\lceil m_c\rceil} Q_\Psi(\{u,i\})\geq  ((\rho-\rho^{5/4})\lambda_u)^{2}\left(1-\frac{2}{\lceil m_c\rceil+1}\right)
\eeq


We can generalize this result as follows. Let $1\le m\leq \lceil m_c\rceil-1$ and write $x_+=\max\{x,0\}$. If $\be(u)\leq\lceil m_c\rceil-m-1$ and $\be(0)\geq 2m$, then the states $(\mathcal A^{u})^j\be$, $1\leq j\leq m$, are in $M$, since $\left((\mathcal A^{u})^{j-1}\be\right)^*(u)=\be^*(u)+j-1\leq\lceil m_c\rceil-1<m_c$, and \eqref{properf2}, \eqref{properf3} give
\begin{align*}
|f((\mathcal A^{u})^m\be)|^2&\geq\lambda_u^{2m}|\Lambda|^{-2m}\prod_{j=0}^{m-1}\left(\be(0)-2j\right)\left(\be(0)-2j-1\right)|f(\be)|^2\\
&\geq \lambda_u^{2m}|\Lambda|^{-2m}\left(\be(0)-2m\right)^{2m}|f(\be)|^2 .
\end{align*}
Since $(\mathcal A^{u})^m$ is one to one and $m\leq ((\mathcal A^{u})^m\be)(u)\leq\lceil m_c\rceil$, iterating the argument in proving  \eqref{lastlb2.1}
and \eqref{lastlb2} we have 
\beqa\label{lastlb4}
&&\lambda_u^{2m}|\Lambda|^{-2m}\sum_{\be(u)=0}^{\lceil m_c\rceil-m-1}
\left(\be(0)-2m\right)_+^{2m}|f(\be)|^2 \nonumber\\
&&\qquad\le \sum_{i=m}^{\lceil m_c\rceil} Q_\Psi(\{u,i\})
\eeqa
Again, using Jensen's inequality for the convex function $x\mapsto (x-2m)_+^{2m}$, Proposition \ref{prop6} with $k=\lceil m_c\rceil-m-1$, and \eqref{mcmlbum1}, we have 
\beq
\sum_{i=m}^{\lceil m_c\rceil} Q_\Psi(\{u,i\})\geq  ((\rho-\rho^{5/4})\lambda_u)^{2m}\left(1-\frac{m+1}{\lceil m_c\rceil+1}\right)
\eeq
So with the fact $m_c=\rho^{-\eta}$, $Q_\Psi(u)$ can be bounded as follows, 
\beqa\nonumber
 Q_\Psi(u)= \sum_{m=1}^\infty\sum_{i=m}^\infty  Q_\Psi(\{u,i\})
\geq&& \sum_{ m=1}^{\lceil\sqrt{m_c}\rceil}((\rho-\rho^{5/4})\lambda_u)^{2m}\left(1-\frac{m+1}{\lceil m_c\rceil+1}\right)\\\label{ubpsiuPL}
\ge&&(1-2\rho^{\eta/2})  \sum_{i=1}^{\lceil\sqrt{m_c}\rceil}(\rho\lambda_u)^{2i}\\\nonumber
%=&&\frac{\rho^2\lambda_u^2}{1-\rho^2\lambda_u^2}(1-\eps-(\rho\lambda_u)^{2\sqrt{m_c}})
\eeqa
where we have used that, for $m\leq\lceil\sqrt{m_c}\rceil\leq\sqrt{m_c}+1$, $(m+1)/(\lceil m_c\rceil+1)\leq\rho^{\eta/2}+2\rho^{\eta}$ and $(1-\rho^{1/4})^{2m}\geq 1-4\rho^{1/4-\eta/2}$. 
%where $\e$ is independent of $u$ and $\e \to 0$ if $m_{c} \to \infty$.  
Since $\lceil\sqrt{m_c}\rceil\to\infty$ as $\rho\to0$ and all terms are nonnegative, for every fixed $n$ the right hand side of \eqref{ubpsiuPL} is at least $(1-2\rho^{\eta/2})\sum_{i=1}^{n}(\rho\lambda_u)^{2i}$ for $\rho$ small. 
Now the summation  over $u \in P_{L}$ was carried out in Corollary \ref{cortail} (first $\rho\to0$, then $n\to\infty$)
and  we have proved   (\ref{upt01}).  Since the total number of particle is $N$, 
the  bounds on $Q_\Psi(0)$ follows from (\ref{upt01}) and Proposition \ref{particlebound}. 
This concludes Lemma \ref{boundspsi0}. 

\bigskip 

%last term on the right hand side i the same as 
%The summation on the right hand side  can be estimated as 
% in \eqref{temp4.17}. 
 
The previous  method can be applied to yield the following estimates which will be useful later on. 
\begin{lem} For $u\in P_L$ and  $\rho$ sufficiently small, the following two bounds hold: 
\beq\label{ubpsiuPLmc}
	\sum_{m\geq m_c-2} m Q_\Psi(\{u,m\})\leq \frac{\rho^2\lambda_u^2}{1-\rho^2\lambda_u^2}\,\rho^{\eta/2}
\eeq
\beq\label{lbpsiuPL0}
\begin{split}
	&\sum_{\al(u)\leq m_c-2}|f(\al)|^2\al(0)(\al(0)-1)\al(u)\\
	&\qquad\geq N^2\frac{\rho^2\lambda_u^2}{1-\rho^2\lambda_u^2}(1-2\rho^{\eta/2}-|\rho\lambda_u|^{2\sqrt{m_c}}).
\end{split}
\eeq
The sum in \eqref{ubpsiuPLmc} has at most three nonzero terms, $m_c-2\leq m\leq\lceil m_c\rceil$. 
 \end{lem}
 \begin{proof}
Because $Q_\Psi(\{u,m\}) $ is monotonic decrease in $m$, we have, for $m\geq1$,
\[
\frac{m(m+1)}2\,Q_\Psi(\{u,m\})\leq\sum_{i=1}^m iQ_\Psi(\{u,i\})\leq Q_\Psi(u) .
\]
At most three integers $m$ satisfy $m_c-2\leq m\leq\lceil m_c\rceil$, and $Q_\Psi(\{u,m\})=0$ for $m>\lceil m_c\rceil$. Hence 
 \beq
 \sum_{m\geq m_c-2} m Q_\Psi(\{u,m\})\leq \frac{6} {m_c-1} Q_{\Psi}(u)\leq  \frac{\const} {m_c}Q_{\Psi}(u)
 \eeq
%we have, following the argument in \eqref{ubpsiuPL},
%\beq\nonumber
%\sum_{m=0}^{m_c-2} m Q_\Psi(\{u,m\}) \ge (1-\rho^{\eta/2})  \sum_{i=1}^{\sqrt{m_c}+1}(\rho\lambda_u)^{2i}
%\geq \frac{\rho^2\lambda_u^2}{1-\rho^2\lambda_u^2}(1-\rho^{\eta/2}-(\rho\lambda_u)^{2\sqrt{m_c}})
%\eeq
Together with the upper bound \eqref{ubpsiuroughPL} on $Q_\Psi(u)$, we have proved \eqref{ubpsiuPLmc}.


To prove \eqref{lbpsiuPL0}, we follow the argument in \eqref{lastlb4}. Let $1\leq m\leq m_{c}-2$. If $\be(u)\leq m_c-m-2$ and $\be(0)\geq 2m+1$, then $\al=(\mathcal A^{u})^m\be\in M$, $m\leq\al(u)\leq m_c-2$ and $\al(0)(\al(0)-1)\geq(\be(0)-2m-1)^2$. Hence 
\begin{align*}
&\lambda_u^{2m}|\Lambda|^{-2m}\sum_{\be(u)\leq m_c-m-2}
\left(\be(0)-2m-1\right)_+^{2m+2}|f(\be)|^2 \\
&\qquad\le \sum_{m\leq \al(u)\leq m_c-2} \al(0)(\al(0)-1)|f(\al)|^2
\end{align*}
Again, using Jensen's inequality, Proposition  \ref{prop6} and \eqref{mcmlbum1} (which gives $\sum_{\be(u)\leq m_c-m-2}|f(\be)|^2\geq 1-(m+3)/(\lceil m_c\rceil+1)$), and summing over $1\leq m\leq\lceil\sqrt{m_c}\rceil$, we have 
\beqa\nonumber
&&\sum_{\al(u)\leq m_c-2}|f(\al)|^2\al(0)(\al(0)-1)\al(u)\\\nonumber
&&\qquad=\sum_{m\geq1}\;\sum_{m\leq\al(u)\leq m_c-2} \al(0)(\al(0)-1)|f(\al)|^2\\
&&\qquad\geq (1-2\rho^{\eta/2})  \sum_{i=1}^{\lceil\sqrt{m_c}\rceil}(\rho\lambda_u)^{2i}N^2
\eeqa
Since $\sum_{i=1}^{n}x^i=x(1-x^n)/(1-x)$ and $|\rho\lambda_u|^{2\lceil\sqrt{m_c}\rceil}\leq|\rho\lambda_u|^{2\sqrt{m_c}}$, this implies \eqref{lbpsiuPL0}.
 \end{proof}




Lemma \ref{boundspsi0} can be extended to the following estimate: 

\begin{lem}\label{boundspsi02}
With the assumptions in Lemma \ref{boundspsi0},  $Q_\Psi(0,0)$ satisfies the estimate
\beqa
\big(|\Lambda|\rhom \big)^2\leq Q_\Psi\big(0, 0 \big)\leq \big(|\Lambda|\rhop\big)^2
\eeqa
\end{lem}
%After bounding $Q_\Psi(0)$ and $Q_\Psi(0,0)$'s, we are going to prove the upper bounds on $ Q_\Psi(u)$'s. 


\begin{proof}
By  Jensen's inequality and Lemma \ref{boundspsi0},   we have the lower bound
\[
Q_\Psi(0,0)\geq \left[Q_\Psi(0)\right]^2\geq \big(|\Lambda|\rhom \big)^2
\]

For the upper bound,  we start with  
\beqa\label{temp6.56}
Q_\Psi(0,0)
&= & N^2-2N\sum_{u\neq 0} Q_\Psi(u)+\sum_{u,v\neq 0} Q_\Psi(u,v)\\\nonumber
&\leq & (Q_\Psi(0))^2+\sum_{u,v\neq 0} Q_\Psi(u,v)
\eeqa
Since the number of particles with momentum $u\in P_L$ is at most  $\lceil m_c\rceil$,
\beqa\label{6.6}
\sum_{u\in P_L,v\neq 0} Q_\Psi(u,v) &&\leq\left(\sum_{u\in P_L}\lceil m_c\rceil\right)\left(\sum_{v\neq0}Q_\Psi(v)\right)
\eeqa
By definition of $P_L$, the number of its elements is asymptotically equal to $(6\pi^2)^{-1}(\eta_L^{-3}-\eps_L^3)\rho^{3/2}|\Lambda|$, hence at most $\const\eta_L^{-3}\rho^{3/2}|\Lambda|$, so that we have $\sum_{u\in P_L}\lceil m_c\rceil\leq \const m_c\eta_L^{-3}\rho^{3/2}|\Lambda|$.
The last factor in \eqref{6.6} can be estimated by Proposition \ref{particlebound}. Thus we have
\beq\label{temp6.58}
\sum_{u\in P_L,v\neq 0} Q_\Psi(u,v)=o(\rho^{5/2}|\Lambda|^2)
\eeq
For the terms $\sum_{u\in P_I\cup P_H,v\neq 0}$, 
the upper bound on the total number of particles in $P_I$ and $P_H$ in Proposition \ref{particlebound} yields that 
\beqa\label{temp6.59}
\sum_{u\in P_I\cup P_H,v\neq 0} Q_\Psi(u,v)&&\leq\sum_{u\in P_I\cup P_H} Q_\Psi(u) N =o(\rho^{5/2}|\Lambda|^2)
\eeqa
Inserting \eqref{temp6.58}, \eqref{temp6.59} into  \eqref{temp6.56}
and using the upper bound in Lemma \ref{boundspsi0},  we obtain the upper bound on $Q_\Psi(0,0)$. 
\end{proof}








\section{Estimates on  Kinetic Energy}

In this section, we will prove the kinetic energy estimate Lemma \ref{lemkinetic}.
Since $\Psi$ is supported on $M$, we have $Q_\Psi(u)=0$ unless $u\in P$; in particular there are no particles with momenta $0<|u|<\eps_L\rho^{1/2}$, and
\[
\left\langle\sum_{i=1}^N-\Delta_i\right\rangle_{\Psi}=\sum_{u\in P_L\cup P_I\cup P_H}u^2Q_\Psi(u) .
\]
This lemma follows from summing the estimates \eqref{boundspsiu2PMtotal}--\eqref{boundspsiu2PHtotal} of the next lemma, as follows. Since $\lambda_u^2=w_u^2$ for $u\in P_I\cup P_H$, the terms subtracted from $Q_\Psi(u)$ in \eqref{boundspsiu2PMtotal}--\eqref{boundspsiu2PHtotal}, summed with the weights $|\Lambda|^{-1}u^2$, add up to
\[
\rho_0^2|\Lambda|^{-1}\sum_{u\in P_L\cup P_I\cup P_H}u^2w_u^2+\frac{4g^{3/2}_0}{3\pi^2}\rho_0^{5/2}|\Lambda|^{-1}\sum_{u\in P_H}u^2w_u^2 .
\]
By \eqref{gp}, $u^2w_u^2=|g_u|^2|u|^{-2}\geq0$, so both sums are bounded by the sum over $u\in\Lambda^*\setminus\{0\}$. The function $|g_p|^2|p|^{-2}$ satisfies the assumptions of \eqref{riemann} with $a=2$ and any $b$ (since $g$ is smooth and, like $V$, decays fast, $g_p$ and $\nabla g_p$ decay faster than any power of $|p|$). Applying \eqref{riemann} with $\delta=2\pi/L=2\pi\rho^{25/24}$, and then Parseval's identity, we obtain
\begin{align*}
|\Lambda|^{-1}\sum_{u\in \Lambda^*\setminus\{0\}}\frac{|g_u|^2}{|u|^{2}}&=\int_{\R^3}\frac{|g_p|^2}{|p|^{2}}\frac{\rd^3p}{(2\pi)^3}+O(\rho^{25/24}\log\rho^{-1})\\
&=\|\nabla w\|_2^2+O(\rho^{25/24}\log\rho^{-1}),
\end{align*}
where we used $|p|^2|w_p|^2=|g_p|^2|p|^{-2}$. Since $\rho_0\leq\rho$, the subtracted terms are at most $\rho_0^2\|\nabla w\|_2^2+\frac{4g^{3/2}_0}{3\pi^2}\|\nabla w\|_2^2\rho^{5/2}+o(\rho^{5/2})$, and \eqref{proofthem1-1} follows. (The series $\sum_u u^2Q_\Psi(u)$ converges: on $P_H$ this follows from the bounds \eqref{temp7.16}--\eqref{temp7.20} below.)
%The following Lemma provides estimates on the kinetic energy of $\Psi$.


\begin{lem}\label{boundspsiu} 
In the limit $\rho\to0$, $Q_\Psi(u,\,v)$ can be bounded  above by 
\beq\label{boundspsiuvtotal}
\lim_{\rho\to 0} \left(\rho^{-5/2}|\Lambda|^{-2}\sum_{u,v\neq 0} Q_\Psi(u,v)\right) \le 0
\eeq
Furthermore, $\sum u^2  Q_\Psi(u)$ can be bounded  above as follows
\beqa\label{boundspsiu2PMtotal}
&&\overline\lim_{\rho\to 0} \left(\rho^{-5/2}|\Lambda|^{-1}\sum_{u\in P_I}u^2\left( Q_\Psi(u)-(\rho_0\lambda_u
)^2   \right)\right)\le 0\\
\label{boundspsiu2PLtotal}
&&\overline\lim_{\rho\to 0} \left(\rho^{-5/2}|\Lambda|^{-1}\sum_{u\in P_L}u^2\left( Q_\Psi(u)-(\rho_0w_u)^2\right)\right)\leq-\frac{8}{5\pi^2}g_0^{5/2}\\
\label{boundspsiu2PHtotal}
&&\overline\lim_{\rho\to 0} \left(\frac{\rho^{-5/2}}{|\Lambda|}\!\!\sum_{u\in P_H}u^2\!\!\left( Q_\Psi(u)\!-\!\left(\rho_0^2+\frac{4g^{3/2}_0}{3\pi^2}\rho_0^{5/2}\right)\!\lambda^2_u   \!\right)\!\!\right) \le 0
\eeqa
%Here $\widetilde\lambda_u$ is defined as follows,
%\[\widetilde\lambda_u=\frac{1-\sqrt{1+\frac{4\rho_0g_0}{u^2}}}{1+\sqrt{1+\frac{4\rho_0g_0}{u^2}}}\]
\end{lem}

%
%\begin{cor} \label{totalkineticPsi} 
%The total kinetic energy of $\Psi$ is bounded  above by 
% %\par For any $\eps> 0$, there exists $m_c^*(\eps)$ and $\rho^*(m_c,\eps)$, when $m_c\geq m_c^0$ and $\rho\geq\rho^*(m_c,\eps)$, the total kinetic energy of $\Psi$ is bounded from above as
%\beqa\label{resultcor1}
%&&\overline\lim_{m_c\to \infty}\overline\lim_{\rho\to 0}
%\left(\left\langle\sum_{i=1}^N-\Delta_i\right\rangle_{\Psi}-\Lambda\rho_0^2 \|\nabla w\|_2^2\right)\rho^{-5/2}|\Lambda|^{-1}\\\nonumber
%\leq&&  \left(\frac{4}{3\pi^2} \|\nabla w\|_2^2g_0^{-1}-\frac{8}{5\pi^2}\right)g_0^{5/2}
%\eeqa
%\end{cor}
%\begin{proof}
%Combine the results (\ref{boundspsiu2PMtotal}, \ref{boundspsiu2PLtotal}, \ref{boundspsiu2PHtotal})
%and use  $\lambda_u=-w_u$ for $u\in P_I\cup P_H$. 
%\end{proof}


\bigskip 



%We now prove Lemma \ref{boundspsiu}. 
%\subsubsection{Upper bounds in Lemma \ref{boundspsiu}}
\begin{proof} The bound  (\ref{boundspsiuvtotal}) was proved in \eqref{temp6.58} and \eqref{temp6.59}. 
We now prove  (\ref{boundspsiu2PMtotal}) concerning $ u\in P_I$.

The upper bound of $ Q_\Psi(u)$   in (\ref{ubpsiuroughPM}) can be rewritten as  
\beq
 Q_\Psi(u) \leq (\rho\lambda_u)^2+\frac{(\rho\lambda_u)^4}{1-(\rho\lambda_u)^2}
\eeq
Recall $\rho_{0} = \rho(1 + O(\sqrt \rho))$  and 
the bounds on $\lambda$ in (\ref{ublambda}).  
%We can bound the right side of the last inequality by
%\beqa
%&&  (\rho_0\lambda_u)^2(1+\const\rho^{1/2})+\const (\rho_0\lambda_u)^4\\\nonumber
%&\leq& (\rho_0\lambda_u)^2+\const\rho^{5/2}|u|^{-4}+\const\rho^3 |u|^{-6}
%\eeqa
%Here we used $\rho\lambda_u\ll1$ for $u\in P_I$.
Since $\rho^{1/2}\ll|u|\ll 1$ when $u\in P_I$, see Definition \ref{def1}, the error term of the last bound can be estimated by 
\beqa
&&\overline\lim_{\rho\to0} |\Lambda|^{-1} \rho^{-5/2} \sum_{u: \rho^{1/2}\ll|u|\ll 1}u^2\frac{(\rho\lambda_u)^4}{1-(\rho\lambda_u)^2}=0
\eeqa
This proves  (\ref{boundspsiu2PMtotal}).


\par We now prove  \eqref{boundspsiu2PLtotal} concerning $ u\in P_L$.
Following the strategy of the previous argument, we first  use $1\ge 1-(\rho\lambda_u)^2 \ge \const \e_{L}$
in \eqref{ublambda} and \eqref{lublambdaPL0} to
rewrite  the upper bound of $ Q_\Psi(u)$ in (\ref{ubpsiuroughPL}) as
\beq
 Q_\Psi(u)%\leq \frac{(\rho\lambda_u)^2}{1-(\rho\lambda_u)^2}(1+\const\frac{\rho m_c}{\eps_H})\\\nonumber
\leq \frac{(\rho\lambda_u)^2}{1-(\rho\lambda_u)^2}+\const\frac{\rho m_c}{\eps_H\e_L}
\eeq
The error terms are negligible in the sense that 
\[
\sum_{u\in P_L}u^2\frac{\rho m_c}{\eps_H\e_L}=o(\rho^{5/2}|\Lambda|)
\]
Since  $w_u=g_u |u|^{-2}$, $\rho_0-\rho=O(\rho^{3/2})$ and $|g_u-g_0|\leq \const |u|$,   we have
\beq
\lim_{\rho}\sum_{u\in P_L}u^2\left( (\frac{\rho g_0}{u^2} )^2-(\rho_0w_u)^2\right)\rho^{-5/2}|\Lambda|^{-1}=0
\eeq
Summarize what we have proved, we have the following  inequality:
\beqa\label{temp5.40}
&&\overline\lim_{\rho}\sum_{u\in P_L}u^2\left( Q_\Psi(u)-(\rho_0w_u)^2\right)\rho^{-5/2}|\Lambda|^{-1}\\\nonumber
\leq &&\overline\lim_{\rho}\sum_{u\in P_L}u^2\left(\frac{(\rho\lambda_u)^2}{1-(\rho\lambda_u)^2}-(\frac{\rho g_0}{u^2} )^2\right)\rho^{-5/2}|\Lambda|^{-1}
\eeqa
%Then, we are going to prove the right side of the above inequality is no more than $-\frac{8}{5\pi^2}g^{5/2}_0$. 
Let $u=\sqrt{\rho}k$ and $h(k)=\sqrt{1+4g_0|k|^{-2}}$ as in \eqref{defhk}. The integrand below behaves like $-g_0^2|k|^{-2}$ near $0$ and is $O(|k|^{-4})$ at infinity, so \eqref{riemann} applies with $a=2$, $b=4$ and $\delta=2\pi\rho^{13/24}$, with an error $O(\rho^{13/24}\log\rho^{-1})=O(\rho^{1/2})$. Then the right hand side  of \eqref{temp5.40} is estimated as
\beq\nonumber
 \frac{1}{(2\pi)^3}\int_{\eps_L\leq |k|\leq\eta_L^{-1}}k^2\left(\frac{1+2g_0|k|^{-2}}{2h(k)}-\frac{1+2(g_0|k|^{-2})^2}2\right)dk^3+O(\rho^{1/2})
\eeq
By the same behavior of the integrand, extending the integral to $\R^3$ changes it by $O(\eps_L+\eta_L)=O(\rho^\eta)$. Direct calculation yields that
\beq
\frac{1}{(2\pi)^3}\int_{k\in \R^3}k^2\left(\frac{1+2g_0|k|^{-2}}{2h(k)}-\frac{1+2(g_0|k|^{-2})^2}2\right)dk^3=-\frac{8}{5\pi^2}g^{5/2}_0,
\eeq
Inserting this result into (\ref{temp5.40}), 
we obtain  the desired result (\ref{boundspsiu2PLtotal}).

Finally, we  prove  \eqref{boundspsiu2PHtotal} concerning $ u\in P_H$.
Recall the bound  (\ref{temp6.29}) on the ratio of $ Q_\Psi(\{u,m\})/ Q_\Psi(\{u,m-1\})$. Since  
$|\lambda_u|\leq g_0 |u|^{-2}$ (\ref{ublambda}) and $u \in P_{H}$, the factor on the right hand side  of
(\ref{temp6.29}) can be bounded by $\rho^{3/2}$.
Thus we have 
%\beq\label{temp5.16}
%  Q_\Psi(\{u,m\})\leq \const \lambda_u |u|^{-2}\rho^2 m_c  Q_\Psi(\{u,m-1\}){\rm\,\,\,for\,\,\,} m\geq 1 
%\eeq
\beq\label{temp7.16}
 Q_\Psi(u)=\sum_{m}m  Q_\Psi(\{u,m\})\leq
\sum_{m\geq 1}  Q_\Psi(\{u,m\})(1+O(\rho^{3/2}))
\eeq
We now repeat the argument from \eqref{4.1} to \eqref{temp6.29} but refine the proof by using Proposition 
\ref{particlebound}. 
Hence for any $u\in P_H$, we have %can bound $\sum_{m\geq 1}  Q_\Psi(\{u,m\})$ by
\beqa\nonumber
&&\sum_{m\geq 1}  Q_\Psi(\{u,m\})\\\nonumber
&\leq&\sum_{\be} \left(\frac{\be(0)^2}{|\Lambda|^2}\lambda_u^2+\sum_{v\in P_L}4\frac{\be(0)}{|\Lambda|}\frac{\be(v)}{|\Lambda|}\left|\lambda_{u}\lambda_{-u+v}\right|\right)|f(\be)|^2\\\label{proofubpsiu12}
&\leq& |\Lambda|^{-2}\lambda_u^2Q_\Psi(0,0)+\sum_{v\in P_L}\rho |\Lambda|^{-1}\left(4Q_\Psi(v)\left|\lambda_{u}\lambda_{-u+v}\right|\right)
\eeqa
By mean value theorem and $\lambda_k=-g_k|k|^{-2}$ for $k\in P_I\cup P_H$, we have that  $\exists\tilde u\in \R^3: |\tilde u+u|\leq |v|$ s.t.
\beqa
|\lambda_{-u+v}-\lambda_{-u}|&\leq &\const \left(\left|\frac{\partial g_{\tilde u}}{\partial {\tilde u}}\right|{\tilde u}^{-2}+|g_{\tilde u}|{\tilde u}^{-3}\right)|v|
\eeqa
Here $\partial g/\partial\tilde u$ denotes the gradient $\nabla g_{\tilde u}$, and $\lambda_{-u}=\lambda_u$.
{F}rom  the estimates \eqref{ublambda} on $\lambda_u$ and $|u|\sim |\tilde u|$ (as $|v|\leq\eta_L^{-1}\rho^{1/2}\ll\eps_H<|u|$), we obtain, recalling that $g$ is even:
\beqa
|\lambda_u||\lambda_{-u+v}-\lambda_{-u}|&\leq &\const \left(\left|\frac{\partial g_{\tilde u}}{\partial {\tilde u}}g_u\right|u^{-4}+|g_{\tilde u}||g_u|u^{-5}\right)|v|\nonumber\\
&\leq &\const |u|^{-2}\eps_H^{-3}G(u)|v|,
\eeqa 
where by Schwarz inequality, we have:
\beq
G(u)=\max_{u':|u'-u|\leq\eta_L^{-1}\rho^{1/2}}\left\{\left|\frac{\partial g_{ u'}}{\partial {u'}}\right|^2+|g_{u'}|^2\right\}
\eeq
Since $g$ is smooth and decays fast, we have $|g_p|+|\nabla g_p|\leq \const(1+|p|)^{-2}$; as $\eta_L^{-1}\rho^{1/2}\leq1$, this gives $G(u)\leq\const(1+|u|)^{-4}$, and hence
\[
|\Lambda|^{-1}\sum_{u\in P_H}G(u)\leq\const|\Lambda|^{-1}\sum_{u\in\Lambda^*}(1+|u|)^{-4}\leq\const,
\]
uniformly in $\rho$ and in $L\geq1$. Together with the results on the total number of $P_L$ particles in (\ref{boundspsiuPLtotal}),
we obtain that, for $\rho$ small enough  and $u\in P_H$, the last term in (\ref{proofubpsiu12}) is bounded above by 
\beqa\label{temp7.20}
 \lambda_u^2\rho^{5/2}\left(\frac{4g^{3/2}_0}{3\pi^2}+\rho^\eta\right)+ \frac{\const\rho^{3}}{u^2\eps_H^{3}\eta_L}G(u)
\eeqa
The $Q_\psi(0,0)$ in the last second term of \eqref{proofubpsiu12} is bounded by Lemma \ref{boundspsi02}. Inserting (\ref{temp7.20}) and \eqref{proofubpsiu12} into (\ref{temp7.16}) and using  $\lambda_u^2=w_u^2$ for $u\in P_H$, we obtain that,
\beq
\overline\lim_{\rho\to0}\sum_{u\in P_H}u^2\left( Q_\Psi(u)\!-\!(\rho_0w_u)^2\left(1+\left[\frac{4g^{3/2}_0}{3\pi^2}\right]\rho_0^{1/2}\right)\!\right)\frac{\rho^{-5/2}}{|\Lambda|}\leq 0
\eeq
This proves  (\ref{boundspsiu2PHtotal}). 
\end{proof}
%\par We note: usting the method in (\ref{proofubpsiu1}) and fact $\lambda_u\leq g_0/u^2$, one can obtain a useful upper bound on $  Q_\Psi(\{u,m\})$ for $u\in P_H$ as follows,
%\beq\label{uppsium}
%  Q_\Psi(\{u,m\})\leq (\const\lambda_u \rho^2 /u^2 )^m
%\eeq










\section{Estimates on Pair Interaction Energies}




\subsection{Proof of Lemma \ref{lemHS1}}
First, with the fact $a_u^\dagger a_u^\dagger a_ua_u\leq (a_u^\dagger a_u)^2$ and $0\leq |V_u|\leq V_0$ for any $u$, we can bound $H_{S1}$ as follows
\beqa
H_{S1}\leq&& V_0|\Lambda|^{-1}\sum_{u,v} a_u^\dagger a_u a_v^\dagger a_v +|\Lambda|^{-1}\sum_{u\neq v} V_{u-v}a_u^\dagger a_u a_v^\dagger a_v\\\nonumber
\leq &&V_0N\rho+V_0|\Lambda|^{-1}\sum_{u\neq v}a_u^\dagger a_v^\dagger a_va_u=2V_0N\rho-V_0|\Lambda|^{-1}\sum_{u}(a_u^\dagger a_u)^2
\eeqa
Therefore we can bound the expectation value $\langle H_{S1}\rangle$:
\beq
\langle H_{S1}\rangle_\Psi\leq 2V_0N\rho-V_0|\Lambda|^{-1}\sum_{u}Q_\Psi(u,u)\leq 2V_0N\rho-V_0|\Lambda|^{-1}Q_\Psi(0,0)
\eeq
By the lower bounds of $Q_\Psi(0,0)$ in Lem. \ref{boundspsi02} and the definition of $\rho_0$ in \eqref{defrho0}, we have proved Lemma \ref{lemHS1}.


\subsection {Proof of Lemma \ref{lemHS2}}


%Let one and only one of $u,v$ be zero. In this subsection we focus on the four-particle-vertex that annihilate a pair of particles with the momenta $(u, -u)$ and create a pair of particles with the momenta $(v, -v)$.

We start the proof with  the following identity for $\langle\Psi|a^\dagger_{u_1}a^\dagger_{u_2}a_{u_3} a_{u_4}|\Psi\rangle$. 
\begin{lem}\label{id4a}
For any fixed $u_{1,2,3,4}\in \Lambda^*$ and $\al\in M$,  define  $T(\al)$ to be the state 
\beq
|T(\al)\rangle=Ca^\dagger_{u_1}a^\dagger_{u_2}a_{u_3} a_{u_4}|\al\rangle, 
\eeq 
where $C$ is the positive normalization constant when  $| T(\al)\rangle\not =0$. Then we have
\beq\label{resultid4a}
\langle\Psi|a^\dagger_{u_1}a^\dagger_{u_2}a_{u_3} a_{u_4}|\Psi\rangle=\sum_{\al\in M}f(\al)\overline {f(T(\al))}
\sqrt{\langle\al|a^\dagger_{u_4}a^\dagger_{u_3}a_{u_2}a_{u_1} a^\dagger_{u_1}a^\dagger_{u_2}a_{u_3} a_{u_4}|\al\rangle}
\eeq
\end{lem}
The map $T$ depends on $u_{1,2,3,4}$ and in principle it has to carry them as subscripts. We omit these subscripts since it will be clear
from the context  what they are. 

\begin{proof} %Recall $f(\be(\al))=0$ when $|\be(\al)\rangle=0$  or $\be(\al)\notin M$.
For any $u_{1,2,3,4}\in \Lambda^*$ fixed, by definition of $\Psi$, we have
\beq\label{proofid4a1}
\langle\Psi|a^\dagger_{u_1}a^\dagger_{u_2}a_{u_3} a_{u_4}|\Psi\rangle=\sum_{\al,\be\in M}f(\al) {\overline{f(\be)}}\langle\be|a^\dagger_{u_1}a^\dagger_{u_2}a_{u_3} a_{u_4}|\al\rangle
\eeq
%Here $\overline {f(\be)}$ is the complex conjugation of $f(\be)$. In our definition, for any $\be\in M$, $f(\be)\in \R$, so we can drop the bar.
Since the vectors $|\be\rangle$, $\be\in\widetilde M$, are orthonormal occupation-number states and $a^\dagger_{u_1}a^\dagger_{u_2}a_{u_3} a_{u_4}|\al\rangle$ is a multiple of a single such state, we have
\beq\label{proofid4a2}
\langle\be|a^\dagger_{u_1}a^\dagger_{u_2}a_{u_3} a_{u_4}|\al\rangle\neq 0\Rightarrow \be=T(\al)
\eeq
%On the other hand, with the definition  of $\be(\al)$, we have 
%\beqa\label{proofid4a3}
%&&\langle\be(\al)|a^\dagger_{u_1}a^\dagger_{u_2}a_{u_3} a_{u_4}|\al\rangle\\\nonumber
%=&&\sqrt{\langle\al|a^\dagger_{u_4}a^\dagger_{u_3}a_{u_2}a_{u_1} |a^\dagger_{u_1}a^\dagger_{u_2}a_{u_3} a_{u_4}|\al\rangle\langle\be(\al)|\be(\al)\rangle}\\\nonumber
%=&&\sqrt{\langle\al|a^\dagger_{u_4}a^\dagger_{u_3}a_{u_2}a_{u_1} |a^\dagger_{u_1}a^\dagger_{u_2}a_{u_3} a_{u_4}|\al\rangle}
%\eeqa
Since  $|T(\al)\rangle$ is normalized, the identity in Lemma \ref{id4a} is obvious. 
\end{proof}


Lemma \ref{lemHS2} follows from the following lemma and $\lambda_u=-w_u$ for $u\in P_H\cup P_I$.
Notice that the factor $2$ in the estimate of Lemma \ref{lemHS2} is due to the complex conjugate in the definition of 
$H_{S2}$. Similar factor also appears in Lemma \ref{lemHAS1}. 
% like   $$\langle\Psi|V_{u}|\Lambda|^{-1}a^\dagger_ua^\dagger_{-u}a_0 a_0|\Psi\rangle\,\,\,{\rm for  }\,\,\,u\neq 0.$$
\begin{lem}\label{symno20}
\par %The terms $\langle\Psi|V_{u}|\Lambda|^{-1}a^\dagger_ua^\dagger_{-u}a_0 a_0|\Psi\rangle$ are bounded as follows
\beqa\nonumber
&&\overline\lim_{k_c\to\infty}\overline\lim_{\rho\to0}\sum_{u\in P_I\cup P_H}\left(\langle V_u|\Lambda|^{-1}a^\dagger_ua^\dagger_{-u}a_0a_0\rangle- \rho_0^{2}V_u  \lambda_u\right)\rho^{-5/2}|\Lambda|^{-1}\\\label{u-uPMH}
&&\qquad=0
\\\nonumber
&&\overline\lim_{k_c\to\infty}\overline\lim_{\rho\to0}\sum_{u\in P_L}\left(\langle V_u|\Lambda|^{-1}a^\dagger_ua^\dagger_{-u}a_0a_0\rangle+\rho_0^2V_uw_u\right)\rho^{-5/2}|\Lambda|^{-1}\\\label{u-uPL}
&&\qquad\leq \frac{V_0g_0^{3/2}}{\pi^2}
\eeqa
\end{lem}

%\begin{cor} \label{potentialsymo20}
%Some symmetric parts of  interaction potential energy: $\langle\sum_{u\neq 0} V_u|\Lambda|^{-1}a^\dagger_ua^\dagger_{-u}a_0a_0\rangle_\Psi$ is bounded from below as follows
%\beq
%\overline\lim_{m_c,\,\rho}\left(\langle \sum_{u\neq 0}V_u|\Lambda|^{-1}a^\dagger_ua^\dagger_{-u}a_0a_0\rangle+\rho_0^2\Lambda \|Vw\|_1\right)\rho^{-5/2}|\Lambda|^{-1}=\frac{V_0g_0^{3/2}}{\pi^2}
%\eeq
%\end{cor} 
%\begin{proof}
%Combine \eqref{u-uPL} and \eqref{u-uPMH}. We note $\lambda_u=-w_u$ when $u\in P_H\cup P_I$
%\end{proof}
%\subsection{Proof for Lemma \ref{symno20}}
\begin{proof}
We first prove \eqref{u-uPMH} concerning with $u\in P_I\cup P_H$. By Lemma \ref{id4a}, we have
\beqa\label{prle81}
&&\langle V_u|\Lambda|^{-1}a^\dagger_ua^\dagger_{-u}a_0a_0\rangle\\\nonumber
=&&V_u|\Lambda|^{-1}\!\!\!\!\!\!\sum_{\al:\al\in M,  \mathcal A^u \al \in M }\!\!\!f(\al)\overline{f( \mathcal A^u \al )}\sqrt{(\al(0)^2-\al(0))(\al(u)+1)(\al(-u)+1)}
\eeqa
The case that $\al\in M$ and $ \mathcal A^u \al \notin M$ can only happen when $\al(0)=0$ or $1$
and thus has no contribution. 
{F}rom the relation between $f(\al)$ and $f( \mathcal A^u \al )$ in \eqref{properf1},  we have
\beq\label{temp8.15}
\eqref{prle81}=\lambda_uV_u|\Lambda|^{-2}\sum_{\al\in M }|f(\al)|^2\al(0)(\al(0)-1)\sqrt{(\al(u)+1)(\al(-u)+1)}
\eeq
By the Schwarz inequality, we have  
\beqa\label{temp8.16}
&&\left|\sum_{\al}\al(0)(\al(0)-1)\left(\sqrt{(\al(u)+1)(\al(-u)+1)}-1\right)|f(\al)|^2\right|\nonumber\\
\leq&&N^2\left|\sum_{\al}\frac{\al(u)+\al(-u)}2|f(\al)|^2\right|=N^2 Q_\Psi(u)
\eeqa
Inserting \eqref{temp8.16} into \eqref{temp8.15},  using $|\rho\lambda_u|\leq 1$ and $|V_u|\leq V_0$, and  summing  over $u\in P_I\cup P_H$,  we obtain
\beqa
&&\sum_{u\in P_I\cup P_H}\left|\langle V_u|\Lambda|^{-1}a^\dagger_ua^\dagger_{-u}a_0a_0\rangle-V_u\lambda_u\frac{Q_\Psi(0,0)-Q_\Psi(0)}{|\Lambda|^{2}}\right|\\\nonumber
\leq&&V_0 \rho\sum_{u\in P_I\cup P_H}Q_\Psi(u)
\eeqa
{F}rom  the upper bounds of $\sum Q_\Psi(u)$ in \eqref{temp4.17} and \eqref{4.18},  the right hand side  of above inequality is
bounded by $o(\rho^{5/2}|\Lambda|)$.  By  the bounds on $Q_\Psi(0,0)$ in  Lemma \ref{boundspsi02}, $Q_\Psi(0)\leq N$ and $|\Lambda|^{-1}\sum_{u\in P_I\cup P_H}|V_u\lambda_u|\leq\const$, we have proved  \eqref{u-uPMH}.



To prove \eqref{u-uPL} concerning  $u\in P_L$, we note that  \eqref{prle81} still holds, but
$ \mathcal A^u \al \notin M$ when $\al^*(u)=\lceil m_c\rceil$ (recall $\al^*(u)\leq\lceil m_c\rceil$, and $\al^*(u)<m_c$ is equivalent to $\al^*(u)\neq\lceil m_c\rceil$). Therefore, for $u\in P_L$, \eqref{prle81} is equal to
\beq\nonumber
V_u|\Lambda|^{-1}\sum_{\al:\al\in M, \al^*(u)<m_c}f(\al)\overline{f( \mathcal A^u \al )}\sqrt{\al(0)(\al(0)-1)(\al(u)+1)(\al(-u)+1)}
\eeq
We can express $f( \mathcal A^u \al )$ in terms of $f(\al )$; in both cases: $\al\in M_u^s$ or $\al\in M_u^a$, we have the following identity:
\beqa
&& f(\al)\overline{f( \mathcal A^u \al )}\sqrt{(\al(u)+1)(\al(-u)+1)} \\&=& \lambda_u |f(\al)|^2|\Lambda|^{-1} \sqrt{\al(0)(\al(0)-1)}(\al^*(u)+1)
%{\bf replaced \, by  \frac {\al^*(u)+1} {\al^*(u)}}
\eeqa
Hence, for $u\in P_L$, 
\beq\label{temp8.23}
\eqref{prle81}= \sum_{\al:\al\in M, \al^*(u)<m_c}\lambda_uV_u|\Lambda|^{-2}|f(\al)|^2\al(0)(\al(0)-1)(\al^*(u)+1)
\eeq
We note $\lambda_u<0$ and  $V_u\approx V_0>0$, for $u\in P_L$. For any $\alpha \in M$, $\al^*(u)-\al(u)\leq 1$ by definition.    Hence we can replace the summation $\al^*(u)<m_c$ 
by  $\al(u)\leq m_c-2$ to have an upper bound. Summing over $u\in P_L$ of \eqref{temp8.23}, we have
\beqa\nonumber
&&\langle \sum_{u\in P_L}V_u|\Lambda|^{-1}a^\dagger_ua^\dagger_{-u}a_0a_0\rangle\\\nonumber
 &\leq& \sum_{u\in P_L}
\sum_{\al(u)\leq m_c-2}\lambda_uV_u|\Lambda|^{-2}|f(\al)|^2\al(0)(\al(0)-1)\al(u)\\\label{temp8.24}
&& +\sum_{u\in P_L} \sum_{\al(u)\leq m_c-2}\lambda_uV_u|\Lambda|^{-2}|f(\al)|^2\al(0)(\al(0)-1)
\eeqa
The last term is equal to 
\beqa\label{temp8.25}
%&& \sum_{u\in P_L} \sum_{\al(u)\leq m_c-2}\lambda_uV_u|\Lambda|^{-2}|f(\al)|^2\al(0)(\al(0)-1) \\
& &\sum_{u\in P_L} \lambda_uV_u|\Lambda|^{-2}(Q_\Psi(0,0)-Q_\Psi(0))\\\nonumber
&-&\sum_{u\in P_L} \sum_{i>m_c-2}\lambda_uV_u|\Lambda|^{-2}\big(Q_\Psi(0,0|\{u,i\})-Q_\Psi(0|\{u,i\})\big)Q_\Psi(\{u,i\})
\eeqa
Here we have used $\al(u)\leq\al^*(u)\leq\lceil m_c\rceil$; since $m_c$ need not be an integer, the values $i$ removed by the restriction $\al(u)\leq m_c-2$ are all integers in $(m_c-2,\lceil m_c\rceil]$.
Since  $0\leq Q_\Psi(0,0|\{u,i\})-Q_\Psi(0|\{u,i\})\leq N^{2}$, $|V_u|\leq V_0$ and $|\rho\lambda_u|\leq 1$,  the last term in \eqref{temp8.25} is bounded from above by
\beq
\sum_{u\in P_L} V_0\rho^2|\lambda_u| \sum_{i> m_c-2}Q_\Psi(\{u,i\}) \le V_0\rho^{1+\frac\eta2}\sum_{u\in P_L}\frac{(\rho\lambda_u)^2}{1-(\rho\lambda_u)^2},
\eeq
where we have used \eqref{ubpsiuPLmc} (note that $i\geq1$ in the sum, since $m_c-2\geq1$). By \eqref{3.1}--\eqref{sum}, the right hand side is $o(\rho^{5/2}|\Lambda|)$.
For the first term of \eqref{temp8.25}, we can bound it by using  Lemma \ref{boundspsi02}.
We now use \eqref{lbpsiuPL0} to estimate the first term on the right hand side of \eqref{temp8.24}. 
Combining these results, we have 
\beqa\nonumber
&&\langle \sum_{u\in P_L}V_u|\Lambda|^{-1}a^\dagger_ua^\dagger_{-u}a_0a_0\rangle\\\nonumber
&\leq&\sum_{u\in P_L}\lambda_uV_u\rho^2\frac{(\rho\lambda_u)^2}{1-(\rho\lambda_u)^2}(1-2\rho^{\frac\eta2}-|\rho\lambda_u|^{2\sqrt{m_c}})\\\label{temp8.26}
&&+\sum_{u\in P_L}\lambda_uV_u\rho_0^2+o(\rho^{5/2}|\Lambda|)
\eeqa
Since $\left|\lambda_u\rho\right|\leq 1$ and $|V_u|\leq V_0$, we have
\beq\label{temp8.27}
\sum_{u\in P_L}|\lambda_uV_u|\rho^2\frac{(\rho\lambda_u)^2}{1-(\rho\lambda_u)^2}\leq \sum_{u\in P_L}V_0\rho\frac{(\rho\lambda_u)^2}{1-(\rho\lambda_u)^2}\leq \const\rho^{5/2}|\Lambda|
\eeq
By  \eqref{temp6.35}, we have
\beq\label{temp8.28}
\sum_{u\in P_L}|\lambda_uV_u|\rho^2\frac{(\rho\lambda_u)^2}{1-(\rho\lambda_u)^2}|\rho\lambda_u|^{2\sqrt{m_c}}\leq  o(\rho^{5/2}|\Lambda|)
\eeq
Inserting \eqref{temp8.27}-\eqref{temp8.28} into \eqref{temp8.26}, we have
\beqa\label{temp8.29}
&&\sum_{u\in P_L}\left(\langle V_u|\Lambda|^{-1}a^\dagger_ua^\dagger_{-u}a_0a_0\rangle+w_uV_u\rho_0^2\right)
\nonumber \\
&&  \le  \sum_{u\in P_L}(\lambda_u+w_u)V_u\rho^2_0 +\sum_{u\in P_L}V_u\rho^2\frac{\lambda_u^3\rho^2}{1-\rho^2\lambda_u^2}+o(\rho^{5/2}|\Lambda|)
\eeqa
Since  $|g_u-g_0|+ |V_{u}-V_{0}| \leq \const |u|$,  we can replace $w_u$ and $V_u$ by $g_0|u|^{-2}$ and $V_0$ in last inequality so that the rhs of \eqref{temp8.29} is bounded by 
\beqa
&&V_0\rho^2_0 \sum_{u\in P_L}(\lambda_u+g_0|u|^{-2}) +V_0\rho^2 \sum_{u\in P_L}\frac{\lambda_u^3\rho^2}{1-\rho^2\lambda_u^2}+o(\rho^{5/2}|\Lambda|)\\\nonumber
=&&V_0\rho^2_0 \sum_{u\in P_L}\left(\lambda_u+g_0|u|^{-2}+\frac{\lambda_u^3\rho^2}{1-\rho^2\lambda_u^2}\right)+o(\rho^{5/2}|\Lambda|)
\eeqa
Let $u=\sqrt \rho k$. We have
\beqa\label{intPLlambdag}
&&\lim_{\rho\to 0}\sum_{u\in P_L}\left(\lambda_u+g_0|u|^{-2}+\frac{\lambda_u^3\rho^2}{1-\rho^2\lambda_u^2}\right)\rho^{-1/2}|\Lambda|^{-1}\\\nonumber
=&&\lim_{\rho\to 0}\frac{1}{(2\pi)^3}\int_{\eps_L\leq |k|\leq \eta_L^{-1}}
 g_0|k|^{-2}\left(1-\frac{1}{\sqrt{1+4g_0|k|^{-2}}}\right)dk^3\\\nonumber
 =&&g_0^{3/2}\pi^{-2}
\eeqa
So the leading term of the right hand side of \eqref{temp8.29} is $V_0g_0^{3/2}\pi^{-2}\rho^{5/2}|\Lambda|$. This completes the proof of \eqref{u-uPL}.
\end{proof}


\subsection{Proof of Lemma \ref{lemHS3}}

%In this subsection we focus on estimating  $\langle\Psi|a_u^\dagger a_{-u}^\dagger a_v  a_{-v}|\Psi\rangle$. 
Define $P(u,v)$ by
\beq\label{defPuv}P(u,v)\equiv\sum_{\gamma\in M}\overline{f( \mathcal A^u\gamma)}f( \mathcal A^v\gamma)
\sqrt{(\gamma(u)+1)(\gamma(-u)+1)(\gamma(v)+1)(\gamma(-v)+1)}
\eeq
Recall  $f(\al)=0$ when $|\al\rangle=0$ or $\al\notin M$.
\begin{lem}\label{id4a2}
Let $u$, $v\in \Lambda^*$, $u\neq v$ and  $u,v\neq 0$. 
\par If one of $u$ and $v$  $\in P_L\cup P_I$, we have the following identity.
\beqa\label{resultid4a2}
&&\langle\Psi|a_u^\dagger a_{-u}^\dagger a_v  a_{-v}|\Psi\rangle=P(u,v)
\eeqa
\par If $u,v\in P_H$, we have
\beqa\label{resultid4a2o}|\langle\Psi|a_u^\dagger a_{-u}^\dagger a_v  a_{-v}|\Psi\rangle-P(u,v)|\leq \const\rho^4\left|\lambda_u\lambda_v\right|
\eeqa
\end{lem}
\begin{proof}
We first assume $u\neq -v$ (the case $u=-v$ is treated at the end of the proof). We  prove \eqref{resultid4a2} and assume without loss of generality  that $v\in P_L\cup P_I$. Using Lemma \ref{id4a} and $u\neq\pm v$, we rewrite $\langle a_u^\dagger a_{-u}^\dagger a_v  a_{-v}\rangle_\Psi$ as
\beq\label{id4a21}
\langle a_u^\dagger a_{-u}^\dagger a_v  a_{-v}\rangle_\Psi=\sum_{\al\in M}f(\al)\overline{f(T(\al))}\sqrt{(\al(u)+1)(\al(-u)+1)\al(v)\al(-v)}
\eeq
Here $|T(\al)\rangle=Ca_u^\dagger a_{-u}^\dagger a_v  a_{-v}|\al\rangle$ and $C$ is positive normalization constant. Since $v\in P_L\cup P_I$ and $\al(v)>0, \al(-v)>0$, by definition of $M$  
there exists unique $\gamma\in M$ such that 
\beq\label{gammaDaal}
 \mathcal A^v\gamma=\al
\eeq
Therefore, with $|T(\al)\rangle=Ca_u^\dagger a_{-u}^\dagger a_v  a_{-v}|\al\rangle$, we have
\beq\label{gammaDabe}T(\al)= \mathcal A^u\gamma.
\eeq
Furthermore, by \eqref{gammaDaal} and $u\neq\pm v$, we have
\beq\label{gammaualu}
\gamma(\pm u)=\al(\pm u) \rmand \gamma(\pm v)=\al(\pm v)-1.
\eeq
Inserting \eqref{gammaDaal}, \eqref{gammaDabe} and \eqref{gammaualu} into \eqref{id4a21}, we 
have proved  \eqref{resultid4a2}.

To  prove \eqref{resultid4a2o}, 
we define $N_v$ as the following set:
\beq
N_v\equiv\left\{\al\in M|\forall \gamma\in M,  \mathcal A^v\gamma\neq \al\right\}
\eeq
Following the previous argument, we have 
\beq
\left|\langle a_u^\dagger a_{-u}^\dagger a_v  a_{-v}\rangle_\Psi-P(u,v)\right|\leq\sum_{\al\in N_v,\be\in N_u} \left|f(\al)f(\be)\langle \be|a_u^\dagger a_{-u}^\dagger a_v  a_{-v}|\al\rangle\right|
\eeq
The right hand side can be divided into two cases: 
\beqa\label{rhs.decompose}
\sum_{\al\in N_v,\be\in N_u, \be(u)\be(-u)\geq \al(v)\al(-v)}\left|f(\al)f(\be)\langle \be|a_u^\dagger a_{-u}^\dagger a_v  a_{-v}|\al\rangle\right|\\\nonumber
+\sum_{\al\in N_v,\be\in N_u, \al(v)\al(-v)>\be(u)\be(-u) }\left|f(\al)f(\be)\langle \be|a_u^\dagger a_{-u}^\dagger a_v  a_{-v}|\al\rangle\right|
\eeqa
By definition of $f$, if $\langle \be|a_u^\dagger a_{-u}^\dagger a_v  a_{-v}|\al\rangle\neq 0$,  we have $|f(\be)|=|\lambda_u/\lambda_vf(\al)|$,  $\be(u)=\al(u)+1$ and $\be(-u)=\al(-u)+1$. Denote by  $N_{v,u}\subset N_v$ the set 
\beq
N_{v,u}\equiv\{\al\in N_v:(\al(u)+1)(\al(-u)+1)\leq \al(v)\al(-v)\}
\eeq
Hence   we can bound  \eqref{rhs.decompose} by
\beqa\label{temp8.41}
&&\sum_{\al\in N_v,\be\in N_u}\left|f(\al)f(\be)\langle \be|a_u^\dagger a_{-u}^\dagger a_v  a_{-v}|\al\rangle\right|\\\nonumber
\leq&&\sum_{\al \in N_{v,u}}\left | \frac{\lambda_u}{\lambda_v} \right | |f(\al)|^2\al(v)\al(-v)+\sum_{\be \in N_{u,v}}\left | \frac{\lambda_v}{\lambda_u} \right | |f(\be)|^2\be(u)\be(-u)
\eeqa



Now we bound $\sum_{\al \in N_{v,u}}|f(\al)|^2\al(v)\al(-v)$. If $\al   \in N_v$ and $\al(v)\al(-v)>0$, then with Proposition  \ref{prop2}, there exist $\al'$, $v'\in P_L$ with  $\al'\in M_{v'}^s$ such that 
\beq\label{alv'2}
\al=\mathcal A^{v', \,v-\frac{v'}{2}}\al'
\eeq
If $\al'\notin N_v$, then there exists $\gamma'$ s.t. $\mathcal A^{v}\gamma'=\al'$. Hence 
$$\mathcal A^v(\mathcal A^{v', \,v-\frac{v'}{2}}\gamma')=\al\Rightarrow\al\notin N_v
$$
and we have a contradiction. 
Hence  we have $\al'\in N_v$ and $\al'(-v)>0$. Again by  Proposition  \ref{prop2}, there exist $\al''$, $v''\in P_L$ such that  $\al''\in M_{v''}^s$ 
\beq\label{al'v''2}
\al'=\mathcal A^{v'', \,-v-\frac{v''}{2}}\al''
\eeq
Combining \eqref{alv'2} and \eqref{al'v''2} and using \eqref{properf4}, we  express $f(\al)$ in terms of  $f(\al'')$, $\al''(v')$, $\al''(v'')$ and $\al''(0)$ and $\lambda$'s. By definition of $M$,
 $\al''(\widetilde v)\leq \lceil m_c\rceil\leq 2m_c$  for any $\widetilde v\in P_L$ and  we obtain %By definition of $P_H$, if $k\in P_H$, $k\geq \eps_H$. With \eqref{temp2.12},
\beq\label{temp8.44}
|f(\al)|^2\leq \const \rho^2 m_c^2|\Lambda|^{-2}\lambda_v^2\left|\lambda_{-v+v'}\lambda_{v+v''}\right||f(\al'')|^2
\eeq
By \eqref{ublambdaPH0} and $-v+v',v+v''\in P_H$, we have
\beq
|f(\al)|^2\leq  \const \rho^2m_c^2|\Lambda|^{-2}\lambda_v^2 \eps_H^{-4}|f(\al'')|^2
\eeq
Summing  over $v'$, $v''\in P_L$ and $\al''\in M$, and using $\#P_L\leq\const\eta_L^{-3}\rho^{3/2}|\Lambda|$, we obtain
\beq
\sum_{\al \in N_v, \al(v)+\al(-v)\geq2}|f(\al)|^2\leq \const \rho^5\eta_L^{-6} m_c^2\lambda_v^2 \eps_H^{-4}\leq (\rho^2\lambda_v)^2
\eeq 
Similarly, one can prove that  
\beq
\sum_{\al\in N_v, \al(v)+\al(-v)\geq m}|f(\al)|^2\leq  (\rho^2\lambda_v)^m
\eeq
Hence, we can obtain
\beq\nonumber
\sum_{\al\in N_{v,u} }\left|\frac{\lambda_u}{\lambda_v} \right||f(\al)|^2\al(v)\al(-v)\leq\sum_{\al\in N_v }\left|\frac{\lambda_u}{\lambda_v} \right||f(\al)|^2\al(v)\al(-v)\leq 2\rho^4\left|\lambda_u\lambda_v \right|
\eeq
Inserting this result into \eqref{temp8.41} and using the symmetry, we obtain
\beq
\sum_{\al\in N_v,\be\in N_u}\left|f(\al)f(\be)\langle \be|a_u^\dagger a_{-u}^\dagger a_v  a_{-v}|\al\rangle\right|\leq \const \rho^4\left|\lambda_u\lambda_v\right|
\eeq
This proves \eqref{resultid4a2o} for $u\neq -v$.

It remains to consider $u=-v$. Then $a_u^\dagger a_{-u}^\dagger a_v  a_{-v}=a_v^\dagger a_v a_{-v}^\dagger a_{-v}$, and by \eqref{number}, $\langle a_u^\dagger a_{-u}^\dagger a_v  a_{-v}\rangle_\Psi=\sum_{\al\in M}\al(v)\al(-v)|f(\al)|^2$. Since $\mathcal A^{-v}=\mathcal A^{v}$, \eqref{defPuv} gives $P(-v,v)=\sum_{\gamma\in M}|f(\mathcal A^v\gamma)|^2(\gamma(v)+1)(\gamma(-v)+1)$, which is the same sum restricted to $\al\notin N_v$. If $v\in P_L\cup P_I$, every $\al\in M$ with $\al(v)\al(-v)>0$ is of the form \eqref{gammaDaal}, and \eqref{resultid4a2} follows. If $v\in P_H$, the difference is $\sum_{\al\in N_v}\al(v)\al(-v)|f(\al)|^2\leq 2\rho^4\lambda_v^2=2\rho^4|\lambda_u\lambda_v|$ by the bounds above and $\lambda_{-v}=\lambda_v$. This completes the proof.
\end{proof}
Using this lemma, we can estimate the term $\langle a^\dagger_ua^\dagger_{-u}a_va_{-v}\rangle$ as follows.

\begin{lem}\label{u-uv-vPsi}
For $u,v\in P_I\cup P_H$ with $u\neq\pm v$,
\beqa\label{u-uv-vPMH}
&&\left|\langle a_u^\dagger a_{-u}^\dagger a_va_{-v}\rangle-\lambda_u\lambda_v\frac{Q_\Psi(0,0)-Q_\Psi(0)}{|\Lambda|^2}\right|\\\nonumber
\leq &&\left|\lambda_u\lambda_v\right|\rho^2(\left(Q_\Psi(u,v)+Q_\Psi(u,-v)\right)/2+Q_\Psi(u)+Q_\Psi(v)+\const\rho^2)
\eeqa 
 For $u\in P_L$, $v\in P_I\cup P_H$,
\beqa\label{u-uv-vPLMH}
&&\left|\langle a_u^\dagger a_{-u}^\dagger a_va_{-v}\rangle-\lambda_u\lambda_v\left(\frac{Q_\Psi(0,0)-Q_\Psi(0)}{|\Lambda|^2}+\frac{\rho^4\lambda_u^2}{1-\rho^2\lambda_u^2}\right)\right|\\\nonumber
\leq &&\left|\lambda_u\lambda_v\right|\rho^2\big(\left(Q_\Psi(u,v)+Q_\Psi(u,-v)\right)/2+2Q_\Psi(v)\\\nonumber
+&&\frac{4\rho^2\lambda_u^2}{1-\rho^2\lambda_u^2}(\rho^{\eta/2}+|\rho\lambda_u|^{2\sqrt{m_c}})\big)
\eeqa
 For $u,v\in P_L$ with $u\neq\pm v$,
\beqa\label{u-uv-vPL}
&&\langle a_u^\dagger a_{-u}^\dagger a_va_{-v}\rangle-\lambda_u\lambda_v\frac{Q_\Psi(0,0)-Q_\Psi(0)}{|\Lambda|^2}\\\nonumber
\leq &&\left|\lambda_u\lambda_v\right|\rho^2\left(Q_\Psi(u,v)+2Q_\Psi(u)+2Q_\Psi(v)+3\right)
\eeqa
We note that there is no absolute value on the left hand side  of  the inequality  when $u,v\in P_L$.
\end{lem}


\begin{proof}
We first prove \eqref{u-uv-vPMH} concerning $u,v\in P_I\cup P_H$.  By  Lemma \ref{id4a2}, we have  
\beq\label{u-uv-vPMHproof1}
\left|\langle a_u^\dagger a_{-u}^\dagger a_va_{-v}\rangle-P(u,v)\right|\leq \const \rho^4\left|\lambda_u\lambda_v\right|,
\eeq
where $P(u,v)$ is defined in \eqref{defPuv}. By the property of $f$ in \eqref{properf1}, we can rewrite $P(u,v)$ as 
\beqa\label{u-uv-vPMHproof2}
\sum_{\gamma\in M, \mathcal A^{u}\gamma\in M, \mathcal A^{v}\gamma\in M}&&\lambda_u\lambda_v|f(\gamma)|^2\frac{\gamma(0)^2-\gamma(0)}{|\Lambda|^2}\\\nonumber
\times&&\sqrt{(\gamma(u)+1)(\gamma(-u)+1)(\gamma(v)+1)(\gamma(-v)+1)}
\eeqa
The situation that $\gamma\in M$ and $\mathcal A^{u(v)}\gamma\notin M$ can only happen when $\gamma(0)=1$ or $0$. But in this case,  $\gamma(0)^2-\gamma(0)=0$ and the term vanishes.  Hence the summation 
of $\gamma$ in \eqref{u-uv-vPMHproof2} can be replaced by $\sum_{\gamma\in M}$. Therefore, 
for  $u,v\in P_I\cup P_H$, we have 
%\beq
%P(u,v)=\sum_{\gamma\in M}\lambda_u\lambda_v|f(\gamma)|^2\frac{\gamma(0)^2-\gamma(0)}{|\Lambda|^2}\sqrt{(\gamma(u)+1)(\gamma(-u)+1)(\gamma(v)+1)(\gamma(-v)+1)}
%\eeq
%Furthermore, with the notations $Q_\Psi$'s, we have 
\beqa\label{u-uv-vPMHproof3}
&&\left|P(u,v)-\lambda_u\lambda_v\frac{Q_\Psi(0,0)-Q_\Psi(0)}{|\Lambda|^2}\right|\\\nonumber
\leq&&\sum_{\gamma\in M}\left|\lambda_u\lambda_v\right||f(\gamma)|^2\frac{\gamma(0)^2}{|\Lambda|^2}\\\nonumber
&&\times\left|\sqrt{(\gamma(u)+1)(\gamma(-u)+1)(\gamma(v)+1)(\gamma(-v)+1)}-1\right|
\eeqa
{F}rom  $\gamma(0)\leq N$ and the Schwarz inequality,  the rhs. is bounded by  
\beq\label{u-uv-vPMHproof3.5}
%\eqref{u-uv-vPMHproof3}\leq
\sum_{\gamma\in M}\left|\lambda_u\lambda_v\right||f(\gamma)|^2\rho^2
\left [ \left( \frac {\gamma(u)+\gamma(-u)}  2+1\right)\left(\frac {\gamma(v)+\gamma(-v)} 2 +1 \right) -1 \right ] 
\eeq
By symmetry, we have $Q_\Psi(u)=Q_\Psi(-u)$ and $Q_\Psi(u,v)=Q_\Psi(-u,-v)$. So we have
\beq
\eqref{u-uv-vPMHproof3.5}\leq\left|\lambda_u\lambda_v\right|\rho^2\left(\frac12\left(Q_\Psi(u,v)+Q_\Psi(u,-v)\right)+Q_\Psi(u)+Q_\Psi(v)\right)
\eeq
Together with \eqref{u-uv-vPMHproof1}, we have proved   \eqref{u-uv-vPMH}.

We now prove \eqref{u-uv-vPLMH} concerning $u\in P_L$, $v\in P_I\cup P_H$. Following arguments in the previous paragraph and using  \eqref{properf2} and \eqref{properf3}, we can rewrite $P(u,v)$ as
\beqa\label{u-uv-vPLMHproof1}
\sum_{\gamma\in M, \mathcal A^{u} \gamma \in M}&&\lambda_u\lambda_v|f(\gamma)|^2\frac{\gamma(0)^2-\gamma(0)}{|\Lambda|^2}\\\nonumber
&&\times\sqrt{(\gamma^*(u)+1)(\gamma^*(-u)+1)(\gamma(v)+1)(\gamma(-v)+1)}
\eeqa 
Notice that no matter we use \eqref{properf2} or \eqref{properf3}, the final result is the same. 
For $\gamma\in M$ with $\gamma(0)\geq 2$,  the case $\mathcal A^{u}\gamma\notin M$  can only happen when $\gamma^*(u)=\lceil m_c\rceil$. Hence, the summation of $\gamma$ in \eqref{u-uv-vPLMHproof1} can be replaced by $\sum_{\gamma^*(u)\neq \lceil m_c\rceil}$. Since  $\gamma^*(u)=\gamma^*(-u)$, for  $u\in P_L,v\in P_I\cup P_H$ we have
\beq
P(u,v)=\!\!\!\!\!\sum_{\gamma^*(u)\neq \lceil m_c\rceil}\!\!\!\!\!\lambda_u\lambda_v|f(\gamma)|^2\frac{\gamma(0)^2\!-\!\gamma(0)}{|\Lambda|^2}(\gamma^*(u)+1)\sqrt{(\gamma(v)+1)(\gamma(-v)+1)}
\eeq
%
%Then move the following term $$\sum_{\gamma}\lambda_u\lambda_v\frac{Q_\Psi(0,0)-Q_\Psi(0)}{|\Lambda|^2}-\sum_{\gamma}\lambda_u\lambda_v|f(\gamma)|^2\frac{\gamma(0)^2-\gamma(0)}{|\Lambda|^2}\gamma^*(u)$$ from rhs. to lhs. 
Since $\gamma(0)\leq N$, we have
\beqa\nonumber
&&\!\!\!\left|P(u,v)-\lambda_u\lambda_v\frac{Q_\Psi(0,0)\!-\!Q_\Psi(0)}{|\Lambda|^2}-\sum_{\gamma}\lambda_u\lambda_v|f(\gamma)|^2\frac{\gamma(0)^2-\gamma(0)}{|\Lambda|^2}\gamma^*(u)\right|
\\\nonumber
\leq&&\sum_{\gamma^*(u)\neq \lceil m_c\rceil}\left|\lambda_u\lambda_v\right||f(\gamma)|^2\rho^2(\gamma^*(u)+1)\left|\sqrt{(\gamma(v)+1)(\gamma(-v)+1)}-1\right|\\\label{u-uv-vPLMHproof2}
+&&\sum_{\gamma^*(u)=\lceil m_c\rceil}\left|\lambda_u\lambda_v\right||f(\gamma)|^2\rho^2\left(\gamma^*(u)+1\right)
\eeqa
We can replace $\sum_{\gamma^*(u)\neq \lceil m_c\rceil}$ in the first term of rhs.  by  $\sum_{\gamma\in M}$ to have an upper bound. Since  $\sqrt{(\gamma(v)+1)(\gamma(-v)+1)}-1\le [\gamma(v)+\gamma(-v)]/2$ and $\gamma^*(u)\leq \gamma(u)+1$, we can bound 
the right hand side of \eqref{u-uv-vPLMHproof2} by 
\beqa\label{u-uv-vPLMHproof3}
%\eqref{u-uv-vPLMHproof2}\leq&& 
\left|\lambda_u\lambda_v\right|\rho^2\left[\half\left(Q_\Psi(u,v)+Q_\Psi(u,-v)\right)+2Q_\Psi(v)+\!\!\!\!\sum_{\gamma(u)\geq m_c-1}\!\!\!\!2|f(\gamma)|^2\gamma(u)\right]
\eeqa
Here we used that $\gamma^*(u)=\lceil m_c\rceil$ implies $\gamma(u)\geq\lceil m_c\rceil-1$, i.e., $\gamma(u)\geq m_c-1$.
The last term   is bounded  in \eqref{ubpsiuPLmc}, i.e.,
\beq\label{u-uv-vPLMHproof3.5}
\sum_{\gamma(u)\geq m_c-1}|f(\gamma)|^2\gamma(u)\leq \frac{\rho^2\lambda_u^2}{1-\rho^2\lambda_u^2}\, \rho^{\eta/2}
\eeq



The estimate \eqref{u-uv-vPLMH} follows from last three inequalities and \eqref{u-uv-vPLMHproof1},   provided 
that we can establish the following estimate
\beq\label{8.2}
\sum_{\gamma} |f(\gamma)|^2\frac{\gamma(0)^2-\gamma(0)}{|\Lambda|^2}\gamma^*(u) 
=\frac{\rho^{4}\lambda_u^2}{(1-(\rho\lambda_u)^2)}[1 + O(\rho^{\eta/2})+O(|\rho\lambda_u|^{2\sqrt{m_c}})]
\eeq
To prove this, we first divide the summation of $\gamma$ into $\gamma\in M_u^s$ and $\gamma\in M_u^a$. 
For the case $\gamma\in M_u^s$, we have 
\beq
\sum_{\gamma \in M_u^s}|f(\gamma)|^2\frac{\gamma(0)^2-\gamma(0)}{|\Lambda|^2}\gamma^*(u)\leq\rho^2Q_\Psi(u)
\le  \rho^2\frac{(\rho\lambda_u)^2}{(1-(\rho\lambda_u)^2)}(1+\rho^{2/3}),
\eeq
where we have used \eqref{ubpsiuroughPL} in the last inequality. 
For the case $\gamma\in M_u^a$, using \eqref{ubpsiuroughPL1}, we have 
\beqa\nonumber
\sum_{\gamma \in M_u^a}|f(\gamma)|^2\frac{\gamma(0)^2-\gamma(0)}{|\Lambda|^2}\gamma^*(u)&&\leq
\const\rho^2 \frac{\rho m_{c}}{\eps_H}\sum_{\gamma\in M_u^s}|f(\gamma)|^2\gamma(u)\\
&&\le  \rho^{\frac{8}{3}}\frac{(\rho\lambda_u)^2}{(1-(\rho\lambda_u)^2)}.
\eeqa
This proves the upper bound part of \eqref{8.2}. The lower bound 
follows from  \eqref{lbpsiuPL0} since $\gamma^*(u) \ge \gamma(u)$. 




Finally,  we prove  \eqref{u-uv-vPL} concerning $u,v\in P_L$. Similar to the previous argument, by \eqref{properf2} and \eqref{properf3}, we can rewrite $P(u,v)$ as
\beqa\label{u-uv-vPLproof1}
\sum_{\gamma\in M, \mathcal A^{u}  \gamma \in M, \mathcal A^{v}  \gamma \in M}  && \lambda_u\lambda_v|f(\gamma)|^2\frac{\gamma(0)^2-\gamma(0)}{|\Lambda|^2} \\
&& \times  \sqrt{(\gamma^*(u)+1)(\gamma^*(-u)+1)(\gamma^*(v)+1)(\gamma^*(-v)+1)}\nonumber 
\eeqa 
Since  $\lambda_u\lambda_v\geq0$ and $\gamma^*(u)=\gamma^*(-u)$, we have for $u,v\in P_L$,
\beqa\nonumber
&&P(u,v)-\lambda_u\lambda_v\frac{Q_\Psi(0,0)-Q_\Psi(0)}{|\Lambda|^2}\\
&\leq& \sum_{\gamma\in M} \lambda_u\lambda_v\rho^2|f(\gamma)|^2\left|(\gamma^*(u)+1)(\gamma^*(v)+1)-1\right|
\eeqa
Using $\gamma^*-\gamma\leq 1$, we have proved  \eqref{u-uv-vPL}.
\end{proof}


We can now prove Lemma \ref{lemHS3}. %state the main estimate on the pair interaction energy. 
%
%\begin{lem}\label{totalu-uv-v}
%\beq\label{proofthem1-4}
%\overline\lim_{m_c,\,\rho}\left(\sum_{u,\,v\neq 0}
%\frac{V_{u-v}}{|\Lambda|^2}\langle a^\dagger_ua^\dagger_{-u}a_va_{-v}\rangle
%-\rho_0^2\|w^2V\|_1\right)\rho^{-5/2}=-\frac{2\|wV\|_1}{\pi^2}g_0^{3/2}
%\eeq
%\end{lem}

\begin{proof}
Denote by $\sum'_{u,v}$ the sum over $u,v\in P_L\cup P_I\cup P_H$ with $u\neq\pm v$. Since $\Psi$ only contains states with momenta in $P$, these are the only nonzero terms of $|\Lambda|^{-1}\langle H_{S3}\rangle_\Psi$.
Summing  \eqref{u-uv-vPMH}, \eqref{u-uv-vPLMH} and \eqref{u-uv-vPL} over these $u,v$ (for $u\in P_I\cup P_H$, $v\in P_L$ we use \eqref{u-uv-vPLMH} with $u$ and $v$ exchanged, since $\langle a^\dagger_ua^\dagger_{-u}a_va_{-v}\rangle=\overline{\langle a^\dagger_va^\dagger_{-v}a_ua_{-u}\rangle}$ and $V_{u-v}=V_{v-u}$), we obtain that
\beq
\sideset{}{'}\sum_{u,v}
\frac{V_{u-v}}{|\Lambda|^2}\langle a^\dagger_ua^\dagger_{-u}a_va_{-v}\rangle
\le A+ B + \Omega
\eeq
where
\[
A= \frac{Q_\Psi(0,0)-Q_\Psi(0)}{|\Lambda|^2}  \sideset{}{'}\sum_{u,v}\frac{V_{u-v}}{|\Lambda|^2}\lambda_u\lambda_v
\]
\[
B= 2\sum_{u\in P_L,v\in P_I\cup P_H}\frac{V_{u-v}}{|\Lambda|^2}\lambda_u\lambda_v\frac{\rho^4\lambda_u^2}{1-\rho^2\lambda_u^2}
\]
\beqa\label{longexpression1}
\Omega &= &\frac{1}{|\Lambda|^2}\bigg(\sideset{}{'}\sum_{u,v}V_0|\lambda_u\lambda_v|\rho^2Q_\Psi(u,v)+\!\!\!\sum_{u\in P_I\cup P_H,v\neq 0}\!\!\!\!\!\!4|\lambda_u\lambda_vV_{u-v}|\rho^2Q_\Psi(u)\nonumber \\
&+&\sum_{u,\,v\in P_L}|V_{u-v}||\lambda_u\lambda_v|\rho^2(4Q_\Psi(u)+3)+\sum_{u,v\in P_I\cup P_H}\const\rho^4|\lambda_u\lambda_v||V_{u-v}|\nonumber \\
&+&\sum_{u\in P_L,v\in P_I\cup P_H}|\lambda_u\lambda_v||V_{u-v}|\frac{8\rho^4\lambda_u^2}{1-\rho^2\lambda_u^2}(\rho^{\eta/2}+|\lambda_u\rho|^{2\sqrt{m_c}})\bigg)
\eeqa
%Since $V_{u-v} \sim V_0 \ge 0$ when $u,v \in P_L$, we can remove the absolute value of $V$ in the definition of $B$.

The error term $\Omega$ can be bounded by using  the following facts, (1): $|\rho\lambda_u|\leq 1$, (2): $\sum_{v\neq 0}|\lambda_vV_{u-v}|\leq \const|\Lambda|$, (3): $|V_u|\leq V_0$, (4): $|\lambda_u|\leq g_0|u|^{-2}$for any $u\neq 0$ and (5): $\sum_{u,v}|\lambda_uV_{u-v}\lambda_v|\leq \const |\Lambda|^2$:
\beqa\nonumber
\Omega  &\le & \frac{\const}{|\Lambda|^2}\bigg(\sum_{u,\,v\neq 0}Q_\Psi(u,v)+\sum_{u\in P_I\cup P_H} Q_\Psi(u)\rho|\Lambda|+\sum_{u,\,v\in P_L}\frac{Q_\Psi(u)+1}{u^2v^2}\rho^2\\\label{temp8.70}
&+&\rho^4|\Lambda|^2+\sum_{u\in P_L}\frac{\rho^3\lambda_u^2}{1-\rho^2\lambda_u^2} |\Lambda| (\rho^{\eta/2}+|\lambda_u\rho|^{2\sqrt{m_c}})\bigg)
\eeqa
By  \eqref{boundspsiuvtotal} and \eqref{boundspsiuPMPHtotal}, the first two terms on the right hand side
are bounded by $o(\rho^{5/2})$. Using the trivial bound  $Q_\Psi(u)\leq \lceil m_c\rceil$ for $u\in P_L$,
the third term is also bounded by $o(\rho^{5/2})$. By  \eqref{temp8.27} and \eqref{temp8.28}, the last term is also $o(\rho^{5/2})$.  Hence the error terms are bounded by $\Omega \leq o(\rho^{5/2})$. 


We now estimate $A$ and $B$.  Notice that
$
(Q_\Psi(0,0)-Q_\Psi(0))|\Lambda|^{-2}
= \rho_{0}^{2} + o(\rho^{5/2})
$. 
Hence we shall replace this factor in $A$ by $\rho_{0}^{2}$.  
Since  $\lambda_{u}= - w_{u}$ for $u \in P_{I}\cup P_{H}$ and $V_{u-v}=V_{v-u}$, we have
\begin{align*}
\sideset{}{'}\sum_{u,v}V_{u-v}\lambda_u\lambda_v={}&\sideset{}{'}\sum_{u,v}V_{u-v}w_uw_v-2\sideset{}{'}\sum_{u\in P_L,v}V_{u-v}(\lambda_u+w_u)w_v\\
&+\sideset{}{'}\sum_{u, v\in P_L}V_{u-v}(\lambda_u+w_u)(\lambda_v+w_v)
\end{align*}
We can now decompose $A$ into
\beq \label{temp8.73}
A = \|w^2V\|_1 \rho_{0}^{2}+A_{1}+A_{2}+A_{3} + o(\rho^{5/2})
\eeq
where
\[
A_{1}=
-2\rho_{0}^{2} \sum_{u\in P_L,v\in P_I\cup P_H}\frac{V_{u-v}}{|\Lambda|^2}(\lambda_u+w_u)w_v
\]
\[
A_{2}=
-2\rho_{0}^{2} \sum_{u,v\in P_{L},\,u\neq\pm v}\frac{V_{u-v}}{|\Lambda|^2}(\lambda_u+w_u)w_v
\]
\[
A_{3}= \sum_{u,v\in P_L,\,u\neq\pm v}\frac{V_{u-v}}{|\Lambda|^2}(\lambda_u+w_u)(\lambda_v+w_v)\rho_{0}^{2}
\]
In \eqref{temp8.73} we used that the terms missing from $|\Lambda|^{-2}\sum'_{u,v}V_{u-v}w_uw_v$, compared with the sum over all $u,v\neq 0$ (those with $0<|u|<\eps_L\rho^{1/2}$ or $0<|v|<\eps_L\rho^{1/2}$, and those with $u=\pm v$), are bounded by
\[
\const|\Lambda|^{-1}\sum_{0<|u|<\eps_L\rho^{1/2}}|u|^{-2}+\const|\Lambda|^{-2}\sum_{|u|\geq\eps_L\rho^{1/2}}|u|^{-4}\leq \const\eps_L\rho^{1/2};
\]
after multiplication by $\rho_0^2$ they are $O(\rho^{5/2+\eta})$.

To bound $A_2$ and $A_3$, we first show that
\beq\nonumber
|\Lambda|^{-1}\sum_{u\in P_L}|\lambda_u+w_u|\leq\const\rho^{1/2}.
\eeq
For $u\in P_L$, the definition of $\lambda$ gives
\[
\frac{\lambda_u}{1-\rho^2\lambda_u^2}+g_0|u|^{-2}=g_0|u|^{-2}\Big(1-\frac{1}{\sqrt{1+4\rho g_0|u|^{-2}}}\Big)\geq 0,
\]
and $|w_u-g_0|u|^{-2}|\leq\const|u|^{-1}$ since $|g_u-g_0|\leq\const|u|$. We write
\[
\lambda_u+w_u=\Big(\frac{\lambda_u}{1-\rho^2\lambda_u^2}+g_0|u|^{-2}\Big)-\frac{\rho^2\lambda_u^3}{1-\rho^2\lambda_u^2}+\big(w_u-g_0|u|^{-2}\big).
\]
After summation over $u\in P_L$ and division by $|\Lambda|$, the first term contributes at most $\const\rho^{1/2}$ by \eqref{intPLlambdag}; since $|\rho^2\lambda_u^3|\leq\rho^{-1}(\rho\lambda_u)^2$, the second contributes at most $\const\rho^{1/2}$ by \eqref{3.1}--\eqref{sum}; and the third at most $\const|\Lambda|^{-1}\sum_{u\in P_L}|u|^{-1}\leq\const\eta_L^{-2}\rho$. Since moreover $|V_{u-v}|\leq V_0$ and $|\Lambda|^{-1}\sum_{v\in P_L}|w_v|\leq\const\eta_L^{-1}\rho^{1/2}$ (from $|w_v|\leq\const|v|^{-2}$), we obtain $|A_{3}|\leq\const\rho^{3}$ and $|A_{2}|\leq\const\eta_L^{-1}\rho^{3}$. Hence $A_{2}$ and $A_{3}$ are $o(\rho^{5/2})$.

If we replace $\rho^{2}$ in $B$ by  $\rho_{0}^{2}$, which is equal to $\rho^2-O(\rho^{5/2})$, we have 
\beq
B+ A_{1}=
-2\sum_{u\in P_L,v\in P_I\cup P_H}\frac{V_{u-v}}{|\Lambda|^2}w_v\rho^2\left(\lambda_u\frac{\rho^2\lambda_u^2}{1-\rho^2\lambda_{u}^2}+\lambda_u+w_u\right) 
\eeq
Using  $|V_{u-v}-V_v|\leq \const |u|$ for $u \in P_{L}$ and $v \in P_{I}\cup P_{H}$,  we can simplify $B+A_1$ as
\beq\label{temp8.75}
B+A_1\leq -\frac{2\|Vw\|_1}{|\Lambda|}\rho^2\sum_{u\in P_L}\left(\frac{\lambda_u}{1-\rho^2\lambda_{u}^2}+w_u\right)+o(\rho^{5/2})
\eeq

Since  $|g_u-g_0|\leq \const|u|$, we have $| w_u- g_0|u|^{-2}|\leq \const \rho^{-1/2}\eps_L^{-1}$. Then we can replace $w_u$ with $g_0|u|^{-2}$ in \eqref{temp8.75}. Setting  $u=\rho^{1/2}k$, we have, by definition of $\lambda$, 
\beqa
&&\lim_{\rho\to 0}(\rho^{1/2}|\Lambda|)^{-1}\sum_{u\in P_L}\left(\frac{\lambda_u}{1-\rho^2\lambda^2_u}+g_0|u|^{-2}\right)\\\nonumber
=&&\frac{1}{8\pi^3}\int_{k\in \R^3}g_0|k|^{-2}\left(\frac{\sqrt{1+4g_0|k|^{-2}}-1}{\sqrt{1+4g_0|k|^{-2}}}\right)dk^3=\frac{g_0^{3/2}}{\pi^2}
\eeqa
Inserting this result  into \eqref{temp8.75} and \eqref{temp8.73}, we have proved \eqref{proofthem1-4}.
\end{proof}
\section{Proof of Lemma \ref{lemHAS1}}

In this section, we  prove Lemma \ref{lemHAS1} concerning  potential energy terms with one $a_{0}$. 
%Without loss  of generality, we can assume the four-point-vertex 
%is written as $a^\dagger_{v_0}a^\dagger_{v_1}a_{v_2}a_{v_3}$ with  $v_0=0$ the only zero momentum.  
Let  $v_j\in \Lambda^*$ and $v_{j}\neq 0$ for $j = 1, 2, 3$.  Define $P_{H,\,c}$ as the following  subset of $P_H$:
\beq\label{defPHc}
P_{H,\,c}=\{k\in P_H: |k|\leq k_c\}.
\eeq
The following lemma classify all possible scenarios of $v_1,\,v_2,\,v_3$. Through out  this section, we assume that $v_{i} \not = 0$
for $i=1, 2, 3$. 

\begin{lem}\label{posscombone0}
Suppose  $\be, \alpha \in M$ and $\langle \al| a^\dagger_{0}a^\dagger_{v_1}a_{v_2}a_{v_3}|\be\rangle\neq 0$. 
Then there are only three possibilities:
\begin{enumerate}
	\item 
	\beq\label{one0firstcase}
	v_1\in P_L,\,\,\,v_2, v_3\in P_{H,\,c},   v_{i} \not = \pm v_{j} \text{ for } i \not = j .
	\eeq
		\item 
	\beq\label{one0secondcase}
	v_1\in P_{H,\,c}, \,\,\,v_{2}\in P_L, \,\,\,v_{3}\in P_{H,\,c},  v_{i} \not = \pm v_{j} \text{ for } i \not = j; \text{ \rm or }  2 \leftrightarrow 3.
	\eeq
	\item  
	\beq\label{one0thirdcase}
	v_1\in P_L,\,\,\,v_2\in P_L,\,\,\,v_3\in P_L.
	\eeq
\end{enumerate}
\end{lem}
\begin{proof}
Since particles with momenta in   $P_I$ are always created in pair, e.g.,  $(u,-u)$, 
either none of $v_i$'s belongs to $P_I$ or two of them belong to $P_I$. Thus  we have:
 \beq 
 v_1,v_{2}\in P_I\Rightarrow v_1=v_{2}, \text{ \rm or }  2 \leftrightarrow 3
 \eeq
 \beq
 v_2,v_3\in P_I\Rightarrow v_2=-v_3.
 \eeq
If two of $v_i$'s are in $P_I$, by the momentum conservation $v_1=v_2+v_3$ the other one must be equal to zero, which is  a  contradiction. Therefore
\beq
 v_i\notin P_I,\,\,\,{\rm for}\,\,\,1\leq i\leq 3 \, .
\eeq  
The restriction  $|v_i|\leq k_c$ follows from the construction of $M$. 
 Therefore, we have
\beq\label{PLPHc}
 v_i\in P_L\cup P_{H,c},\,\,\,{\rm for}\,\,\,1\leq i\leq 3
 \eeq  
 
Since particles in  $P_{H,c}$ are always created in soft pair creations which generated two particles in  
$P_{H,c}$, the number of particles in $P_{H,c}$ is even. So either none of $v_i$'s are in $P_{H,c}$ or two of them are in $P_{H,c}$. Together with \eqref{PLPHc}, and momentum conservation, we prove the lemma.
 \end{proof}
 
 
For fixed $v_1,v_2,v_3$, define 
\beq\label{defFal}
F(\al)\equiv \sum_{i: v_i\in P_L, i = 1, 2, 3} |\al(v_i)-\al(-v_i)|
\eeq
\begin{lem}\label{FalFbe0}
For any $\al, \be\in M$  if  $ \langle \al| a^\dagger_{0}a^\dagger_{v_1}a_{v_2}a_{v_3}|\be\rangle\neq 0$ and $v_i\neq \pm v_j$, we have:
 \beq\label{fab=no}
 F(\al)+F(\be)= \# \{i =1, 2, 3:  v_i\in P_L \}
 \eeq  
Furthermore,  the ratio between  $f(\al)$ and $f(\be)$ is bounded as follows. 
 \beq\label{relationalbe}
 \rho^{\frac{1}{20}}\sqrt{N}^{F(\al)-F(\be)}\leq \left|\frac{f(\be)\sqrt{\lambda_{v_1}}}{f(\al)\sqrt{\lambda_{v_2}\lambda_{v_3}\al(0)/|\Lambda|}}\right|\leq \rho^{\frac{-1}{20}}\sqrt{N}^{F(\al)-F(\be)}
 \eeq
% and the sign of $f(\al)f(\be)$ is equal to 
% \beq\label{signrelation}
% (-1)^{\sum_{v_i\in P_L}\left(\al^*(v_i)+\be^*(v_i)\right)}
% \eeq
\end{lem}
 \begin{proof} Since  $v_i\neq \pm v_j$,  for each $i$ fixed,  if $\alpha \in M^a_{v_i}$, 
 then $\beta \in  M^s_{v_i}$ and vice verse.    This proves  \eqref{fab=no}.
 
 
 Recall  the definition of $f$ in \eqref{deffal}. Then one can check the ratio involving $f(\be)/f(\al)$ in \eqref{relationalbe}  depends only on the last factor 
 \[
 \prod_{u\in P_L, \alpha^\ast (u) - \alpha (u) = 1}\sqrt{4\al^*(u)\lambda_u|\Lambda|^{-1}}
\]
We now use \eqref{lublambdaPL0} 
to bound $\lambda$ in this expression. Since $F(\alpha)$ counts how many times this factor appears, 
this proves \eqref{relationalbe}. 
 \end{proof}
 
 
Using the definitions of $\eta_L$ and $m_c$, the bound $\alpha(0)/|\Lambda| \le \rho$ 
and lemma \ref{id4a}, we have   
 \beqa\label{fal0}
 &&\left|f(\al)f(\be)\langle \al| a^\dagger_{0}a^\dagger_{v_1}a_{v_2}a_{v_3}|\be\rangle\right|\\\nonumber
 \leq&& \sqrt{N}^{F(\al)-F(\be)+1}\rho^{\frac{-1}{20}} \sqrt{ \rho \left|\frac{\lambda_{v_2}\lambda_{v_3}}{\lambda_{v_1}}\right|}\sqrt{\al(v_1)(\al(v_2)+1)(\al(v_3)+1)}|f(\al)|^2
 \eeqa
 and
 \beqa
 &&\label{fbe0}
 \left|f(\al)f(\be)\langle \al| a^\dagger_{0}a^\dagger_{v_1}a_{v_2}a_{v_3}|\be\rangle\right|\\\nonumber
 \leq&&\sqrt{N}^{F(\be)-F(\al)+1}\rho^{\frac{-1}{20}} \sqrt{\left|\lambda_{v_1}\lambda_{v_2}^{-1}\lambda_{v_3}^{-1}\right|\rho^{-1}}\sqrt{(\be(v_1)+1)\be(v_2)\be(v_3)}|f(\be)|^2
 \eeqa
 

Lemma \ref{lemHAS1} follows from summing the three inequalities of the following Lemma. 
Since $H_{A1}$ contains the complex conjugate of each of its terms, only the real parts of the expectation values below enter $\langle H_{A1}\rangle_\Psi$. These expectation values can be complex: for $v_2,v_3\in P_{H,c}$ the numbers $\lambda_{v_2}$ and $\lambda_{v_3}$ may have opposite signs, and then $\sqrt{\lambda_{v_2}}\sqrt{\lambda_{v_3}}$ (with the convention for $\sqrt{x}$, $x<0$, of \eqref{deffal}) is purely imaginary; see \eqref{temp8.103} below.
 \begin{lem}
In the limit $k_c\to \infty$, $\rho\to 0$, we have 
 \beq\label{0resultfirstcase}
	\overline\lim_{k_c\to\infty}\overline\lim_{\rho\to0}  |\Lambda|^{-2} \rho^{-5/2}\, {\rm Re}  \sum_{\eqref{one0firstcase}}\left\langle  V_{v_2}a^\dagger_0a^\dagger_{v_1}a_{v_2}a_{v_3}\right\rangle
	= - 2 \|Vw\|_1\frac{g_0^{3/2}}{3\pi^2}
	\eeq
	\beq\label{0resultsecondcase}
	\overline\lim_{k_c\to\infty}\overline\lim_{\rho\to0} |\Lambda|^{-2} \rho^{-5/2} \sum_{\eqref{one0secondcase}}\left|  V_{v_2} \left\langle a^\dagger_0a^\dagger_{v_1}a_{v_2}a_{v_3}\right\rangle\right|=0
	\eeq
	\beq\label{0resultthirdcase}
	\overline\lim_{k_c\to\infty}\overline\lim_{\rho\to0}|\Lambda|^{-2}\rho^{-5/2} \sum_{\eqref{one0thirdcase}}\left|  V_{v_2} \left  \langle a^\dagger_0a^\dagger_{v_1}a_{v_2}a_{v_3}\right\rangle\right|=0
\eeq
 \end{lem}
 \begin{proof} 
 We first  prove (\ref{0resultfirstcase}) concerning  \eqref{one0firstcase}, which implies that 
 $F(\alpha)+ F(\beta)=1$.   By the bounds on $\lambda_u$ in \eqref{lublambdaPL0} and $\al^*(u)\leq \lceil m_c\rceil$ for $u\in P_L$, we have, for $ F(\be)=0$  the following slightly modified version  of  \eqref{fbe0} 
 \beq\label{fbe0dup1}
 |f(\al)f(\be)|\left|\langle \al| a^\dagger_{0}a^\dagger_{v_1}a_{v_2}a_{v_3}|\be\rangle\right|\leq\rho^{\frac{-1}{10}} \rho^{-1}\sqrt{\left|\lambda_{v_2}^{-1}\lambda_{v_3}^{-1}\right|}\sqrt{\be(v_2)\be(v_3)}|f(\be)|^2
 \eeq
Here we replaced  $\rho^{-1/20}$ in \eqref{fbe0} by $\rho^{-1/10}$ to accommodate small errors. 
Summing over $\beta$ with $F(\be)=0$, we have 
\beq\label{sumfbe01}
\sum_{F(\be)=0}|f(\be)f(\al)|\left|\langle \al| a^\dagger_{0}a^\dagger_{v_1}a_{v_2}a_{v_3}|\be\rangle\right|\leq \rho^{-11/10}\sqrt{\left|\lambda_{v_2}^{-1}\lambda_{v_3}^{-1}\right|}Q_\Psi(v_2,v_3)
\eeq
Using the bound \eqref{boundpsifixeduv} on $Q_\Psi(v_2,v_3)$ (recall $v_2\neq\pm v_3$ in the case \eqref{one0firstcase})  and $|\lambda_u|\leq g_0|u|^{-2}$, we obtain that $\eqref{sumfbe01}=o(\rho^2)$. 

Since $F(\alpha)+ F(\beta)=1$, the other case is $F(\alpha)=0$.  
Hence we have 
\beqa\label{temp8.96}
&&\langle  a^\dagger_{0}a^\dagger_{v_1}a_{v_2}a_{v_3}\rangle= A_{1}+ A_{2}+ o(\rho^{3/2})\\\nonumber
&&A_{1} = \sum_{F(\al)=0}\sqrt{\al(0)\al(v_1)}\,\overline{f(\al)}f(\be)\\\nonumber
&&A_{2} = \sum_{F(\al)=0}\sqrt{\al(0)\al(v_1)}\left(\sqrt{(\al(v_2)+1)(\al(v_3)+1)}-1\right)\overline{f(\al)}f(\be)
\eeqa
Here $\be=\mathcal A^{v_1,\,v_2-v_1/2}\al$ is the state with $\langle \al| a^\dagger_{0}a^\dagger_{v_1}a_{v_2}a_{v_3}|\be\rangle\neq0$, and the complex conjugate comes from Lemma \ref{id4a}.
By the estimate  \eqref{relationalbe} and the Schwarz inequality 
$ | 2 (\sqrt{(a+1)(b+1)}-1)|\leq a+b$, we  have 
\begin{align}\label{temp8.100}
|A_{2}| & \le \rho^{4/5}\sum_{F(\al)=0} \frac {\al(v_2)+\al(v_3)}  2 |f(\al)|^2 \\
& \leq  \rho^{4/5}(Q_\Psi(v_2)+Q_\Psi(v_3)) \leq o(\rho^2)\, , 
\end{align}
where we have used the bounds on $\lambda$'s and $Q_\Psi(u)$ for $u\in P_H$. 

By the property  \eqref{properf4} for $f$, we have 
\beq
A_{1}= 2 \sqrt{\lambda_{v_2}}\sqrt{\lambda_{v_3}}\sum_{F(\al)=0}\al(0)\al(v_1)|\Lambda|^{-1}|f(\al)|^2,
\eeq
where the square roots are taken with the convention of \eqref{deffal} (in general $\sqrt{\lambda_{v_2}}\sqrt{\lambda_{v_3}}\neq\sqrt{\lambda_{v_2}\lambda_{v_3}}$).
We notice 
\beqa\nonumber
&&\sum_{F(\al)=0}\al(0)\al(v_1)|\Lambda|^{-1}|f(\al)|^2\\
&=&Q_{\Psi}(0,v_1)|\Lambda|^{-1}-\sum_{\al\in M_{v_1}^a}\al(0)\al(v_1)|\Lambda|^{-1}|f(\al)|^2
\eeqa


The absolute value of the second term is less than $\rho \lceil m_c\rceil \sum_{\al\in M_{v_1}^a}|f(\al)|^2$. By \eqref{boundMa}, it is less than $\rho^{7/4}$. Then  with $|\sqrt{\lambda_{v_2}}\sqrt{\lambda_{v_3}}|\leq O(\eps_H^{-2})$, we obtain 
\beq\label{temp8.103}
\langle  a^\dagger_{0}a^\dagger_{v_1}a_{v_2}a_{v_3}\rangle= 2 \sqrt{\lambda_{v_2}}\sqrt{\lambda_{v_3}}Q_{\Psi}(0,v_1)|\Lambda|^{-1}+o(\rho^{3/2}).
\eeq


Recall $\lambda_{u} = - w_{u}$ for $u \in P_{I}\cup P_{H}$ and $w_{u}= w_{-u}$ due to our assumption on $V$. 
Since  $v_{1} \in P_{L}$, i.e., $|v_1|\leq\eta_L^{-1}\rho^{1/2}$, and $v_{2}= -v_{3} + v_{1}$ and $v_{2} \in P_{H, c}$, we can check 
that 
\beq
\left|\lambda_{v_2}-\lambda_{v_3}\right|\leq \rho^{1/3}
\eeq
Inserting this in \eqref{temp8.103} and using the elementary inequality $|\sqrt{x}\sqrt{y}-x|\leq 2|x-y|$ for all real $x,y$ (with the convention of \eqref{deffal}), we arrive at 
\beq
\left\langle a^\dagger_0a^\dagger_{v_1}a_{v_2}a_{v_3}\right\rangle= 2 \lambda_{v_2}Q_\Psi(0,v_1)|\Lambda|^{-1}+o(\rho^{5/4})
\eeq
Since $\lambda_{v_2}=-w_{v_2}$ and $|\Lambda|^{-1}\sum_{v_2\in P_{H,c}}V_{v_2}w_{v_2}\to\|Vw\|_1$, in the limit $k_c\to\infty,\rho\to 0$ we have
\beqa\nonumber
&&|\Lambda|^{-2}\sum_{v_1\in P_L, v_2\in P_{H,c}}\left\langle V_{v_2} a^\dagger_0a^\dagger_{v_1}a_{v_2}a_{v_3}\right\rangle\\
&=&-2\|Vw\|_1 |\Lambda|^{-2} \sum_{v_1\in P_L}Q_\Psi(0,v_1)+o(\rho^{5/2})
\eeqa
We note
\beq
|\Lambda|^{-2} \sum_{v_1\in P_L}Q_\Psi(0,v_1)=\rho |\Lambda|^{-1}Q_\Psi(0)- |\Lambda|^{-2}Q_\Psi(0,0)-
|\Lambda|^{-2} \sum_{u\in P_I\cup P_H}Q_\Psi(0,u)
\eeq
The last term is less than $N|\Lambda|^{-2} \sum_{u\in P_I\cup P_H}Q_\Psi(u)\le o(\rho^{5/2})$ by Theorem \ref{thmA}.  Together with Lemma \ref{boundspsi02}, \ref{boundspsi0} on $Q_\Psi(0,0)$ and $Q_\Psi(0)$, we can compute the first two terms, i.e.,
\beq\label{sum0PL}
|\Lambda|^{-2} \sum_{v_1\in P_L}Q_\Psi(0,v_1)=\rho_0(\rho-\rho_0)+o(\rho^{5/2})
\eeq
This yields  \eqref{0resultfirstcase}.


We next prove \eqref{0resultsecondcase} concerning  \eqref{one0secondcase}. Without loss of generality we assume  that 
\beq
v_{1,3}\in  P_{H,c}\rmand v_2\in P_L
\eeq
Following similar arguments  in the previous proof, i.e.,  using  Lemma \ref{id4a}, \eqref{fal0} or  \eqref{fbe0} and the bounds on $\lambda_u$'s, we have
\beqa
|\langle a^\dagger_0a^\dagger_{v_1}a_{v_2}a_{v_3}\rangle|&&\leq \sum_{F(\al)=0}\rho^{-\frac{1}{10}} \sqrt{\al(v_1)(\al(v_3)+1)\left|\lambda_{v_1}^{-1}\right|}|f(\al)|^2\\\nonumber
+&&\sum_{F(\be)=0}\rho^{-\frac{1}{10}} \sqrt{\be(v_3)(\be(v_1)+1)\left|\lambda_{v_3}^{-1}\right|}|f(\be)|^2
\eeqa
For the upper bound, we can replace $\sum_{F(\al)=0}$ by  $\sum_{\al \in M}$. Using the upper bounds (\ref{ubpsiuroughPH}) and  (\ref{boundpsifixeduv}) on $Q_\Psi(u)$ and  $Q_\Psi(u,v)$  for $u,v\in P_H$ (with $u=v_1$, $v=v_3$; note $v_1\neq\pm v_3$ in the case \eqref{one0secondcase}), we obtain 
$
\left|\left\langle a^\dagger_0a^\dagger_{v_1}a_{v_2}a_{v_3}\right\rangle\right| \leq \const\rho^{3/2}
$. 
This proves  (\ref{0resultsecondcase}).

We now prove \eqref{0resultthirdcase} concerning   $v_i\in P_L$ satisfying   $F(\al)+F(\be)=3$.   
It is easy to prove that the contribution from the special cases,  $v_1=-v_2$(or $v_3$) or $v_2=v_3$,  is negligible, 
\beq
\overline\lim_{\rho}\sum_{\rm special\,\,\,cases}
  \left |  V_{v_2} \langle a^\dagger_0a^\dagger_{v_1}a_{v_2}a_{v_3}\rangle \right | 
\rho^{-5/2}|\Lambda|^{-2}=0
\eeq
So from now on we assume that $v_i\neq \pm v_j$ for $i\neq j$.  As before, we rewrite  $\langle a^\dagger_0a^\dagger_{v_1}a_{v_2}a_{v_3}\rangle$ by using  Lemma \ref{id4a} and \eqref{fal0} or \eqref{fbe0}.  Together with the bounds on $\lambda_u$'s  and $\al(v_i)\leq \lceil m_c\rceil$, we have
\beqa
|\langle a^\dagger_0a^\dagger_{v_1}a_{v_2}a_{v_3}\rangle|
\leq &&\sum_{F(\al)=0}N^{-1}\rho^{-\frac{1}{10}} |f(\al)|^2+\sum_{F(\al)=1}\rho^{-\frac{1}{10}} |f(\al)|^2\\\nonumber
+&&\sum_{F(\be)=0}N^{-1}\rho^{-\frac{1}{10}} |f(\be)|^2+\sum_{F(\be)=1}\rho^{-\frac{1}{10}} |f(\be)|^2
\eeqa
By symmetry, we only need to estimate the first two terms on the rhs. The first term is less than $N^{-1}\rho^{-\frac{1}{10}}$. For the second term,  we note $F(\al)=1$ implies  that there exists $i, 1\leq i\leq 3$
such that  $\al \in M_{v_i}^a$.  By \eqref{boundMa}, we have
\beq
\sum_{F(\al)=1}|f(\al)|^2\leq \rho^{3/4}
\eeq
This implies $|\langle a^\dagger_0a^\dagger_{v_1}a_{v_2}a_{v_3}\rangle|\le \rho^{1/2}$ and \eqref{0resultthirdcase}, which 
complete the proof.
\end{proof}
 
 




 
\section{Interaction Energy with Four Nonzero Momenta: The Classification} 

In the next three sections, we will prove Lemma \ref{lemHAS2} involving interaction energy without $a_0$. 
We will show that the only contribution to the accuracy we need comes from four high momentum particles, 
to  be computed in next section. In this section, we start the procedure of identifying the error terms. 

For $\al, \be \in M$, we have the following lemma, similar to  Lemma \ref{posscombone0} and Lemma \ref{FalFbe0}. 
Since it can be proved by same method, we will only state the result. 

\begin{lem}\label{fourlegslemma1}
Suppose   $v_i\neq 0, 1\leq i\leq 4$ and  $v_{1} + v_{2} \not = 0$, $v_1\neq v_3$ and $v_1\neq v_4$. If  $\langle \al| a^\dagger_{v_{1}}a^\dagger_{v_2}a_{v_3}a_{v_4}|\be\rangle\neq 0$ for some $\alpha, \beta \in M$, then there are exactly the following five cases:
\begin{enumerate}
	\item All of $v_i\in P_L$ for $1\leq i\leq 4$.
	\item
%	\beq\label{no0secondcase}
$	v_{1}, v_{2}\in P_{L}, \,\,\,v_{3}, v_{4}\in P_{H,\,c}$.
%	\eeq
	\item[$2'$.]
$	v_{1}, v_{2}\in P_{H,\,c}, \,\,\,v_{3}, v_{4}\in P_{L}$.
	\item One of $v_{1}, v_{2}$ is in $P_L$ and the other is in $P_{H,c}$;  one of $v_{3}, v_{4}$ is in $P_L$ and the other is in $P_{H,c}$.
	\item All of $v_i\in P_{H,c}$ for $1\leq i\leq 4$.
\end{enumerate}
If $v_i\neq \pm v_j$, for $1\leq i<j\leq 4$, we have
\beq 
\rho^{\frac{1}{20}}\sqrt{N}^{F(\al)-F(\be)}\leq \left|\frac{f(\be)\sqrt{\lambda_{v_1}\lambda_{v_2}}}{f(\al)\sqrt{\lambda_{v_3}\lambda_{v_4}}}\right|\leq \rho^{\frac{-1}{20}}\sqrt{N}^{F(\al)-F(\be)},
\eeq
%With Lemma \ref{id4a}, it implies 
 \beqa\label{no0Fal}
&&\left|f(\al)f(\be)\langle \al| a^\dagger_{v_1}a^\dagger_{v_2}a_{v_3}a_{v_4}|\be\rangle\right|\\\nonumber
 &&\leq \sqrt{N}^{F(\al)-F(\be)}\rho^{\frac{-1}{20}} \sqrt{\frac{\lambda_{v_3}\lambda_{v_4}}{\lambda_{v_1}\lambda_{v_2}}}\sqrt{\al(v_1)\al(v_2)(\al(v_3)+1)(\al(v_4)+1)}|f(\al)|^2
\eeqa
and
\beqa\label{no0Fbe}
&&\left|f(\al)f(\be)\langle \al| a^\dagger_{v_1}a^\dagger_{v_2}a_{v_3}a_{v_4}|\be\rangle\right|\\\nonumber
 &&\leq \sqrt{N}^{F(\be)-F(\al)}\rho^{\frac{-1}{20}} \sqrt{\frac{\lambda_{v_1}\lambda_{v_2}}{\lambda_{v_3}\lambda_{v_4}}}\sqrt{(\be(v_1)+1)(\be(v_2)+1)\be(v_3)\be(v_4)}|f(\be)|^2.
\eeqa
\end{lem}




\begin{prop} For $u\in P_L$ and $v\in P_{H,c}$, we have, for $\rho$ small enough (depending on $k_c$), the following inequality
\beq\label{temp8.119}
\sum_{\al\in M_u^a}\al(v)\left|f(\al)^2\right|\leq \left|\lambda_v\right|\rho^{3-\frac{1}{10}}
\eeq\end{prop}
\begin{proof} By definition of $M$ \eqref{defMsau}, for any $\al\in M_u^a$, there exist $\be\in M_u^s$, a sign $\sigma\in\{1,-1\}$ and $k$ such that $\mathcal A^{\sigma u,k}\be=\al$ and $\pm k+\sigma u/2\in P_{H,c}$; here $\sigma=1$ if $\al(-u)=\al(u)+1$ and $\sigma=-1$ if $\al(u)=\al(-u)+1$, since the soft pair creation annihilates a particle of momentum $\sigma u$. Clearly, for any $v \in P_{H}$ we have
$\alpha(v) \le \beta(v) +1$ and the case we need the constant $1$ occurs only when  $v=k+\sigma u/2 $ or
$v= - k+\sigma u/2$. The two values of $\sigma$ are treated in the same way, and we write the details for $\sigma=1$. The contribution of $\sigma=1$ to the left hand side of \eqref{temp8.119} is bounded by
\beq\label{10.1}
 \sum_{\be}\sum_{k:\pm k+u/2\in P_{H,c}}\be(v)|f(\mathcal A^{u,k}\be)^2|+\sum_{\be}\sum_{k:\pm k+u/2=v}|f(\mathcal A^{u,k}\be)^2|
\eeq
Recall   \eqref{properf4},   $\be(0)\leq N$ and $\be(u)\leq \lceil m_c\rceil\leq 2m_c$ imply that
\beq
|f(\mathcal A^{u,k}\be)|^2\leq 8|f(\be)|^2\rho m_c|\Lambda|^{-1}\left|\lambda_{k+u/2}\lambda_{-k+u/2}\right|
\eeq
By  the bound \eqref{ublambdaPH0} on $\lambda$, we obtain that
\beqa\nonumber
\eqref{10.1}&&\leq \sum_{\be}\be(v)|f(\be)|^2\frac{8\rho m_c}{|\Lambda|}\left[\sum_{\pm k+u/2\in P_{H,c}}\left|\lambda_{k+u/2}\lambda_{-k+u/2}\right|\right]\\\nonumber
&&\quad+\const\left|\lambda_v\right||\Lambda|^{-1}\\
&&\leq \const Q_\Psi(v)\rho m_c\eps_H^{-4}k_c^3+\const\left|\lambda_v\right||\Lambda|^{-1}
\eeqa
Using the bound \eqref{ubpsiuroughPH} on $Q_\Psi(v)$, the same bound for $\sigma=-1$ and $\rho$ small enough (depending on $k_c$), we have proved  \eqref{temp8.119}.
\end{proof}

\begin{lem}We have the following estimates on the interaction energies: 
\beq\label{resultasymno24L}
\overline\lim_{k_c\to\infty}\,\overline\lim_{\rho\to0}
 \; \rho^{-5/2} {|\Lambda|}^{-2}  \sum_{v_1,v_2,v_3,v_4\in P_L} \left | V_{v_{1}-v_{3}}\left\langle a_{v_1}^\dagger a_{v_2}^\dagger a_{v_3}a_{v_4}\right\rangle_{\Psi}\right| =0,
\eeq
\beq\label{resultasymno22L2H1}
\overline\lim_{k_c\to\infty}\,\overline\lim_{\rho\to0}
\; \rho^{-5/2} {|\Lambda|}^{-2}  \sum_{\substack{v_1+v_2\neq0, v_1,v_2\in P_L\\ v_3,v_4\in P_H}}  \left | V_{v_{1}-v_{3}}\left\langle a_{v_1}^\dagger a_{v_2}^\dagger a_{v_3}a_{v_4}\right\rangle_{\Psi}\right|  = 0
\eeq
\beq\label{resultasymno22L2H2}
\overline\lim_{k_c\to\infty}\,\overline\lim_{\rho\to0}
 \; \rho^{-5/2} {|\Lambda|}^{-2}  \sum_{v_1,v_3\in P_L,v_2,v_4\in P_H} \left | V_{v_{1}-v_{3}}\left\langle a_{v_1}^\dagger a_{v_2}^\dagger a_{v_3}a_{v_4}\right\rangle_{\Psi}\right|  =0
\eeq
The same estimates hold for the sums over the configurations obtained from those in \eqref{resultasymno22L2H1} and \eqref{resultasymno22L2H2} by exchanging $(v_1,v_2)$ with $(v_3,v_4)$, or $v_1$ with $v_2$, or $v_3$ with $v_4$.
In other words, the contributions from cases 1, 2, $2'$ and 3 in Lemma \ref{fourlegslemma1} are negligible for our purpose.
\end{lem}
\begin{proof}
We first prove the \eqref{resultasymno24L} concerning  $v_i\in P_L$. By Lemma \ref{id4a}, we have 
\beq
\left|\left\langle a_{v_1}^\dagger a_{v_2}^\dagger a_{v_3}a_{v_4}\right\rangle_{\Psi}\right|\leq \sum_{\al}\left|f(\al)
f(T(\al))\right| m_c^4
\eeq
Using the Schwarz inequality, we have $\left|\left\langle a_{v_1}^\dagger a_{v_2}^\dagger a_{v_3}a_{v_4}\right\rangle_{\Psi}\right|\leq m_c^4$. The summation  over the $v_i$ with $v_i=\pm v_j$ for some $1\leq i<j\leq 4 $
is negligible in the sense that  
\beq\label{specialviPL}
{|\Lambda|}^{-2} \sum_{v_1,v_2,v_3,v_4\in P_L, v_i=\pm v_j} \left | \left\langle a_{v_1}^\dagger a_{v_2}^\dagger a_{v_3}a_{v_4}\right\rangle_{\Psi}  \right | \leq o(\rho^{5/2})
\eeq
{F}rom now on, we assume that  $v_i\neq \pm v_j$ for any  $1\leq i<j\leq 4 $.  

Using \eqref{no0Fal}, \eqref{no0Fbe}  and the bounds  \eqref{lublambdaPL0} on $\lambda$, we have 
\beqa\nonumber
\left|\left\langle a_{v_1}^\dagger a_{v_2}^\dagger a_{v_3}a_{v_4}\right\rangle_{\Psi}\right|&&\leq \sum_{F(\al)\leq 1}\rho^{\frac{-1}{10}}N^{-1}|f(\al)^2|\\\nonumber
&&+\sum_{F(\al)=2}\rho^{\frac{-1}{10}}|f(\al)^2|+\sum_{F(\be)\leq 1}\rho^{\frac{-1}{10}}N^{-1}|f(\be)^2|
\eeqa
By \eqref{boundMaMa}, we have $\left|\left\langle a_{v_1}^\dagger a_{v_2}^\dagger a_{v_3}a_{v_4}\right\rangle_{\Psi}\right|\leq \rho^{9/5} .  $ 
 %{\bf ($1+ 1/8$ is the min needed)} 
Together with \eqref{specialviPL} and  $\Lambda=\rho^{-25/8}$, we can sum over $v_{j}$ to have  
\beq
{|\Lambda|}^{-2} \sum_{v_1,v_2,v_3,v_4\in P_L}\left\langle a_{v_1}^\dagger a_{v_2}^\dagger a_{v_3}a_{v_4}\right\rangle_{\Psi}\leq o(\rho^{5/2})
\eeq


Before proving \eqref{resultasymno22L2H1} and \eqref{resultasymno22L2H2}, we dispose of the terms with $v_i=\pm v_j$ for some $i\neq j$, to which \eqref{no0Fal} and \eqref{no0Fbe} do not apply. A momentum in $P_L$ and a momentum in $P_{H}$ are never equal up to sign. In the case of \eqref{resultasymno22L2H1}, $v_3=\pm v_4$ is impossible, since $v_3+v_4=v_1+v_2\neq 0$ and $v_3=v_4$ would give $|v_3|=|v_1+v_2|/2<\eps_H$; so the only coincidence is $v_1=v_2$. In the case of \eqref{resultasymno22L2H2}, the only coincidences are $v_1=\pm v_3$, since $v_2=v_4$ implies $v_1=v_3$ and $v_2=-v_4$ would give $|v_1-v_3|=2|v_4|>2\eps_H$. In both cases the term is determined by one momentum in $P_L$ and one in $P_{H,c}$, so there are at most $\const \eta_L^{-3}\rho^{3/2}k_c^3|\Lambda|^2$ such terms. By the Schwarz inequality,
\[
\left|\left\langle a_{v_1}^\dagger a_{v_2}^\dagger a_{v_3}a_{v_4}\right\rangle_{\Psi}\right|\leq \|a_{v_1}a_{v_2}\Psi\|\,\|a_{v_3}a_{v_4}\Psi\|\leq \const m_c\rho^2\eps_H^{-4} .
\]
Here we have used, for $v_1=v_2$, $\|a_{v_1}a_{v_1}\Psi\|\leq 2m_c$ (since $\al(u)\leq\lceil m_c\rceil\leq 2m_c$ for $u\in P_L$) and $\|a_{v_3}a_{v_4}\Psi\|^2=Q_\Psi(v_3,v_4)$ with \eqref{boundpsifixeduv}; and, for $v_1=\pm v_3$, $\|a_{v_1}a_{v_2}\Psi\|^2=Q_\Psi(v_1,v_2)\leq 2m_cQ_\Psi(v_2)$, the same bound for $\|a_{v_3}a_{v_4}\Psi\|^2$, and \eqref{ubpsiuroughPH}. Since $|V_u|\leq V_0$, the contribution of these terms to the left hand sides of \eqref{resultasymno22L2H1} and \eqref{resultasymno22L2H2} is bounded by $\const k_c^3\rho^{1-8\eta}$ before the limits, which vanishes as $\rho\to0$. {F}rom now on, we assume that  $v_i\neq \pm v_j$ for any  $1\leq i<j\leq 4 $.

We now prove \eqref{resultasymno22L2H1} concerning  $v_{1,2}\in P_L$ and $v_{3,4}\in P_{H,c}$.  As before, by \eqref{no0Fal}, \eqref{no0Fbe},  \eqref{lublambdaPL0} and \eqref{ublambdaPH0}, we have
\beqa\nonumber
&&\left|\left\langle a_{v_1}^\dagger a_{v_2}^\dagger a_{v_3}a_{v_4}\right\rangle_{\Psi}\right|\leq\sum_{F(\al)= 0}N^{-1}\rho^{\frac{9}{10}}\sqrt{(\al(v_3)+1)(\al(v_4)+1)}|f(\al)|^2\\\nonumber
+&&\sum_{F(\be)\leq 1}\rho^{\frac{-11}{10}}\sqrt{\frac{\be(v_3)\be(v_4)}{|\lambda_{v_3}\lambda_{v_4}|}}|f(\be)|^2
\eeqa
Since $\sqrt{(x+1)(y+1)}\leq 1+x+y$ for $x,y\geq0$ and $Q_\Psi(v_3), Q_\Psi(v_4)\leq 1$ by \eqref{ubpsiuroughPH}, the first term in rhs. is bounded by $N^{-1}\rho^{\frac{9}{10}}\left(1+Q_\Psi(v_3)+Q_\Psi(v_4)\right)\leq 3\rho^3$. Since $v_3, v_4\in P_H$ and $v_3\neq\pm v_4$, by \eqref{boundpsifixeduv} we obtain that the second term in rhs. is $o(\rho^{11/4})$. So
\beq
\left|\left\langle a_{v_1}^\dagger a_{v_2}^\dagger a_{v_3}a_{v_4}\right\rangle_{\Psi}\right|\leq \rho^{\frac{11}4}
\eeq
%{\bf ($5/2+ 1/8$ is the min needed)}
Summing over $v_{j}$'s, we have proved  \eqref{resultasymno22L2H1}.


Finally, we prove \eqref{resultasymno22L2H2} concerning  $v_{1,3}\in P_L$ and $v_{2,4}\in P_H$. Again, with \eqref{no0Fal}, \eqref{no0Fbe} and the bounds on $\lambda$'s in \eqref{lublambdaPL0} and \eqref{ublambdaPH0}, we have 
\beqa\label{temp8.128}
&&\left|\left\langle a_{v_1}^\dagger a_{v_2}^\dagger a_{v_3}a_{v_4}\right\rangle_{\Psi}\right|\leq 
Q_{1}+ Q_{2}+ Q_{3} \\
Q_{1} &= &\sum_{F(\al)= 0}N^{-1}\rho^{\frac{-1}{10}}\sqrt{\frac{\al(v_2)(\al(v_4)+1)}{\left|\lambda_{v_2}\right|}}|f(\al)|^2\\\nonumber
Q_{2}&=&\sum_{F(\be)= 0}N^{-1}\rho^{\frac{-1}{10}}\sqrt{\frac{\be(v_4)(\be(v_2)+1)}{\left|\lambda_{v_4}\right|}}|f(\be)|^2\\
Q_{3} & = & \sum_{F(\al)= 1}\rho^{\frac{-1}{10}}\sqrt{\frac{\al(v_2)(\al(v_4)+1)}{\left|\lambda_{v_2}\right|}}|f(\al)|^2
\eeqa 
By  Theorem \ref{thmA} and the fact $\sqrt{x}\leq x$ for $x\in\N$, we have 
$$
Q_{1}\le  N^{-1}\rho^{\frac{-1}{10}}\lambda_{v_2}^{-1/2}\left(Q_{\Psi}(v_2)+Q_{\Psi}(v_2,v_4)\right)
\le \rho^{3},
$$ 
where we have used the bounds \eqref{ubpsiuroughPH} and \eqref{boundpsifixeduv} on $Q_\Psi(u)$ and $Q_\Psi(u,v)$. Similarly,  we have  $Q_{2}\le \rho^3$.  
Again using the fact $\sqrt{x}\leq x$ for $x\in\N$ , we have 
\beqa\nonumber
Q_{3} & \le & \sum_{F(\al)= 1}\rho^{\frac{-1}{10}}\al(v_2)\left|\lambda_{v_2}\right|^{-1/2}|f(\al)|^2+\rho^{-\frac{1}{10}}\left|\lambda_{v_2}\right|^{-1/2}Q_{\Psi}(v_{2},v_{4}) \\
&\le & \sum_{F(\al)= 1}\rho^{\frac{-1}{10}}\al(v_2)\left|\lambda_{v_2}\right|^{-1/2}|f(\al)|^2+\rho^{3} , \nonumber
\eeqa
where we have used  \eqref{boundpsifixeduv}.
We can estimate the first term in rhs. by  \eqref{temp8.119}. Collecting all these bounds, we have proved that 
\beq
\left|\left\langle a_{v_1}^\dagger a_{v_2}^\dagger a_{v_3}a_{v_4}\right\rangle_{\Psi}\right| \leq \rho^{2.7}
\eeq
Summing over $v_{j}$'s, we have proved  \eqref{resultasymno22L2H2}.

The other configurations listed after \eqref{resultasymno22L2H2} are reduced to the two cases above. Since $a^\dagger_{v_1}a^\dagger_{v_2}=a^\dagger_{v_2}a^\dagger_{v_1}$, $a_{v_3}a_{v_4}=a_{v_4}a_{v_3}$ and the proofs above only use $|V_u|\leq V_0$, they apply when $v_1$ and $v_2$, or $v_3$ and $v_4$, are exchanged. For case $2'$, i.e., $v_1,v_2\in P_H$ and $v_3,v_4\in P_L$, we use
\[
\left\langle a_{v_1}^\dagger a_{v_2}^\dagger a_{v_3}a_{v_4}\right\rangle_{\Psi}=\overline{\left\langle a_{v_3}^\dagger a_{v_4}^\dagger a_{v_1}a_{v_2}\right\rangle_{\Psi}}, \qquad V_{v_1-v_3}=V_{v_3-v_1},
\]
so that the corresponding sum is equal to the sum in \eqref{resultasymno22L2H1}.
\end{proof}


\section{Interaction  Energy with Four High Momentum Legs I: The Main Term} 


We now  estimate of the interaction energy in the  case 4  of Lemma \ref{fourlegslemma1}, i.e., 
$k_i, i=1, 2, 3, 4 $ satisfy  
\beq\label{conditionk1}
k_1+k_2=k_3+k_4,\,\,\, k_1+k_2\neq 0,\,\,\,k_1\neq k_3,\,\,\,k_1\neq k_4, \,\,\,k_i\in P_{H,c}
\eeq
 In the remainder of this paper,  all $p_i$'s, $q_i$'s, $k_i$'s belong to $P_{H,c}$ and $u_i$, $v_i$'s belong to $P_L$.  We start with some special cases.

\begin{lem} Suppose  $k_i$ satisfy  \eqref{conditionk1}. Then  we have
\beqa\label{4kspecialresult1}
&&\sum_{k_1,k_3}\left|  V_{k_{1}-k_{3}}\left\langle a^\dagger_{k_1}a^\dagger_{k_1}a_{k_3}a_{k_4}\right\rangle\right|=o\left(\rho^{5/2}|\Lambda|^2\right)\\\label{4kspecialresult2}
&&\sum_{k_1,k_2}\left|V_{2k_{1}} \left\langle a^\dagger_{k_1}a^\dagger_{k_2}a_{-k_1}a_{k_4}\right\rangle\right|=o\left(\rho^{5/2}|\Lambda|^2\right)
\eeqa
\end{lem}

\begin{proof} We first show that for $p,q\in P_{H,c}$ with $p+q\neq 0$, 
\beq
\|a_pa_q\Psi\|^2\leq \const\rho^4\eps_H^{-8}.
\eeq
For $p\neq \pm q$, $\|a_pa_q\Psi\|^2=Q_\Psi(p,q)$ and this follows from \eqref{boundpsifixeduv} and \eqref{ublambdaPH0}. For $p=q$, $\|a_pa_p\Psi\|^2=\sum_{m\geq 2}m(m-1)Q_\Psi(\{p,m\})$ and this follows from \eqref{temp5.19} and \eqref{ublambdaPH0}. 
Since $\langle a^\dagger_{k_1}a^\dagger_{k_2}a_{k_3}a_{k_4}\rangle=\langle a_{k_2}a_{k_1}\Psi,\, a_{k_3}a_{k_4}\Psi\rangle$, the Schwarz inequality gives, for any $k_i\in P_{H,c}$ with $k_1+k_2=k_3+k_4\neq 0$,
\beq
\left|\left\langle a^\dagger_{k_1}a^\dagger_{k_2}a_{k_3}a_{k_4}\right\rangle\right|\leq \|a_{k_1}a_{k_2}\Psi\|\,\|a_{k_3}a_{k_4}\Psi\|\leq \const\rho^4\eps_H^{-8}.
\eeq
The summation in \eqref{4kspecialresult1} is over $k_1,k_3\in P_{H,c}$ (with $k_4=2k_1-k_3$) and the summation in \eqref{4kspecialresult2} is over $k_1,k_2\in P_{H,c}$ (with $k_4=2k_1+k_2$), so each of them has at most $(\const k_c^3|\Lambda|)^2$ terms. Since $|V_u|\leq V_0$ and $\eps_H=\rho^\eta$, the left hand sides of \eqref{4kspecialresult1} and \eqref{4kspecialresult2} are bounded by
\beq
\const V_0 k_c^6|\Lambda|^2\rho^{4-8\eta}=o\left(\rho^{5/2}|\Lambda|^2\right). 
\eeq
 \end{proof}
 
 
The same proof applies to all $(k_1, \ldots, k_4)$ satisfying \eqref{conditionk1} with $k_i=\pm k_j$ for some $i\neq j$: such a term is determined by two of the $k_i$'s, so there are at most $(\const k_c^3|\Lambda|)^2$ of them. Hence all these special cases, such as $k_1=-k_4$, are negligible. So from now on we focus on the cases
\beq\label{conditionk2}
k_1+k_2=k_3+k_4,\,\,\,k_i\in P_{H,c},\,\,\,k_i\neq\pm k_j\,\,\,{\rm for}\,\,\, i\neq j
\eeq 
This condition will be imposed for the rest of this section. 
Denote by  $M(k_1,k_2)$ the set of all states created by  a soft pair creation $\mathcal A^{k_1+k_2,\,\,\, k_1/2-k_2/2}$ from another  state, i.e., 
\beq
M(k_1,k_2)\equiv\left\{\be\in M | \exists \al\in M^s_{k_1+k_2} {\rm \; such \; that \;}  \mathcal A^{k_1+k_2,\,\,\, k_1/2-k_2/2}\al=\be\right\}
\eeq
if $k_1+k_2 \in P_L$. Otherwise, we set  $M(k_1,k_2)=\emptyset$.
Notice that $$\left|\mathcal A^{k_1+k_2,\,\,\, k_1/2-k_2/2} \alpha\right\rangle = C  a^\dagger_{k_1}a^\dagger_{k_2}a_{k_1+k_2}a_0 | \alpha\rangle$$
for some normalization constant $C$. Hence for $\beta, \gamma \in M$, if $$\left\langle\be| a^\dagger_{k_1}a^\dagger_{k_2}a_{k_3}a_{k_4}|\gamma\right\rangle\neq 0,$$ 
 %if we let 
% \[
%\alpha= C a_{k_3}a_{k_4} |\gamma\rangle  \in M^{s}_{k_{1}+k_{2}} , 
%\]
we have $k_1+k_2=k_3+k_4$ and 
\beq\label{11.1}
\mathcal A^{k_1+k_2,\,\,\, k_1/2-k_2/2}\al=\be\Leftrightarrow  \; \mathcal A^{k_3+k_4,\,\,\, k_3/2-k_4/2}\al=\gamma,   
\eeq

The main contribution of the four nonvanishing leg term is identified in the next lemma. 

\begin{lem}\label{101}
\beqa\nonumber
&&\overline\lim_{k_c\to\infty}\,\overline\lim_{\rho\to0}    \rho^{-5/2}|\Lambda|^{-2}  \sum_{\eqref{conditionk2}}\sum_{\be\in M(k_1,k_2)}\\\nonumber
&&\qquad\qquad V_{k_1-k_3}\overline{f(\be)}f(\gamma)\left\langle\be| a^\dagger_{k_1}a^\dagger_{k_2}a_{k_3}a_{k_4}|\gamma\right\rangle
\\\label{resultinMk1k2}
&&\leq 4 \|w^2V\|_1 \frac{1}{3\pi^2}g_0^{3/2}
\eeqa
The sum on the left side is real: by \eqref{11.1}, complex conjugation maps the term of $(k_1,k_2,k_3,k_4,\be,\gamma)$ to the term of $(k_3,k_4,k_1,k_2,\gamma,\be)$. 
\end{lem}
\begin{proof}
By  \eqref{11.1}, we have 
\beqa\label{temp8.140}
&&\sum_{\be\in M(k_1,k_2)}\overline{f(\be)}f(\gamma)\left\langle\be| a^\dagger_{k_1}a^\dagger_{k_2}a_{k_3}a_{k_4}|\gamma\right\rangle\\\nonumber
=&&\theta\sum_{\al\in M^s_{k_1+k_2}}4|f(\al)|^2|\Lambda|^{-2}\al(0)\al(k_1+k_2)\prod_{i=1}^4\sqrt{(\al(k_i)+1)}
\eeqa
where we have used \eqref{properf4} and the notation
\[
\theta\equiv\overline{\sqrt{\lambda_{k_1}}\sqrt{\lambda_{k_2}}}\,\sqrt{\lambda_{k_3}}\sqrt{\lambda_{k_4}} 
\]
(recall the convention $\sqrt{x}=\sqrt{|x|}i$ for $x<0$ in \eqref{deffal}). We claim that \eqref{temp8.140} is very close to the following expression:
\beqa\label{temp5.145}
\theta\sum_{\al\in  M^s_{k_1+k_2} }4|f(\al)|^2|\Lambda|^{-2}\al(0)\al(k_1+k_2)
\eeqa 
For $x_i\geq 0$, we have  
\beq\label{inequalityx1234}
1\leq \sqrt{(x_1+1)(x_2+1)(x_3+1)(x_4+1)}\leq \frac{1}{4}(x_1+x_2+2)(x_3+x_4+2), 
\eeq 
Since $\al(0)\leq N$ and $\al(k_1+k_2)\leq \lceil m_c\rceil\leq 2m_c$, we have 
\beq\label{temp8.144}
\frac{\left|\eqref{temp8.140}-\eqref{temp5.145}\right|}{\left|\theta\right|}\leq \frac{8m_c\rho}{|\Lambda|}\left(\sum_{i}Q_\Psi(k_i)+\sum_{i,j}Q_\Psi(k_i, k_j)\right)
\le \frac{\rho^2}{|\Lambda|}
\eeq
where we have used \eqref{ubpsiuroughPH} and  \eqref{boundpsifixeduv}.

By definition, $Q_\Psi(0,k_1+k_2)=\sum_{\al\in M}\al(0)\al(k_1+k_2)|f(\al)|^2$. Together  with  $\al(0)\leq N$ and $\al(k_1+k_2)\leq 2m_c$, we have 
\beq
\left|\frac{\eqref{temp5.145}}{\theta}-4|\Lambda|^{-2}Q_\Psi(0,k_1+k_2)\right|\leq\frac{8m_c\rho}{|\Lambda|}\sum_{\al\in M^a_{k_1+k_2}}|f(\al)|^2
\eeq
Using the bound \eqref{boundMa} concerning $\sum_{\al\in M^a_{k_1+k_2}}$, we have 
\beq\label{temp8.146}
\left|\frac{\eqref{temp5.145}}{\theta}-4|\Lambda|^{-2}Q_\Psi(0,k_1+k_2)\right|\leq\rho^{3/2}|\Lambda|^{-1}
\eeq
Combining \eqref{temp8.144},  \eqref{temp8.146}, with the bounds on $\lambda$ in \eqref{ublambdaPH0}, we have:
\beq\label{temp8.147}
\left|\eqref{temp8.140}-4\theta|\Lambda|^{-2}Q_\Psi(0,k_1+k_2)\right|\leq \frac{\rho^{5/4}}{|\Lambda|}
\eeq

Since $\lambda_p=-w_p=-g_p|p|^{-2}$ for $p\in P_H$, $\lambda_{-q}=\lambda_q$ ($V$, and hence $g$, is even), and $|\nabla_r(g_r|r|^{-2})|\leq \const\eps_H^{-3}$ for $|r|\geq\eps_H/2$ ($g_r$ and $\nabla_rg_r$ are bounded), we have for $p, q\in P_{H,c}$ with $p+q\in P_L$, since the segment from $p$ to $-q$ stays in $\{|r|\geq \eps_H/2\}$ and $|p+q|\leq\eta_L^{-1}\rho^{1/2}$,
\beq
\left|\lambda_{p}-\lambda_{q}\right|=\left|\lambda_{p}-\lambda_{-q}\right|\leq \const\eps_H^{-3}|p+q|\leq \const\rho^{1/2-4\eta}
\eeq
For real $a,b$ one checks $|\sqrt a\sqrt b-a|\leq 2|a-b|$, and the same bound for $\overline{\sqrt a\sqrt b}$. Applying these results 
to  $\theta$ with  $k_1+k_2=k_3+k_4\in P_L$, together with \eqref{ublambdaPH0},  we have
\beq
\left|\theta-\lambda_{k_1}\lambda_{k_3}\right|\leq \const\rho^{1/2-6\eta}
\eeq
Inserting this inequality into \eqref{temp8.147} and using  $Q_\Psi(0,k_1+k_2)\leq N\lceil m_c\rceil\leq 2Nm_c$, we obtain 
\beq\label{temp8.150}
\left|\eqref{temp8.140}-4\lambda_{k_1}\lambda_{k_3}|\Lambda|^{-2}Q_\Psi(0,v)\right|\leq \rho^{5/4}m_c|\Lambda|^{-1}, \; v=k_1+k_2
\eeq
Summing over $v\in P_L$ and $k_1$, $k_3\in P_{H,c}$, 
we have that the left hand side of \eqref{resultinMk1k2} is equal to 
\beq
\overline\lim_{k_c\to \infty}\,\overline\lim_{\rho\to0 }4\|w^2V\|_1 \sum_{v\in P_L}Q_\Psi(0,v) \rho^{-5/2} |\Lambda|^{-2}
\eeq
With \eqref{sum0PL}, we have proved \eqref{resultinMk1k2}. 
\end{proof}





\section{Interaction  Energy with Four High Momentum Legs II:  The Error Terms} 

Our goal in this section is to prove that the interaction energy associated  with four high momentum legs 
which are not covered by Lemma \ref{101}  is negligible. We state it as the following lemma.  Notice that 
Lemma \ref{lemHAS2} follows from the results in the previous two sections and this lemma. 

\begin{lem}\label{113}
\beqa\nonumber
&&\overline\lim_{k_c\to\infty}\,\overline\lim_{\rho\to0}\left(\rho^{5/2}|\Lambda|\right)^{-1}\sum_{\eqref{conditionk2}}\sum_{\be\notin M(k_1,k_2)}\\\nonumber
&&\qquad\qquad\left|\frac{V_{k_1-k_3}}{|\Lambda|}f(\be)f(\gamma)\left\langle\be| a^\dagger_{k_1}a^\dagger_{k_2}a_{k_3}a_{k_4}|\gamma\right\rangle\right|\\\label{totalresultinMk1k2}=&&0
\eeqa
\end{lem}

We start with the following lemma.
\begin{lem}\label{thehardestlemma}
For any fixed $k_c\geq 1$ and $\rho$ small enough (depending on $k_c$), we have
\beqa\label{resulthardest}
&&\sum_{\eqref{conditionk2}}\sum_{\be, \gamma: \beta \notin M(k_1,k_2)}\left|f(\be)f(\gamma)\right|  \leq |\Lambda|= o(\rho^{5/2}|\Lambda|^2)
\eeqa
where  the summation is restricted to all $\beta, \gamma\in M$ such that 
\beq\label{defnastk1234}
\left\langle\be| a^\dagger_{k_1}a^\dagger_{k_2}a_{k_3}a_{k_4}|\gamma\right\rangle\neq 0
\eeq
\end{lem}


\begin{proof}  In this section, we use the following notations: 
\beq
\mathcal A_{-k, k}\al\equiv \mathcal A^k \al\rmand \mathcal A_{-k+\frac u2,\,k+\frac u2}\al\equiv\mathcal A^{u,k}\al
\eeq
For any $\{v_1,\cdots, v_t\}\subset P_L$ such that  $v_i\neq \pm v_j, 1\leq i\neq j\leq t$, and any $\al\in M$ with $\al\in M_{v_i}^s, 1\leq i\leq t$,  define 
\beq\label{elementform2}
 M(\al,s,\{v_1,\cdots,v_t\})\equiv \{\prod_{i=1}^{t+s}\mathcal A_{q_{i}, q^{'}_{i}}\al,\,\,\,\,\,q_i,\,\,q_i'\in P_{H,c}, q_{i}+ q^{'}_{i}= u_{i}\}
\eeq
where  $u_i = v_{i}, 1\le i \le t$ and $u_{i} = 0$ otherwise. Since $v_{i}\in P_{L}$ and 
all other momenta are in $P_{H, c}$, $\mathcal A_{q_{i}, q'_{i}}$'s commutes with one another. 

\begin{prop}\label{propchi}
For any  $\chi\in M$, there exists $(\al,s,\{v_1,\cdots,v_t\})$ such that 
\beq\label{resultbegaM}
\chi\in M(\al,s,\{v_1,\cdots,v_t\})
\eeq
%and $(\al,s,\{v_1,\cdots,v_t\})$ is independent of  $\left\{\chi(u), u\in P_{H,c}\right\}$.
\end{prop}

\begin{proof}
By definition of $M$, $\chi$ is obtained from $\al_{free}=|N\rangle$ by finitely many pair creations allowed in Definition \ref{M}. Moving the strict pair creations of momenta outside $P_{H,c}$ to the front, we can write the state $|\chi\rangle$ as follows:
\beq\label{temp8.164}
|\chi\rangle=\prod_{i=1}^{t}\mathcal A_{p_i,p'_i}  \prod_{k=1}^{s}\mathcal A_{q_k,-q_k} \prod_{j=1}^w(
\mathcal A_{u_j, -u_{j}})^{n_j}|N\rangle, 
\eeq
where $u_j\notin P_{H,c}$, $v_i:= p_i+p'_i \in P_L$, $p_i, p'_i, q_k\in P_{H,c}$. Furthermore, we require that 
$u_j\neq \pm u_{j'}$ for $ j \not= j'$. Notice that 
$\mathcal A_{p, p'}$ commute with $\mathcal A_{q, -q}$  so that their orderings are not important. 
The reordering keeps all intermediate states in $M$. Indeed, the number of zero momentum particles at each step is at least the number annihilated by the remaining operations; the operations moved to the front change the nonzero momenta only outside $P_{H,c}$, and there they only create particles, while the other operations annihilate particles of nonzero momenta only at the $v_i$, one particle each, after all strict pair creations of momenta $\pm v_i$; at the $n$-th strict pair creation of momenta $\pm u$, $u\in P_L$, the condition of rule 3 reads $n-1<m_c$ in both orders, since a soft pair creation lowers only one of $\al(u)$, $\al(-u)$, by one; and the condition $\al(v_i)=\al(-v_i)$ of rule 4 still holds, since strict pair creations preserve it. 
Hence 
\beq
\al=\prod_{j=1}^w(\mathcal A_{u_j,-u_j})^{n_j}|N\rangle
\eeq
belongs to $M$, $\al(u)=\al(-u)$ for all $u$, and $\al=0$ on $P_{H,c}$. Since a soft pair creation of total momentum $\pm v$ requires $\al(v)=\al(-v)$, which is destroyed by it and restored by no other operation, we have $v_i\neq \pm v_{i'}$ for $i \not = i'$. Finally, $\al(0)=\chi(0)+2s+t\geq 2s+t$ and $\al(v_i)\geq 1$ (the particle of momentum $v_i$ annihilated by the soft pair creation was created by a strict pair creation), so none of the operators in \eqref{elementform2} annihilates the state. 
This yields that $\chi\in M(\al, s, \{v_1,v_2\cdots v_t\})$ and proves the proposition. 
\end{proof}
%
%One can check that for fixed $\chi$,  the numbers $n_j$, $t$ and $s$, the sets $\{v_i,1\leq i\leq t\}$ and  $\{u_j\}$ in the  \eqref{temp8.164} is independent of  $\left\{\chi(u), u\in P_{H,c}\right\}$. We choose
%So, with \eqref{elementform}, we have 
%\beq\label{chimchi}
%\eeq
%Here $(\al, s, \{v_i,1\leq i\leq t\})$ is independent of  $\left\{\chi(u), u\in P_{H,c}\right\}$.



If $\be, \gamma$ satisfy \eqref{defnastk1234}, then $\be(u)=\gamma(u)$ for all $u\notin\{k_1,k_2,k_3,k_4\}$, in particular for $u\notin P_{H,c}$. 
The triple $(\al, s, \{v_i,1\leq i\leq t\})$ constructed in the proof of  Proposition \ref{propchi} depends only on the values of $\chi$ outside $P_{H,c}$: $\{v_1,\cdots,v_t\}=\{n\in P_L: \chi(n)<\chi(-n)\}$; $\al(v_i)=\chi(v_i)+1$, $\al(n)=\chi(n)$ for the other $n\notin P_{H,c}\cup\{0\}$, and $\al=0$ on $P_{H,c}$; $2(s+t)=\sum_{n\in P_{H,c}}\chi(n)$, which is equal to $N-\sum_{n\notin P_{H,c}}\chi(n)$; and $\al(0)=\chi(0)+2s+t$. Applying this to $\chi=\be$ and $\chi=\gamma$, we obtain the same triple $(\al, s, \{v_i,1\leq i\leq t\})$, so that
\beq
\be\rmand \gamma\in M(\al, s, \{v_1,\cdots,v_t\})
\eeq
Notice that $\al\in M_u^s$ for any $u\in P_L$. 

For any $(\al,s,\{v_1,\cdots,v_t\})$, define $N(\al,s,\{v_1,\cdots,v_t\})$ as the set of the pairs $(\beta,\gamma)$ such that 
\begin{enumerate}
	\item $\be,\,\,\,\gamma\in M(\al,s,\{v_1,\cdots,v_t\})$
	\item there exist $k_i, i=1, \ldots , 4$ satisfying \eqref{conditionk2}, $\be\notin M(k_1,k_2)$ and \eqref{defnastk1234} holds. 
\item for any other  $\al',s',\{v'_1,\cdots,v'_{t'}\}$ s.t.  $\be,\,\,\,\gamma\in M(\al',s',\{v'_1,\cdots,v'_{t'}\})$, then
\beq\label{stst}
s+t\leq s'+t' 
\eeq
\end{enumerate}
We assume $\be, \gamma\in M(\al,s,\{v_1,\cdots,v_t\})$ and \eqref{defnastk1234} holds. Clearly,  $s+t=1$ implies that  $\be\in M(k_1,k_2)$. 
If $t=0$, all factors in \eqref{elementform2} are strict pair creations, which preserve $\chi(q)-\chi(-q)$ for every $q$; hence $\be(k_1)-\be(-k_1)=\al(k_1)-\al(-k_1)=\gamma(k_1)-\gamma(-k_1)$, which is impossible, since by \eqref{defnastk1234} and \eqref{conditionk2} we have $\be(k_1)=\gamma(k_1)+1$ and $\be(-k_1)=\gamma(-k_1)$. 
Hence if $N(\al,s,\{v_1,\cdots,v_t\}) $ is not an empty set then 
\beq\label{stgeq2}
s+t\geq 2 \;  \text { and } \;
t\geq 1
\eeq 

By definition of $N(\al,s,\{v_1,\cdots,v_t\})$, we have
\beqa\label{temp8.173}
&&\sum_{\eqref{conditionk2}}\sum_{\be\notin M(k_1,k_2)}\left|f(\be)f(\gamma)\right|\\\nonumber
\leq &&4\sum_{\al, s, \{v_1\cdots v_t\}}\left|N(\al,s,\{v_1,\cdots,v_t\})\right|\max_{\be,\gamma\in M(\al,s,\{v_1,\cdots,v_t\})}{\left|f(\be)f(\gamma)\right|}, 
\eeqa
where $\left|N(\al,s,\{v_1,\cdots,v_t\})\right|$ is the cardinality of $N(\al,s,\{v_1,\cdots,v_t\})$. Here each pair $(\be,\gamma)$ on the left side belongs to some $N(\al,s,\{v_1,\cdots,v_t\})$, since $\be$ and $\gamma$ lie in a common $M(\al,s,\{v_1,\cdots,v_t\})$ and we may take one with minimal $s+t$; and it appears for at most $4$ choices of $(k_1,k_2,k_3,k_4)$, since by \eqref{defnastk1234} it determines $\{k_1,k_2\}$ and $\{k_3,k_4\}$.
By definition of $f$ (see \eqref{properf1} and \eqref{properf4}),  if $\be,\gamma\in M(\al,s,\{v_1,\cdots,v_t\})$ then 
\beq
|f(\beta)f(\gamma)|\leq \left|\frac{\al(0)}{|\Lambda|}\right|^{2s+t} \left|\frac {4\lceil m_c\rceil}{|\Lambda|}\right|^{t }\max_{k\in P_H}\{|\lambda_k|\}^{2s+2t}|f(\al)|^2
\eeq
{F}rom \eqref{ublambdaPH0}, $\al(0)\leq N$, $\lceil m_c\rceil\leq 2m_c$ and $m_c=\rho^{-\eta}$, we have
\beq
\max_{\be,\gamma\in M(\al,s,\{v_1,\cdots,v_t\})}{\left|f(\be)f(\gamma)\right|}\leq (\const\rho^{1-5\eta})^{2s+t}|\Lambda|^{-t}|f(\al)|^2\\\nonumber
\eeq
Together with \eqref{temp8.173}, the right hand side of \eqref{temp8.173} is bounded by (the factor $4$ is absorbed in the constant, since $2s+t\geq 2$ by \eqref{stgeq2})
\beq\label{temp8.176}
\le  \sum_{\al, s, \{v_1\cdots v_t\}}\left|N(\al,s,\{v_1,\cdots,v_t\})\right|(\const\rho^{1-5\eta})^{2s+t}|\Lambda|^{-t}|f(\al)|^2
\eeq
Define $N(\al,s,t)$ and $ N(s,t)$ by 
\beq
 N(\al, s,t)\equiv \max_{ \{v_1,\cdots,v_t\}}\left\{\left| N(\al, s,\{v_1,\cdots,v_t\})\right|\right\}
\eeq
\beq
 N(s,t)\equiv \max_{\al}\left\{ N(\al, s,t)\right\}
\eeq
With \eqref{temp8.176}, we can bound \eqref{temp8.173} by 
\beqa\nonumber
\eqref{temp8.173}&&\leq \sum_{\al, s,t}|f(\al)|^2\sum_{\{v_1\cdots v_t\}} N(\al, s,t)(\const\rho^{1-5\eta})^{2s+t}|\Lambda|^{-t}\\
&&\leq \sum_{s,t}\sum_{\{v_1\cdots v_t\}} N( s,t)(\const\rho^{1-5\eta})^{2s+t}|\Lambda|^{-t}
\eeqa
For fixed $t$ the total number of set $\{v_1\cdots v_t, v_i \in P_L\}$ is bounded by
$$
\sum_{\{v_1\cdots v_t\}}1\leq (\Lambda\rho^{3/2}\eta_L^{-3})^t(t!)^{-1}\leq (\rho^{1-5\eta})^{\frac{3t}2}|\Lambda|^{t}(t!)^{-1}
$$
{F}rom  $t\leq (\Lambda\rho^{3/2}\eta_L^{-3})\leq \rho^{-1.65}$ and  \eqref{stgeq2}, we have 
\beq\label{temp8.178}
\sum_{\eqref{conditionk2}}\sum_{\be\notin M(k_1,k_2)}\left|f(\be)f(\gamma)\right | \leq \sum_{t=1}^{\rho^{-1.65}}\sum_{s: s+t\geq 2} N( s,t)(\const\rho^{1-5\eta})^{2s+\frac{5t}2}(t!)^{-1}
\eeq


\begin{lem}\label{lem11.1}
For any fixed $k_c\geq 1$ there is $\bar\rho(k_c)>0$ such that, for $\rho\leq\bar\rho(k_c)$ and any $N(\al, s,\{v_1,\cdots,v_t\})$ with $s+t\geq 2$ and $t\geq 1$, we have 
\beq\label{0.025}
\left|N(\al, s,\{v_1,\cdots,v_t\})\right|\leq t\,!\,t^{(\frac {t}2+1)}|\Lambda|^{\frac{s+t}2+1} (\rho^{-5\eta})^{t+s}
\eeq
\end{lem}

{F}rom this Lemma and  $|\Lambda|=\rho^{-\frac{25}{8}}$, for $\rho\leq\bar\rho(k_c)$ the right hand side of \eqref{temp8.178} is bounded above by 
\beqa\nonumber
&&\sum_{t\geq 1}^{\rho^{-1.65}}\sum_{s: s+ t>1}\,t\left(\rho^{5/2}|\Lambda|^{1/2}t^{1/2}\right)^t\left(\rho^{2}|\Lambda|^{1/2}\right)^s \left(\const\rho^{-35\eta/2}\right)^{t+s}|\Lambda|\\\nonumber
=&&\sum_{t\geq 1}^{\rho^{-1.65}}t(\const\rho^{0.85}t^{1/2})^{t}\sum_{s: s+ t>1}(\const\rho^{0.35})^s|\Lambda|\leq |\Lambda|
\eeqa
This proves Lemma \ref{thehardestlemma}.  

 \end{proof}




We now prove Lemma \ref{lem11.1}.
\begin{proof}
Since  $(\beta,\gamma)\in N(\al, s, \{v_1,\cdots, v_t\})$, we can express them as  
\beq\label{112}
\be =\prod_{j=t+1}^{s+t}\mathcal A_{ q_{2j-1},\, q_{2j}}\prod_{i=1}^{t}A_{ q_{2i-1},\, q_{2i}}\al,  \; \gamma=\prod_{j=t+1}^{s+t}\mathcal A_{\widetilde q_{2j-1},\,\widetilde q_{2j}}\prod_{i=1}^{t}A_{\widetilde q_{2i-1},\,\widetilde q_{2i}}\al
\eeq
and $ q_{2i-1}+ q_{2i} = v_i = \widetilde q_{2i-1} + \widetilde q_{2i}$ for $i=1, \ldots, t$, $q_{2j-1}+ q_{2j}= \widetilde q_{2j-1} + \widetilde q_{2j}=0$ for $t+1\leq j\leq s+t$. 
%Here $\{u\}$ is the union of $t$ zeros and the set $\{v_1,\cdots,v_t\}$, so is $\{\widetilde u\}$.
{F}rom \eqref{defnastk1234}, we have 
 \beq
 \{q_1,\cdots,\,q_{2s+2t}\}-\{k_1,k_2\}=\{\widetilde q_1,\cdots,\,\widetilde q_{2s+2t}\}-\{k_3,k_4\}
 \eeq

Denote the common elements in $\{q_i\}$ and $\{\widetilde q_i\}$ by  $p_1$, $p_2,$ $\cdots$, $p_{2s+2t-2}$. Then 
we have  \beq
 \{q_i\}=k_1,\,\,\,k_2,\,\,\,p_1,\,\,\,p_2,\,\,\,\cdots, \,\,\,p_{2s+2t-2},\,\,\,
 \eeq
 \beq
 \{\widetilde q_i\}=k_3,\,\,\,k_4,\,\,\,p_1,\,\,\,p_2,\,\,\,\cdots, \,\,\,p_{2s+2t-2},\,\,\,
 \eeq
We now construct a graph with vertices  $ \{ k_1, k_2, k_3, k_4, p_i, 1\le i \le 2s+2t -2 \} $.  The edges of the graphs consisting of 
$\beta$ edges $(q_{2i-1},q_{2i}), 1\le i \le s+t$ and $\gamma$ edges $(\widetilde q_{2j-1},\widetilde q_{2j}), 1\le i \le s+t$.  
{F}rom  \eqref{defnastk1234},  the graph can be decomposed into two chains and loops. Thus there exist
 $l$, $m_i\in \Z$ and $0<m_1<m_2<......<m_l=s+t$ such that 
%
%Connect $q_1$ and $q_2$ with a line, then $q_{3}$ and $q_{4}$, $\cdots$, $q_{2s+2t-1}$ and $q_{2s+2t}$. Do the same thing to $\widetilde q$'s. In this graph, we can see the degrees of $k_i$'s are equal to one and the  degrees of $p_i$'s are equal to two, which means this graph can be described as $2$ chains and several loops,
% i.e. 
\beqa\label{temp4.178}
&&k_1\longleftrightarrow p_1\longleftrightarrow p_2\longleftrightarrow p_3\cdots p_{2m_1-1}\longleftrightarrow k_2(\rmor\cdots\,k_4)\\\nonumber
&&k_3\longleftrightarrow p_{2m_1}\longleftrightarrow p_{2m_1+1}\cdots p_{2m_2-2}\longleftrightarrow k_4(\rmor\cdots k_2)
\\\nonumber
&&p_{2m_2-1}\longleftrightarrow p_{2m_2}{\longleftrightarrow}p_{2m_2+1}\cdots p_{2(m_3)-2}\longleftrightarrow p_{2m_2-1}
\\\nonumber
&&\cdots\\\nonumber
&&\cdots
\\\nonumber
&&p_{2m_{l-1}-1}\longleftrightarrow p_{2m_{l-1}}\longleftrightarrow p_{2m_{l-1}+1}\cdots p_{2(m_l)-2}\longleftrightarrow   p_{2m_{l-1}-1}
\eeqa
Here we have relabeled the indices of $p$ and do not distinguish $\beta$ edges and $\alpha$ edges. We also disregard the obvious 
symmetry $k_1 \to k_2$ and $k_3 \to k_4$.  
Due to the condition  \eqref{stst}, the length of the loop must be $4$ or more, i.e., for $3\leq i\leq l$
\beq
m_{i-1}+2\leq m_i
\eeq
Indeed, a loop of length $2$ consists of a $\be$ edge and a $\gamma$ edge with the same end points $q,q'$, i.e., $\mathcal A_{q,q'}$ is a factor of both products in \eqref{112}. Since $\al\in M$, $\al\in M^s_{v_i}$, $\al(v_i)\geq 1$ and $\al(0)\geq 2s+t$ (the last two follow from \eqref{112}), rules 2 and 4 of Definition \ref{M} give $\al':=\mathcal A_{q,q'}\al\in M$, and $\al'\in M^s_{v_j}$ for all $v_j\neq q+q'$. Hence $\be,\gamma\in M(\al',s',\{v'_1,\cdots,v'_{t'}\})$ with $s'+t'=s+t-1$, where $\{v'_j\}$ is $\{v_j\}$ with $q+q'$ removed if $q+q'\neq 0$, and $s'=s-1$ if $q+q'=0$. This contradicts \eqref{stst}. The same argument shows that every partial product of the factors in \eqref{112} applied to $\al$ belongs to $M$. 

Moreover, $(k_1,k_2)$ is not a pair of $\be$ and $(k_3,k_4)$ is not a pair of $\gamma$ in \eqref{112}. Indeed, if $(k_1,k_2)$ is a pair of $\be$, then $k_1+k_2=v_i$ for some $i$, and removing this factor from \eqref{112} gives, by the last remark, $\be'\in M$ with $\be'(\pm v_i)=\al(\pm v_i)$, so that $\be'\in M^s_{v_i}$ and $\be=\mathcal A^{v_i,\,k_1/2-k_2/2}\be'$, i.e., $\be\in M(k_1,k_2)$. If $(k_3,k_4)$ is a pair of $\gamma$, then in the same way $\gamma=\mathcal A^{v_i,\,k_3/2-k_4/2}\gamma'$ with $\gamma'\in M^s_{v_i}$, $v_i=k_3+k_4=k_1+k_2$, and by \eqref{11.1} $\be=\mathcal A^{v_i,\,k_1/2-k_2/2}\gamma'\in M(k_1,k_2)$. Both contradict $\be\notin M(k_1,k_2)$. Hence each of the two chains has at least two edges, i.e., $0<m_1<m_2$. 
Together with $m_l=s+t$, we obtain
\beq\label{uboundl}
l\leq (s+t)/2+1, \quad t \ge 1.
\eeq
Without loss of generality, we assume  for $3\leq i<j\leq l$
\beq
m_i-m_{i-1}\leq m_j-m_{j-1}
\eeq
Denote by $N(\al, s, \{v_1,\cdots,v_t\}, l, \{m_1,\cdots,m_l\})$ the set of all pairs  $(\beta,\,\,\,\gamma)$ having the graph above
and we now estimate its cardinality. 


We can add the information between $k_i$'s and $p_i$'s as follows
\beqa
&&k_1\stackrel{w_1}{\longleftrightarrow} p_1\stackrel{\widetilde w_1}{\longleftrightarrow} p_2\stackrel{w_2}{\longleftrightarrow} p_3\cdots p_{2m_1-1}\stackrel {w_{m_1}}{\longleftrightarrow} k_4(\rmor \cdots k_2)\\\nonumber
&&k_3\stackrel{\widetilde w_{m_1}}{\longleftrightarrow} p_{2m_1}\stackrel{w_{m_1+1}}{\longleftrightarrow} p_{2m_1+1}\cdots p_{2m_2-2}\stackrel{{\widetilde w_{m_2}}}{\longleftrightarrow} k_2(\rmor \cdots k_4)
\\\nonumber
&&p_{2m_2-1}\stackrel{w_{m_2+1}}{\longleftrightarrow} p_{2m_2}{\stackrel{\widetilde w_{m_2+1}}{\longleftrightarrow}}p_{2m_2+1}\cdots p_{2(m_3)-2}\stackrel{\widetilde w_{m_3}}{\longleftrightarrow} p_{2m_2-1}
\\\nonumber
&&\cdots\\\nonumber
&&\cdots
\\\nonumber
&&p_{2m_{l-1}-1}\stackrel{w_{m_{l-1}+1}}{\longleftrightarrow} p_{2m_{l-1}}\stackrel{\widetilde w_{m_{l-1}+1}}{\longleftrightarrow} p_{2m_{l-1}+1}\cdots p_{2(m_l)-2}\stackrel{\widetilde w_{m_{l}}}{\longleftrightarrow}   p_{2m_{l-1}-1}\, , 
\eeqa
where  $A\stackrel{c}{\longleftrightarrow} B$ if and only if  $A+B=c$. And $w_i$'s the union of $s$ zero's and $\{v_1,\cdots, v_t\}$, so are $\widetilde w$'s. 
%We can see 
%\beqa\label{temp5.154}
%&&\be =\mathcal A^{0,w_{s+t}}_{p_{2(s+t)-3},\,\,p_{2(s+t)-2}}\cdots\mathcal A^{0,w_2}_{p_2,\,\,p_3}\mathcal A^{0,w_1}_{k_1,\,\,p_1}\al\\\nonumber
%&&\gamma =\mathcal A^{0,w_{s+t}}_{p_{2m_{l-1}-1},\,\,p_{2(s+t)-2}}\cdots\mathcal A^{0,\widetilde w_1}_{p_1,\,\,p_2}\al
%\eeqa
By \eqref{112}, $\beta$ and $\gamma$ is uniquely determined by  $w_i$'s, $\widetilde w_i$'s and one $k_i$ or $p_i$ for  each loop or chain. 
\par 

To bound $| N(\al, s, \{v_1,\cdots,v_t\}, l, \{m_1,\cdots,m_l\}) |$, we note that the sum of momentum in each loop is zero.  
%$j\geq 3$, we have
%\beq\label{conditionml}
%\sum_{i=m_{j-1}+1}^{m_j}\left(w_i-\widetilde w_i\right)=0
%\eeq
Thus we can count the number of graphs  as follows. 
\begin{enumerate}
	\item  choose the positions of zeros in $\beta$ edges. The total number of choices is less than $2^{t+s}$. 
	\item   choose the positions of $v_1\cdots v_t$ in $\beta$ edges. The total number of choices is  $t!$.
	\item   choose the positions of zeros in $\gamma$ edges. The total number of choices is less than $2^{t+s}$. 
	\item   choose the positions of $v_1\cdots v_t$ in $\gamma$ edges.  We call a loop trivial if all 
	the momenta associated with $\gamma$  edges are zero.   The number of trivial loops is at most $s/2$ since 
	there are at least two $\gamma$ edges per loop. Since $l-2$ of the $l$ components of the graph are loops, the number of non-trivial loops is at least $l-2-s/2$. 
	Thus we only have to fix  $v$ in at most $t-(l-2-s/2)$ edges and  the number of choices is at most $t^{t-l+2+s/2}$.
	%\item Finally, fill up the blank we left in loops, with (\ref{conditionml}), the total number of choices is no more than 1. 
\end{enumerate}
Thus we obtain 
\beqa\label{alstvlm}
 &&|N(\al, s, \{v_1,\cdots,v_t\}, l, \{m_1,\cdots,m_l\})|\\\nonumber
  \leq&& (\const)^{t+s}t! t^{(t+s/2-l+2)}\left(k_c^3|\Lambda|\right)^{l}\\\nonumber
 \leq &&(\const)^{t+s}t! t^{(t/2+1)}\left(k_c^3|\Lambda|\right)^{t/2+s/2+1}
\eeqa
where we have used (\ref{uboundl}) and $t\leq k_c^3|\Lambda|$.
Since  $$| N(\al, s, \{v_1,\cdots,v_t\})|=\sum_{l}\sum_{\{m_1,\cdots,m_l\}} | N(\al, s, \{v_1,\cdots,v_t\}, l, \{m_1,\cdots,m_l\}) | 
$$ and 
\beq
\sum_{l}\sum_{\{m_1,\cdots,m_l\}} 1\leq \const^{s+t}
\eeq 
we obtain $| N(\al, s, \{v_1,\cdots,v_t\})|\leq t!\,t^{(t/2+1)}|\Lambda|^{\frac{s+t}2+1}(\const k_c^3)^{s+t}$, where we have used $(k_c^3)^{\frac{s+t}2+1}\leq (k_c^3)^{s+t}$ for $s+t\geq 2$. Since $\const k_c^3\leq \rho^{-5\eta}$ for $\rho\leq\bar\rho(k_c)$, we have proved \eqref{0.025}.
\end{proof}




We now prove Lemma \ref{113}. 
\begin{proof} Let $\be$, $\gamma\in M$ s.t.$\left\langle\be| a^\dagger_{k_1}a^\dagger_{k_2}a_{k_3}a_{k_4}|\gamma\right\rangle\neq 0$. Using Lemma \ref{id4a} and the definition of $f$, we have 
\beqa\nonumber
&&\left|f(\be)f(\gamma)\left\langle\be|  a^\dagger_{k_1}a^\dagger_{k_2}a_{k_3}a_{k_4}|\gamma\right\rangle\right|=\left|f(\be)f(\gamma)\right|\sqrt{\be(k_1)\be(k_2)\gamma(k_3)\gamma(k_4)}\\\label{temp8.162}
\leq &&|f(\be)|^2\sqrt{\left|\frac{\lambda_{k_3}\lambda_{k_4}}{\lambda_{k_1}\lambda_{k_2}}\right|}(\be(k_1)+\be(k_2))\sqrt{\gamma(k_3)\gamma(k_4)}
\eeqa
{F}rom  the bound on  $\lambda_{k_i}$'s in \eqref{ublambdaPH0}, $\gamma(k_i)\leq N$ and  $N=\rho^{-17/8}$, we have
\beqa\nonumber
&&\left|f(\be)f(\gamma)\left\langle\be|  a^\dagger_{k_1}a^\dagger_{k_2}a_{k_3}a_{k_4}|\gamma\right\rangle\right|\leq |f(\be)|^2\left|\lambda_{k_1}\lambda_{k_2}\right|^{-\frac{1}{2}}(\be(k_1)+\be(k_2))\rho^{-\frac{9}{4}}
\eeqa
For fixed $k_1,k_2,k_3$, the sum of the right side over $\be$ with $\be(k_1)>3$ is bounded by $\rho^{-9/4}|\lambda_{k_1}\lambda_{k_2}|^{-1/2}\sum_{m\geq 4,\,n\geq 1}(m+n)Q_\Psi(\{k_1,m\},\{k_2,n\})$, since $\be(k_2)=\gamma(k_2)+1\geq 1$. By \eqref{boundpsiuvmn} (recall $k_1\neq\pm k_2$), whose factor $|\lambda_{k_1}|^m|\lambda_{k_2}|^n$ absorbs $|\lambda_{k_1}\lambda_{k_2}|^{-1/2}$, and by \eqref{ublambdaPH0}, this is at most $\const\rho^{31/4}\eps_H^{-18}$. The sum over $\be$ with $\be(k_2)>3$ is bounded in the same way. Since there are at most $(\const k_c^3|\Lambda|)^3$ choices of $(k_1,k_2,k_3)$, we have 
\beq
\sum_{\eqref{conditionk2}}\sum_{\be(k_1)>3 \rmor \be(k_2)>3}\left|f(\be)f(\gamma)\left\langle\be| a^\dagger_{k_1}a^\dagger_{k_2}a_{k_3}a_{k_4}|\gamma\right\rangle\right|\leq o(\rho^{5/2}|\Lambda|^2)
\eeq
By symmetry, we have
\beq
\sum_{\eqref{conditionk2}}\sum_{\gamma(k_3)>3 \rmor \gamma(k_4)>3}\left|f(\be)f(\gamma)\left\langle\be| a^\dagger_{k_1}a^\dagger_{k_2}a_{k_3}a_{k_4}|\gamma\right\rangle\right|\leq o(\rho^{5/2}|\Lambda|^2)
\eeq
To prove \eqref{totalresultinMk1k2}, we only have  to focus on the case  $\be(k_{i})\leq 3, i=1, 2$ and $\gamma(k_{i})\leq 3, i=3, 4$. In this case, by \eqref{temp8.162},  we have
\beq\left|f(\be)f(\gamma)\left\langle\be|  a^\dagger_{k_1}a^\dagger_{k_2}a_{k_3}a_{k_4}|\gamma\right\rangle\right|\leq \left|\const f(\be)f(\gamma)\right|
\eeq
Using Lemma \ref{thehardestlemma}, we arrive at the desired result \eqref{totalresultinMk1k2}.
\end{proof}





\section{Proof of Lemma \ref{lemma1}}
\label{sec:lemma22}



The proof of Lemma \ref{lemma1} follows a standard approach and we only give the main steps.
We first  construct an isometry between functions with periodic
boundary condition in $[0, L]^3$ and functions with Dirichlet boundary condition in $[-\ell,\,\,\,L+\ell]^3$, where $0<2\ell<L$.
Denote the coordinates of $\x$ by   $\x = (x^{(1)}, x^{(2)}, x^{(3)})$.
Fix a smooth function $\zeta:\R\to[0,1]$ with $\zeta(\tau)=0$ for $\tau\leq 1/4$ and $\zeta(\tau)=1$ for $\tau\geq 3/4$, and let
$\vartheta(x)=\frac{\pi}{2}\, \zeta\big((\ell-x)/2\ell\big)$; thus $\vartheta(x)=\pi/2$ for $x\leq -\ell/2$ and $\vartheta(x)=0$ for $x\geq \ell/2$.
Let  $h(\x)$ be the function $h(\x) = q(x^{(1)}) q(x^{(2)}) q(x^{(3)})$ where
\beq
q(x)=\left\{
\begin{array}{ll}
\cos\vartheta(x), & |x|\leq \ell\\
1, & \ell<x<L-\ell\\
\sin\vartheta(x-L), & |x-L|\leq \ell\\
0, & \text{otherwise}
\end{array}\right.
\eeq
The four cases are disjoint since $2\ell<L$. The function  $q$ is smooth, $q=1$ on $[\ell/2,L-\ell/2]$ and $q=0$ outside $(-\ell/2,L+\ell/2)$, so
$h$ is smooth and supported in $[-\ell/2,L+\ell/2]^3$, a compact subset of the open cube $(-\ell,L+\ell)^3$.
(The choice $\zeta(\tau)=\tau$, i.e., $\vartheta(x)=(\ell-x)\pi/4\ell$, would also do, but the corresponding $q$ is only Lipschitz continuous.)
Since $q(x)^2+q(x+L)^2=\cos^2\vartheta(x)+\sin^2\vartheta(x)=1$ for $|x|\leq\ell$, we have $\sum_{m\in\Z}q(x+mL)^2=1$ for all $x\in\R$, hence
\[
\sum_{n\in\Z^3}h(\x+nL)^2=1,\qquad \x\in\R^3 .
\]
Writing an integral over $\R^3$ as the integral over $[0,L]^3$ of the sum over the translates by $L\Z^3$, we obtain, for any  function $\phi$ with
the period   $L$,
\beq\label{proh1}
\int_{\x\in[-\ell,\,L+\ell\,]^3}|h \phi(\x)|^2=\int_{\x\in[0,L]^3}|\phi(x)|^2
\eeq
Thus the map $ \phi\longrightarrow h \phi$
is an isometry: $$
L^2_{
{\rm Periodic}}\left([0,L]^3\right)\to L^2_{
{\rm Dirichlet}}\left([-\ell,L+\ell]^3\right).$$



Let $\chi (\x)$ be the characteristic function of the $\ell$-boundary of $[0,L]^3$, i.e.,
$\chi (\x)=1$ if   $|x^{(\alpha)} |\le \ell$ for some $\alpha = 1, 2$ or $3$ where $|x^{(\alpha)} |$
is  the distance on the torus.
We have $|\nabla (h\phi)|^2=h^2|\nabla\phi|^2+|\phi|^2|\nabla h|^2+\frac12\nabla(h^2)\cdot\nabla|\phi|^2$.
Folding the integrals onto $[0,L]^3$ as above, the first term gives $\int_{[0,L]^3}|\nabla\phi|^2$, the last one gives zero since
$\sum_n\nabla(h^2)(\x+nL)=\nabla 1=0$, and
\[
\sum_{n\in\Z^3}|\nabla h(\x+nL)|^2=\sum_{\alpha=1}^3\sum_{m\in\Z}q'(x^{(\alpha)}+mL)^2\leq \frac{3\pi^2\sup|\zeta'|^2}{16\,\ell^2}\,\chi(\x) ,
\]
because $q'(x)^2+q'(x+L)^2=\vartheta'(x)^2\leq (\pi/4\ell)^2\sup|\zeta'|^2$ for $|x|\leq \ell$, and $q'=0$ on $[\ell,L-\ell]$.
This yields the following estimate on the kinetic energy of  $h \phi$
\beqa\label{3.51}
&&\int_{\x\in[-\ell,\,L+\ell\,]^3}|\nabla (h  \phi)(\x)|^2\\\nonumber
\leq &&\int_{\x\in[0,\,L]^3}|\nabla \phi(\x)|^2+\const\ell^{-2}\int_{\x\in[0,\,L]^3}  \chi (\x) | \phi(\x)|^2
\eeqa
%which means
%\beq\label{temp2.15}
%\langle -\Delta\rangle_{h\widetilde\phi}\leq \langle -\Delta\rangle_{\phi}+\const\ell^{-2}\langle |x|
%\leq \ell\rangle_{\widetilde\phi}
%\eeq

The generalization of this isometry to higher dimensions is straightforward.  Suppose
$\Psi(\x_1,\cdots,\x_N)$ is a symmetric function with period $L$ in each $\x_i$.
Then  for any $u \in \R^{3}$, the map
 \beq\label{defFu}
{\cal F}^u(\Psi): = \Psi (\x_1,\cdots,\x_N) \prod_{i=1}^N h(\x_i+u)
\eeq
is an  isometry from $L^2_{
{\rm Periodic}}\left([0,L]^{3N}\right)$ to $ L^2_{
{\rm Dirichlet}}\left([-\ell-u,L+\ell-u]^{3N}\right)$; the function ${\cal F}^u(\Psi)$ is symmetric and its support is a compact subset of the open cube $(-\ell-u,L+\ell-u)^{3N}$.
Clearly,   ${\cal F}^u$ has the property \eqref{3.51}:
\[
\int|\nabla {\cal F}^u(\Psi)|^2\leq \int_{[0,L]^{3N}}|\nabla\Psi|^2+\const\ell^{-2}\sum_{i=1}^N\int_{[0,L]^{3N}}\chi(\x_i+u)\,|\Psi|^2 .
\]



We recall from Section 2 the position form of \eqref{ham}. Let $V^P$ be the periodization of $V$,
\[
V^P(\x):=\sum_{n\in\Z^3}V(\x+nL) . \tag{12.4$'$}\label{defVP}
\]
Since $V$ is smooth with fast decay, the series converges uniformly on compact sets, and $V^P$ is a bounded continuous function with period $L$. Moreover,
$\int_{[0,L]^3}V^P(\x)e^{-iu\x}\rd\x=V_u$ for $u\in\Lambda^*$, i.e., $V^P(\x)=|\Lambda|^{-1}\sum_{u}e^{iu\x}V_u$ by the Poisson summation formula.
Hence, in terms of the plane waves $|\Lambda|^{-1/2}e^{ip\x}$, $p\in\Lambda^*$, the Hamiltonian \eqref{ham} on $\mathcal H_N$ is the operator
\[
H_N=\sum_{i=1}^N-\Delta_{\x_i}+\sum_{i\neq j}V^P(\x_i-\x_j)
\]
on symmetric functions on $[0,L]^{3N}$ with periodic boundary condition (each pair of particles is counted twice, as in \eqref{ham}).
With the Dirichlet boundary condition on a cube, the Hamiltonian is, as in Section 2, $H_N=\sum_{i}-\Delta_{\x_i}+\sum_{i\neq j}V(\x_i-\x_j)$ on symmetric functions vanishing on the boundary.
%
% Because $V^P(x-y)$ statisfies periodic boudary condition on $[0,L]^6$, we can extend it to $\widetilde V^P(x-y)$ s.t.
%\beq
%\widetilde V^P(x-y)=\widetilde V^P(x+\Z^3L-y)
%\eeq
Since $V$ is nonnegative, each term of the series \eqref{defVP} is bounded by its sum: $V(\x+nL)\leq V^P(\x)$ for all $n\in\Z^3$.
{F}rom the definition of    ${\cal F}^u$, we conclude that, for $i\neq j$,
\beq\nonumber
\int |{\cal F}^u(\Psi)|^2V(\x_i-\x_j)\prod_{k=1}^N d\x_k \le
\int_{[0,\,L]^{3N}} |\Psi|^2 V^P(\x_i-\x_j)\prod_{k=1}^Nd\x_k
\eeq
Indeed, folding the integral over $(\R^3)^N$ onto $[0,L]^{3N}$, the left side equals
$\int_{[0,L]^{3N}}|\Psi|^2\sum_{n_1,\cdots,n_N\in\Z^3}\prod_{k=1}^Nh(\x_k+u+n_kL)^2\,V(\x_i-\x_j+(n_i-n_j)L)$, and $\sum_{n_1,\cdots,n_N}\prod_kh(\x_k+u+n_kL)^2=1$.

Therefore,  for every normalized symmetric periodic $\Psi$ with finite kinetic energy, the energies of the two boundary conditions are related by
\beq
\left\langle H_N\right\rangle_{{\cal F}^u(\Psi)}\leq \left\langle H_N\right\rangle_{\Psi}+\const\ell^{-2}\sum_{i=1}^N\left\langle \chi (\x_i+u)\right\rangle_{\Psi}
\eeq
where on the left $H_N$ is the Hamiltonian with the Dirichlet boundary condition on $[-\ell-u,L+\ell-u]^3$ and on the right the periodic one on $[0,L]^3$.
Averaging over $u\in[0,L]^3$ and using $\int_{[0,L]^3}\chi(\x+u)\,du=L^3-(L-2\ell)^3\leq 6\ell L^2$, we have
\beq
\int_{[0,L]^3}\left\langle H_N\right\rangle_{{\cal F}^u(\Psi)}du
\leq L^3\left\langle H_N\right\rangle_{\Psi}+\const\ell^{-1}L^2N
\eeq
So for any $\Psi$ there exists an $u$ such that
\beq\label{Fuexists}
\left\langle H_N\right\rangle_{{\cal F}^u(\Psi)}\leq \left\langle H_N\right\rangle_{\Psi}+\const N(\frac{1}{\ell L})
\eeq
If we choose  $\ell$ and $L$ as
\beq\label{temp2.2}
\ell=\rho^{-25/48},\,\,\,  L=\rho^{-25/24},
\eeq
 the error term is negligible to the accuracy we need in proving  Lemma \ref{lemma1}, as we now show.
Here and below $\rho$ runs through the values for which $N=\rho L^3=\rho^{-17/8}$ is an integer, i.e., $\rho=N^{-8/17}$, and $\rho\to0$ means $N\to\infty$.
With the choice \eqref{temp2.2} we have $\ell/L=\rho^{25/48}$, so that $2\ell<L$ for $\rho<2^{-48/25}$, the error term in \eqref{Fuexists} is $\const N\rho^{75/48}$,
and the cube $[-\ell-u,L+\ell-u]^3$ is a translate of $\widetilde\Lambda=[0,\widetilde L]^3$, $\widetilde L=L+2\ell=L(1+2\rho^{25/48})$,
in which the $N$ particles have the density $\widetilde\rho=N/\widetilde L^3=\rho L^3/\widetilde L^3$ of Lemma \ref{lemma1}.

Let $\delta>0$. By the definition of $e_0^P(\rho,\Lambda)$ through the bottom of the spectrum, there is a normalized symmetric periodic $\Psi$ with finite kinetic energy such that
$\langle H_N\rangle_\Psi\leq N(e_0^P(\rho,\Lambda)+\delta)$. For the $u$ of \eqref{Fuexists}, the function ${\cal F}^u(\Psi)$, translated by $(\ell,\ell,\ell)+u$ in each variable,
is a normalized symmetric trial state with the Dirichlet boundary condition on $\widetilde\Lambda$, and the Hamiltonian is translation invariant. Hence
$e_0^D(\widetilde\rho,\widetilde\Lambda)\leq e_0^P(\rho,\Lambda)+\delta+\const\rho^{75/48}$; since $\delta$ is arbitrary,
\[
e_0^D(\widetilde\rho,\widetilde\Lambda)\leq e_0^P(\rho,\Lambda)+\const\rho^{75/48}.
\]
It remains to change the density in the normalization from $\rho$ to $\widetilde\rho$. We have $\rho/\widetilde\rho=(1+2\rho^{25/48})^3\to1$ and
$0\leq\rho-\widetilde\rho\leq 6\rho^{73/48}$, hence
\[
\frac{e^D_0(\widetilde \rho,\widetilde\Lambda)-g_0\widetilde \rho }{ g_0^{5/2} \widetilde \rho ^{\,3/2}}\leq
\left(\frac{\rho}{\widetilde\rho}\right)^{3/2}\frac{e^P_0(\rho,\Lambda)-g_0\rho }{ g_0^{5/2} \rho ^{3/2}}
+\frac{6g_0\rho^{73/48}+\const\rho^{75/48}}{g_0^{5/2}\,\widetilde\rho^{\,3/2}} .
\]
The last term is $O(\rho^{1/48})$ and $(\rho/\widetilde\rho)^{3/2}\to 1$, so the hypothesis of Lemma \ref{lemma1} implies \eqref{resultlemma1}.
This concludes the proof of Lemma \ref{lemma1}.
We remark that the exponents in \eqref{temp2.2} are chosen so that both the boundary energy $\const(\ell L)^{-1}$ and the change of density,
$g_0(\rho-\widetilde\rho)\approx 6g_0\rho\,\ell/L$, are $o(\rho^{3/2})$; any $\ell=\rho^{-\mu}$ with $22/48<\mu<26/48$ would do.



\subsection{Proof of Lemma \ref{thm1}}\label{subsec:lemma21}

By the variational principle, $E_0^P(\rho,\Lambda)\leq\langle H_N\rangle_\Psi$ for every normalized $\Psi\in\mathcal H_N$ with finite energy, in particular for the trial state of Theorem \ref{lemma2}.
Hence Theorem \ref{lemma2} implies the hypothesis of Lemma \ref{lemma1}, and by Lemma \ref{lemma1} the bound \eqref{resultlemma1} holds.

We write the quantities of Lemma \ref{lemma1} in terms of $N$: $\rho_N=N^{-8/17}$ and $L_N=\rho_N^{-25/24}=N^{25/51}$, so that $N=\rho_NL_N^3$, $\widetilde L_N=L_N(1+2\rho_N^{25/48})=N^{25/51}+2N^{25/102}$, $\widetilde\Lambda_N=[0,\widetilde L_N]^3$ and
\[
\widetilde\rho_N=\frac{N}{\widetilde L_N^3}=N^{-8/17}\left(1+2N^{-25/102}\right)^{-3}.
\]
For each $N$ choose a normalized symmetric $\Psi_N$ satisfying the Dirichlet boundary condition on $\widetilde\Lambda_N$ (it may be taken smooth with compact support in the open cube) such that
$\langle H_N\rangle_{\Psi_N}N^{-1}\leq e_0^D(\widetilde\rho_N,\widetilde\Lambda_N)+\widetilde\rho_N^{\,2}$. Since $\widetilde\rho_N^{\,2}=o(\widetilde\rho_N^{\,3/2})$, \eqref{resultlemma1} gives
\beq\label{lemma21N}
\overline\lim_{N\to\infty}\left(\frac{\langle H_N\rangle_{\Psi_N}N^{-1}-g_0\widetilde\rho_N}{g_0^{5/2}\,\widetilde\rho_N^{\,3/2}}\right)\leq\frac{16}{15\pi^2}.
\eeq
This is Lemma \ref{thm1} along the sequence of densities $\widetilde\rho_N$, with the cubes $\widetilde\Lambda_N$; note that $\widetilde L_N\widetilde\rho_N^{\,25/24}=(1+2N^{-25/102})^{-51/24}\to1$.
In this form, with a trial state for every large $N$, the lemma is used in Section \ref{subsec:dup}.

To obtain Lemma \ref{thm1} for all small $\rho$, we read $L\sim\rho^{-25/24}$ as $L(\rho)\rho^{25/24}\to1$ as $\rho\to0$.
For $0<\rho\leq\widetilde\rho_1=1/27$ let $N=N(\rho)$ be the largest integer with $\widetilde\rho_N\geq\rho$ (it exists since $\widetilde\rho_N\to0$);
thus $\widetilde\rho_{N+1}<\rho\leq\widetilde\rho_N$, and $N(\rho)\to\infty$ as $\rho\to0$ since $N(\rho)\geq n$ for $\rho\leq\widetilde\rho_n$. Let $L(\rho)=(N/\rho)^{1/3}$, so that $N=\rho L(\rho)^3$ and $L(\rho)\geq\widetilde L_N$.
Extended by zero, $\Psi_N$ is a trial state of $N$ particles on $[0,L(\rho)]^3$ satisfying the Dirichlet boundary condition, with the same energy.
Since $N^{-25/102}$ is decreasing in $N$,
\[
\frac{\widetilde\rho_N}{\widetilde\rho_{N+1}}\leq\left(\frac{N+1}{N}\right)^{8/17}\leq1+\frac{8}{17N},
\qquad\text{hence}\qquad 1\leq\frac{\widetilde\rho_N}{\rho}\leq1+\frac{8}{17N} .
\]
Moreover $\rho\leq\widetilde\rho_N\leq N^{-8/17}$ and $\rho>\widetilde\rho_{N+1}\geq(N+1)^{-8/17}/27$, so $N(\rho)$ is of order $\rho^{-17/8}$. Consequently
$L(\rho)\rho^{25/24}=(N\rho^{17/8})^{1/3}=(\rho/\widetilde\rho_N)^{17/24}(1+2N^{-25/102})^{-17/8}\to1$, and, with $e=\langle H_N\rangle_{\Psi_N}N^{-1}$,
\[
\frac{e-g_0\rho}{g_0^{5/2}\rho^{3/2}}=\left(\frac{\widetilde\rho_N}{\rho}\right)^{3/2}\frac{e-g_0\widetilde\rho_N}{g_0^{5/2}\,\widetilde\rho_N^{\,3/2}}+\frac{\widetilde\rho_N-\rho}{g_0^{3/2}\rho^{3/2}},
\]
where $0\leq(\widetilde\rho_N-\rho)\rho^{-3/2}\leq 8/(17N\rho^{1/2})=O(\rho^{13/8})$.
Since $N(\rho)\to\infty$, \eqref{lemma21N} implies Lemma \ref{thm1} with $L=L(\rho)$, $N=N(\rho)$ and the trial state $\Psi_{N(\rho)}$.



\subsection{Proof of Theorem \ref{mainthm}: duplication}\label{subsec:dup}

We now carry out the duplication described after Lemma \ref{thm1}: copies of the Dirichlet states $\Psi_n$ of \eqref{lemma21N}, with $n=N_0$ or $n=N_0+1$ particles,
are placed in a periodic array of cubes separated by corridors of width at least $\rho^{-1/2}$, and their symmetrized product is used as a trial state on the large torus $\Lambda=[0,L]^3$.
Throughout, $V^P$ is the periodization \eqref{defVP} with the period $L$ of the large torus, $c_*=16/(15\pi^2)$, $e_0^P(\rho,\Lambda)$ is the ground state energy per particle of $N$ particles on $\Lambda$ as in Theorem \ref{mainthm},
and $g_0>0$ (the ratios in Theorem \ref{mainthm} require $V\not\equiv0$).

\smallskip\noindent{\it Passage to the torus.}
Let $\Phi$ be a normalized symmetric function on $(\R^3)^N$ with finite kinetic energy whose support is a compact subset of $(0,L)^{3N}$.
Its $L$-periodic extension in each variable has finite kinetic energy on the torus, and by the position form of \eqref{ham} in Section \ref{sec:lemma22} and the variational principle
\beq\label{ePDir}
e_0^P(\rho,\Lambda)\leq\frac1N\int_{[0,L]^{3N}}\Big(\sum_{i=1}^N|\nabla_{\x_i}\Phi|^2+\sum_{i\neq j}V^P(\x_i-\x_j)|\Phi|^2\Big).
\eeq

\smallskip\noindent{\it Decay of $V$.}
Since $V$ has fast decay, $C_8:=\sup_{z\in\R^3}|z|^8V(z)<\infty$. Any bound $V(z)\leq C_k|z|^{-k}$ with $k>29/4$ would suffice below.

\smallskip\noindent{\it The small cubes.}
Fix $\e>0$. By \eqref{lemma21N} there is $N_\e$ such that for every $n\geq N_\e$
\beq\label{dupbox}
\langle H_n\rangle_{\Psi_n}\leq n\left(g_0\widetilde\rho_n+(c_*+\e)\,g_0^{5/2}\widetilde\rho_n^{\,3/2}\right).
\eeq

\smallskip\noindent{\it Scales.}
Let $\rho>0$ be small, $d=\rho^{-1/2}$, and let $N_0$ be the largest integer $N\geq1$ with
\beq\label{dupN0}
N+\frac12\geq\rho\left(\widetilde L_{N+1}+2d\right)^3 .
\eeq
Since $\widetilde L_{N+1}\geq(N+1)^{25/51}$, \eqref{dupN0} implies $N+1<\rho^{-17/8}$; on the other hand, for $N=\lfloor\rho^{-17/8}/2\rfloor$ the ratio of the left side of \eqref{dupN0} to the right side tends to $2^{8/17}>1$ as $\rho\to0$.
Hence $N_0$ is well defined for small $\rho$, and
\[
\rho^{-17/8}/4\leq N_0\leq\rho^{-17/8}-1 ;
\]
in particular $N_0\geq N_\e$ for $\rho$ small enough. Let $b=\widetilde L_{N_0+1}$ and $s=((N_0+1/2)/\rho)^{1/3}$. Then $\rho s^3=N_0+1/2$, and \eqref{dupN0} for $N=N_0$ says $s\geq b+2d$.
Since \eqref{dupN0} fails for $N_0+1$, we have $N_0+3/2<\rho(\widetilde L_{N_0+2}+2d)^3$.
The function $N\mapsto\widetilde L_N=N^{25/51}+2N^{25/102}$ is increasing with derivative at most $50/51$ on $[1,\infty)$, so $\widetilde L_{N_0+2}\leq\widetilde L_{N_0}+2$, and $\widetilde L_{N_0}\geq N_0^{25/51}\geq\rho^{-25/24}/2$.
Therefore, for $n\in\{N_0,N_0+1\}$,
\beq\label{dupdensity}
\widetilde\rho_n=\frac{n}{\widetilde L_n^3}\leq\frac{N_0+3/2}{\widetilde L_{N_0}^3}<\rho\left(1+\frac{2+2d}{\widetilde L_{N_0}}\right)^3\leq\rho\,(1+8\rho^{13/24})^3=:\rho\,(1+\kappa),
\eeq
with $\kappa=O(\rho^{13/24})$.

\smallskip\noindent{\it Geometry.}
Let $L>s$, $K=\lfloor L/s\rfloor\geq1$ and $s_L=L/K$; then $s\leq s_L\leq sL/(L-s)$. For $m\in\Z^3$ let
\[
B_m=s_Lm+d\,(1,1,1)+(0,b)^3 .
\]
Since $s_L\geq s\geq b+2d$, the distance between $B_m$ and $B_{m'}$ is at least $|m-m'|_\infty s_L-b\geq 2d\,|m-m'|_\infty$ for $m\neq m'$.
The $K^3$ cubes $B_m$, $m\in I_K=\{0,\cdots,K-1\}^3$, lie in $(0,L)^3$ at distance at least $d$ from its boundary, and, since $L=Ks_L$,
the cubes $B_m$, $m\in\Z^3$, are exactly the translates of these $K^3$ cubes by $L\Z^3$.

\smallskip\noindent{\it Particle numbers.}
If $L\to\infty$ and $N/L^3\to\rho$, then $s_L\to s$, so that $N/K^3=(N/L^3)\,s_L^3$ tends to $\rho s^3=N_0+1/2$. Hence, for all $(N,L)$ sufficiently close to the thermodynamic limit,
$K^3N_0\leq N<K^3(N_0+1)$, i.e., $N=K^3N_0+j$ with $0\leq j<K^3$. We put $n_m=N_0+1$ particles in $j$ of the cubes $B_m$, $m\in I_K$, and $n_m=N_0$ particles in each of the other cubes; thus $\sum_mn_m=N$.

\smallskip\noindent{\it Trial state.}
Since $\widetilde L_{N_0}\leq b$, the states $\Psi_{N_0}$ and $\Psi_{N_0+1}$, extended by zero, satisfy the Dirichlet boundary condition on $(0,b)^3$.
Let $\psi_m(\x_1,\cdots,\x_{n_m})=\Psi_{n_m}(\x_1-y_m,\cdots,\x_{n_m}-y_m)$, $y_m=s_Lm+d\,(1,1,1)$, be the copy of $\Psi_{n_m}$ in $B_m$;
by translation invariance its energy is $\langle H_{n_m}\rangle_{\Psi_{n_m}}$.
Split $\{1,\cdots,N\}$ into consecutive blocks $J_m$ of sizes $n_m$, let $T(\x_1,\cdots,\x_N)=\prod_{m\in I_K}\psi_m(\x_{J_m})$, where $\x_{J_m}$ denotes the variables with indices in $J_m$, and
\[
\Phi=\Big(N!\prod_mn_m!\Big)^{-1/2}\sum_{\sigma\in S_N}T\circ\sigma ,
\]
where $(T\circ\sigma)(\x_1,\cdots,\x_N)=T(\x_{\sigma(1)},\cdots,\x_{\sigma(N)})$.
If $\sigma$ preserves every block, then $T\circ\sigma=T$ since each $\psi_m$ is symmetric; otherwise $T\circ\sigma$ and $T$ have disjoint supports, since the closures of the cubes $B_m$ are disjoint.
Hence $\Phi$ is symmetric and normalized, it has finite kinetic energy, and its support is a compact subset of $(0,L)^{3N}$.
Since the kinetic and the potential energy densities are local, all cross terms between different functions $T\circ\sigma$ vanish, and the energy of $\Phi$ on the right side of \eqref{ePDir} (times $N$) equals that of $T$, namely
\[
\sum_{m\in I_K}\int\sum_{i\in J_m}|\nabla_{\x_i}\psi_m|^2+\int\sum_{i\neq j}V^P(\x_i-\x_j)\,|T|^2 .
\]

\smallskip\noindent{\it Interaction between cubes.}
Let $\x\in{\rm supp}\,T$ and $i\in J_m$, so that $\x_i\in\overline{B_m}$. For $j\in J_{m'}$ and $n\in\Z^3$, the point $\x_j-nL$ lies in $\overline{B_{m'-Kn}}$. Hence in
\[
\sum_{j\neq i}V^P(\x_i-\x_j)=\sum_{j\neq i}\sum_{n\in\Z^3}V\big(\x_i-(\x_j-nL)\big)
\]
the points $\x_j-nL$ in $\overline{B_m}$ are exactly those with $n=0$ and $j\in J_m$, while every other cube $\overline{B_{m''}}$ contains at most $N_0+1$ of the points $\x_j-nL$
and is at distance at least $2d\,|m''-m|_\infty$ from $\x_i$. Since $\#\{m'':|m''-m|_\infty=r\}=24r^2+2$ and $V(z)\leq C_8|z|^{-8}$,
\[
\sum_{j\neq i}V^P(\x_i-\x_j)\leq\sum_{j\in J_m,\,j\neq i}V(\x_i-\x_j)+(N_0+1)\,C_8\,(2d)^{-8}\sum_{r\geq1}(24r^2+2)\,r^{-8},
\]
and the last sum is less than $27$. We integrate against $|T|^2$ and sum over $i$. With \eqref{ePDir}, \eqref{dupbox}, \eqref{dupdensity}, $N_0+1\leq\rho^{-17/8}$ and $(2d)^{-8}=\rho^4/256$ we obtain
\beqa\nonumber
e_0^P(\rho,\Lambda)&\leq&\frac1N\sum_{m\in I_K}\langle H_{n_m}\rangle_{\Psi_{n_m}}+27\,C_8\,(N_0+1)(2d)^{-8}\nonumber\\
&\leq&g_0\rho\,(1+\kappa)+(c_*+\e)\,g_0^{5/2}\rho^{3/2}(1+\kappa)^{3/2}\nonumber\\\label{dupfinal}
&&+\,\frac{27}{256}\,C_8\,\rho^{15/8}.
\eeqa

\smallskip\noindent{\it Conclusion.}
The bound \eqref{dupfinal} holds for all $(N,L)$ sufficiently close to the thermodynamic limit at the density $\rho$; $\Phi$ is a trial state as required in \eqref{resultpro1}. Hence
\[
\overline\lim_{L\to\infty}\left(\frac{e_0^P(\rho,\Lambda)-g_0\rho}{g_0^{5/2}\rho^{3/2}}\right)\leq g_0^{-3/2}\kappa\,\rho^{-1/2}+(c_*+\e)(1+\kappa)^{3/2}+\frac{27\,C_8}{256\,g_0^{5/2}}\,\rho^{3/8} .
\]
Since $\kappa\rho^{-1/2}=O(\rho^{1/24})$, the right side tends to $c_*+\e$ as $\rho\to0$. As $\e>0$ is arbitrary, this proves Theorem \ref{mainthm}.

The error terms come from the corridors and the two particle numbers ($O(\rho^{1/24})$ after the normalization by $g_0^{5/2}\rho^{3/2}$) and from the interaction between different cubes, including the periodic images ($O(\rho^{3/8})$).
With $V(z)\leq C_k|z|^{-k}$ in place of $C_8|z|^{-8}$ the last term is $O(\rho^{k/2-29/8})$, which tends to zero for $k>29/4$.
Finally, since the support of $\Phi$ is a compact subset of $(0,L)^{3N}$, $\Phi$ is also a trial state for the Dirichlet Hamiltonian $\sum_i-\Delta_{\x_i}+\sum_{i\neq j}V(\x_i-\x_j)$ on $\Lambda$, with the same kinetic energy. Since $V\geq0$, we have $V(\x_i-\x_j)\leq V^P(\x_i-\x_j)$, so its energy per particle in $\Phi$ is at most the right side of \eqref{ePDir}.
Hence, by the variational principle, \eqref{dupfinal} holds for $e_0^D(\rho,\Lambda)$ as well. This proves the remark after Theorem \ref{mainthm} on the Dirichlet boundary condition.




\appendix

\section{Proofs of Propositions \ref{prop1} and \ref{prop2}}\label{app:props}

Every $\al\in M$ is obtained from $\al_{free}$ by a finite sequence of steps, each of which is one of the operations of Definition \ref{M}:
a strict pair creation with momentum in $P_I\cup P_H$ (rule 2) or in $P_L$ (rule 3), or a soft pair creation (rule 4), each allowed at the state to which it is applied.
As maps on $\widetilde M$, these operations add fixed vectors to $\al$, so $\al$ depends only on the collection of steps and not on their order.
The order matters only through the conditions of the rules: the occupation numbers must stay nonnegative, rule 3 requires $\max\{\al(u),\al(-u)\}<m_c$ and rule 4 requires $\al(u)=\al(-u)$.
These conditions involve only the occupation numbers at $0$ and in $P_L$.
Particles with momenta in $P_I\cup P_H$ are never annihilated; particles with momenta in $P_L$ are created only by strict pair creations and annihilated only by soft pair creations.

\begin{proof}  (Proof of Proposition \ref{prop1}) Recall $M$ is constructed through the four operations in the Definition \ref{M}. Notice that
particles with momentum in $P_I$ can only be created via a pair creation and no furthermore annihilation
of particles with momentum in $P_I$ is allowed. In other words, once a pair of particles with momenta
$u, -u \in P_I$ are created, they stay there forever and do not participate in further operations.  So the creation
of a pair of particles with momenta
$u, -u \in P_I$ commute with all other operations in creating $\alpha$. We can now define $\beta$ to be the state
with all operations of creating the state  $\alpha$ except deleting one creation of particles with momenta
$u, -u \in P_I$.
Deleting this step changes the states to which the later steps are applied only by $+2$ at the momentum $0$ and by $-1$ at $\pm u\in P_I$.
Since the conditions of the rules involve only the occupation numbers at $0$ and in $P_L$, every remaining step is still allowed, so $\be\in M$.
Since $\be(0)=\al(0)+2\geq2$, $\mathcal A^u\be$ is defined.
Then clearly,  $\mathcal A^{u}\beta=\al$ and $\be(u)=m-1$.

\end{proof}





\begin{proof} (Proof of Proposition \ref{prop2}) Recall $M$ is constructed through the four operations in the Definition \ref{M}. Notice that
a particle with the momentum  $u \in P_H$ can  be created from two operation: 1. The creation of a pair
$u, -u$. 2. A creation of a soft pair, i.e., there is a $v \in P_L$ such that particles $u$ and $v-u$
are created while a particle of zero momentum and a particle of momentum $v$ are annihilated.
Clearly, if the first alternative occurs, then we can construct $\beta$ as in the proof of  Proposition \ref{prop1}, since the conditions of the rules do not involve momenta in $P_H$.
So we can now focus on the second alternative. In this case, $u$, $v-u\in P_H$
and the soft pair creation can be realized by the operator $a^+_{u} a^+_{v-u} a_0 a_v$. Denote by $T_j$ one of the four operations
allowed in the definition of $M$.   Hence we can write the state
$\alpha$ as
\[
|\alpha\rangle =     C     \prod_{j=1}^n  T_j |N\rangle, \quad
T_s= a^+_{u} a^+_{v-u} a_0 a_v
\]
for some $ 1 \le s \le n$, where the factors are ordered so that $T_1$ acts first, and $C$ is the normalization constant.
Let $\al_j$ denote the state after the first $j$ steps, so $\al_0=\al_{free}$ and $\al_n=\al$.
Since $u$, $v-u \in P_H$ and particles in $P_H$ will not be annihilated,
we have $a^+_{u} a^+_{v-u}$ commute with all $T_j$. Clearly, $a_0$ commute with all $T_j$. Hence we have
\[
|\alpha \rangle =   C a^+_{u} a^+_{v-u} a_0        \prod_{j=s+1}^n  T_j  a_v  \prod_{i=1}^{s-1} T_i |N\rangle,
\]
Since in our definition, we allow $v$ to participate in the soft pair creation only when particles with momentum $v$ and
$-v$ perfectly paired, we know that $T_j, j \ge s+1$ will not involve annihilation of particles with momentum $v$ or $-v$.
Indeed, $\al_{s-1}(v)=\al_{s-1}(-v)$ and the step $s$ lowers $\al(v)$ by one; strict pair creations of momentum $\pm v$ do not change $\al(-v)-\al(v)$,
and a soft pair creation from $\pm v$ would require $\al(v)=\al(-v)$. Hence
\[
\al_j(-v)=\al_j(v)+1, \qquad s\leq j\leq n,
\]
and no step after the step $s$ annihilates particles with momentum $\pm v$; in particular $\al\in M^a_v$.
Thus the operators of momentum $v$ in $T_j$, $j\geq s+1$, are creation operators only, say $(a_v^+)^k$ in total.
For any state $|\gamma\rangle$ with $\gamma(v)\geq1$, we have
\[
(a_v^+)^k a_v |\gamma\rangle =  D  a_v (a_v^+)^{k}  |\gamma\rangle
\]
for some constant $D>0$. Applying this to $\prod_{i=1}^{s-1} T_i |N\rangle$, which is a multiple of $|\al_{s-1}\rangle$ and $\al_{s-1}(v)\geq1$, we have
\[
|\alpha \rangle =   \widetilde C a^+_{u} a^+_{v-u} a_0 a_v       \prod_{j=s+1}^n  T_j    \prod_{i=1}^{s-1} T_i |N\rangle,
\]
for some normalization constant $\widetilde C$. We can now define
\[
|\beta \rangle =    C'     \prod_{j=s+1}^n  T_j    \prod_{i=1}^{s-1} T_i |N\rangle,
\]
with $C'$ some normalization constant.
It remains to check that $\be\in M$, i.e., that the steps $T_1,\cdots,T_{s-1},T_{s+1},\cdots,T_n$ are allowed in this order.
For $j<s$ nothing changes. For $j>s$, the state to which $T_j$ is applied is now $\al_{j-1}$ with one more particle at $0$ and at $v$, and one particle less at $u$ and at $v-u$.
The occupation numbers at $0$ and in $P_L$ did not decrease, so no annihilation fails.
At $\pm v$ we now have $\al_{j-1}(v)+1=\al_{j-1}(-v)$ particles, so $\max\{\al(v),\al(-v)\}$ is unchanged and the condition of rule 3 for a strict pair creation of momentum $\pm v$ holds as before;
no step after the step $s$ is a soft pair creation from $\pm v$, and at the other momenta of $P_L$ nothing changed, so the conditions of rules 3 and 4 hold for all steps. Hence $\be\in M$.
Moreover $\be(v)=\al(v)+1=\al(-v)=\be(-v)$, so $\be\in M^s_v$; since $\be(0)\geq1$, $\be(v)\geq1$ and $\e_H<|u|,|v-u|\leq k_c$, the soft pair creation $T_s$ is allowed at $\be$ by rule 4,
and $\al=\mathcal A^{v,\,u-v/2}\be$. Finally $|v|<2\e_H<|2u|$, so $v-u\neq u$ and $\be(u)=\al(u)-1=m-1$.
This proves the proposition.

\end{proof}



%:003














\begin{thebibliography}{99}
%\bibitem{CW} M.H. Anderson, J.R. Ensher,
%M.R. Matthews, C.E. Wieman, and E.A. Cornell, {\it Science} ({\bf 269}),
%198 (1995).



\bibitem{B} N.N. Bogoliubov, {\it On the theory of superfluidity}, (in Russian). Izv. Academii Nauk USSR \textbf{11} (1): 77.  {\it On the theory of superfluidity}, (in English) Journal of Physics \textbf{11} (1): 23--32.
%{\bf 11}, 77 (1947). English translation: J. Phys. (USSR), {\bf 11} 23 (1947). See
%also {\sl Lectures on quantum statistics,} vol. 1. Gordon and Breach (1967)

%\bibitem{Kett} K. B. Davis, M. -O. Mewes, M. R. Andrews, N. J. van Druten, D. S. Durfee, D. M. Kurn and W. Ketterle, {\it Phys. Rev. Lett.} ({\bf 75}), 3969 (1995).

\bibitem{D} F.J. Dyson, {\it Ground-State Energy of a Hard-Sphere Gas},
Phys. Rev. \textbf{106}, 20--26 (1957).



\bibitem{ESY} L. Erd{\H o}s, B. Schlein and H.-T. Yau, {\it Ground-state energy of a low-density Bose gas: A second-order upper bound}, Phys. Rev. A \textbf{78}, 053627 (2008)

\bibitem{GA} M. Girardeau and R. Arnowitt, {\it Theory of Many-Boson Systems: Pair Theory }, Phys. Rev., {\bf 113}, No. 3, 755-761 (1959).

\bibitem{GS} A. Giuliani, R. Seiringer, {\it The ground state energy of the weakly interacting Bose gas at high density},  arXiv:0811.1166v1 (2008).

\bibitem{HY} K. Huang and C.N. Yang, {\it Quantum-Mechanical Many-Body Problem with Hard-Sphere Interaction}, Phys. Rev. {\bf 105}, No. 3., 767--775 (1957).


\bibitem{Landau} L. Landau and E. Lifshitz,   {\it Quantum mechanics,  non-relativistic theory}, 3ed., Pergamon, 1991.


\bibitem{LHY} T.D. Lee, K. Huang and C.N. Yang, {\it Eigenvalues and Eigenfunctions of a Bose System of Hard Spheres and Its Low-Temperature Properties}, Phys. Rev. {\bf 106}, No. 6., 1135--1145 (1957)


\bibitem{LYang} T.D. Lee and C.N. Yang, {\it Many body problem in quantum mechanics and quantum statistical mechanics}, Phys. Rev. {\bf 105}, No. 3, 1119--1120 (1957)

\bibitem{L} E.H. Lieb, {\it Simplified Approach to the Ground State
Energy of an Imperfect Bose Gas}, Phys. Rev. {\bf 130}, 2518--2528
(1963). See also  Phys. Rev. {\bf 133}, A899-A906 (1964) (with A.Y.
Sakakura) and Phys. Rev. {\bf 134}, A312-A315 (1964) (with W. Liniger).




%\bibitem{LSY} E.H. Lieb, R. Seiringer and J. Yngvason, Phys. Rev A \textbf{61} (2000), 043602.

\bibitem{LS1}E.H.Lieb, J.P. Solovej, {\it Ground State Energy of the One-Component Charged Bose Gas}, Commun. Math. Phys. {\bf 217}, 127--163 (2001).

\bibitem{LS2}E.H.Lieb, J.P. Solovej, {\it Ground State Energy of the Two-Component Charged Bose Gas}, Commun. Math. Phys. {\bf 252}, 448--534 (2004).


\bibitem{LY}  E.H. Lieb and J. Yngvason, {\it Ground State Energy of the Low Density Bose Gas}, Phys. Rev. Lett. {\bf 80}, 2504-2507 (1998)

%\bibitem{Pit} L. Pitaevskii and S. Stringari, {\sl Bose-Einstein Condensation}, Oxford Science Publications, Oxford (2003)



\bibitem{Sol} J.P. Solovej, {\it Upper Bounds to the Ground State Energies of the One- and Two-Component Charged Bose Gases}, Comm. Math. Phys. {\bf 266}, No. 3, 797-818 (2006)


\bibitem{Y} C.N. Yang, {\it Dilute Hard "Sphere" Bose Gas in Dimensions 2, 4 and 5}, arXiv:0807.0938
\end{thebibliography}




\end{document}

% LocalWords:  Subadditivity subadditivity
